% DER_XIX: Grothendieck's text with the commentary set in, then the overview, inside a group that applies the chunk's own definitions
\dgrchapter{XIX}{Chapter XIX}{Modèles (4)}
\dgrpart{0}{XIX}{Manuscript pages 1--50}{XIX\enspace\textperiodcentered\enspace Chapter XIX\enspace\textperiodcentered\enspace manuscript pages 1--50}{Chapter XIX \textperiodcentered\ \textup{XIX}}
\dgrchunkbegin
% DER_XIX: the definitions of the chunk's preamble, applied inside its group
\let\printedpage\relax\newcommand{\printedpage}[1]{\marginpar{\footnotesize p.~#1}}
\let\msp\relax\newcommand{\msp}[1]{\marginpar{\footnotesize [#1]}}

% Q-orthogonality arrows (half arrowheads = Q-orthogonal, full arrowheads = equality)
\let\Qo\relax\newcommand{\Qo}{\mathrel{\leftharpoondown\joinrel\relbar\joinrel\rightharpoonup}}   % lower half-head left, upper half-head right
\let\Qorev\relax\newcommand{\Qorev}{\mathrel{\leftharpoonup\joinrel\relbar\joinrel\rightharpoondown}} % mirrored form of (1.2)
\let\Qog\relax\newcommand{\Qog}{\mathrel{\leftarrow\joinrel\relbar\joinrel\rightharpoonup}}         % full head left
\let\Qod\relax\newcommand{\Qod}{\mathrel{\leftharpoondown\joinrel\relbar\joinrel\rightarrow}}       % full head right
\let\Qogd\relax\newcommand{\Qogd}{\mathrel{\leftarrow\joinrel\relbar\joinrel\rightarrow}}            % full heads both sides
% the same arrows crossed by a vertical bar (Quillen factorisation)
\let\Qx\relax\newcommand{\Qx}[2]{\mathrel{#1\joinrel\mapstochar#2}}
% crossed by two vertical bars (functorial factorisation)
\let\Qxx\relax\newcommand{\Qxx}[2]{\mathrel{#1\joinrel\mapstochar\mkern2mu\mapstochar#2}}
% Q-closure marks
\let\tld\relax\newcommand{\tld}{\tilde{}}   % left Q-closure mark (small tilde in exponent)
\let\rcl\relax\newcommand{\rcl}{\bar{}}       % right Q-closure mark (horizontal bar in exponent)
\let\Qorth\relax\newcommand{\Qorth}{\mathrel{\leftharpoondown\mkern-6mu\relbar\mkern-6mu\rightharpoonup}}
\let\Qorthx\relax\newcommand{\Qorthx}{\mathrel{\leftharpoondown\mkern-9mu{+}\mkern-9mu\rightharpoonup}}
\let\Qorthxfull\relax\newcommand{\Qorthxfull}{\mathrel{\leftharpoondown\mkern-9mu{+}\mkern-9mu\rightarrow}}
\let\Qcouple\relax\newcommand{\Qcouple}{\mathrel{\leftarrow\mkern-9mu{+}\mkern-9mu\rightarrow}}
\let\clos\relax\newcommand{\clos}{\longleftrightarrow}
\let\fl\relax\newcommand{\fl}{\longrightarrow}
\let\M\relax\newcommand{\M}{\mathcal{M}}
\let\Fl\relax\newcommand{\Fl}{\mathrm{Fl}}
\let\Ho\relax\newcommand{\Ho}{\mathrm{Ho}}
\let\cc\relax\newcommand{\cc}{\mathrm{c}}
\let\ddarr\relax\newcommand{\ddarr}{\mathrel{\twoheadleftarrow\joinrel\twoheadrightarrow}}
\let\harp\relax\newcommand{\harp}{\mathrel{\leftharpoondown\joinrel\relbar\joinrel\rightharpoonup}}
\let\deq\relax\newcommand{\deq}{\overset{\text{déf}}{=}}
\let\Mono\relax\newcommand{\Mono}{\mathrm{Mono}}
\let\Filt\relax\newcommand{\Filt}{\mathrm{Filt}}
\let\Cof\relax\newcommand{\Cof}{\mathrm{Cof}}
\let\qpair\relax\newcommand{\qpair}{\mathrel{\leftharpoondown\joinrel\relbar\joinrel\rightharpoonup}}
\let\Cat\relax\newcommand{\Cat}{\mathrm{Cat}}
\let\Diag\relax\newcommand{\Diag}{\mathrm{Diag}}
\let\id\relax\newcommand{\id}{\mathrm{id}}
\let\Qs\relax\newcommand{\Qs}{\tilde{\ }}
\let\fact\relax\newcommand{\fact}{\mathrel{\leftarrow\mkern-11mu\twoheadrightarrow}}
\let\limind\relax\newcommand{\limind}{\varinjlim}
\let\limproj\relax\newcommand{\limproj}{\varprojlim}
\let\Is\relax\newcommand{\Is}{\mathrm{Is}}

% ---- from der-ch19_p001-005.tex
\printedpage{(non numérotée)}
% head: (none; title page)

\begin{center}
{\Large\textsf{LES DÉRIVATEURS}}

\bigskip

\textbf{ALEXANDRE GROTHENDIECK}

\bigskip\bigskip

\textbf{Chapitre XIX}

\medskip

\textbf{Modèles (4)}
\end{center}

\newpage

\printedpage{(non numérotée)}
% head: (none; editorial note)

Ce texte a été déchiffré et transcrit en \LaTeXe\ par M. Künzer. Il a été édité par M. Künzer, J. Malgoire, et G. Maltsiniotis. La transcription est aussi fidèle que possible au manuscrit. Seules quelques corrections mineures évidentes ont été effectuées. Pour les rares ajouts ou commentaires des éditeurs, ainsi que pour la numérotation originale des pages du manuscrit, les caractères de machine à écrire \texttt{[typewriter]} entre crochets sont utilisés.

Cette édition est provisoire. Les remarques, commentaires et corrections sont bienvenus. Envoyer un message à :

\textbf{[e-mail address removed]}

\hfill G. Maltsiniotis\hspace*{2cm}

\newpage

\printedpage{1}
% head: A. Grothendieck : Dérivateurs, Ch. XIX, Édition des Universités de Montpellier II et Paris VII
% foot: Version 11 février 2019

\msp{page 1}

\section*{1\quad Rappels sur les couples de Quillen (couples de Quillen spéciaux)}

Je vais reprendre les notations et la terminologie pour les relations de LLP \texttt{[left lifting property]} et RLP \texttt{[right lifting property]} de Quillen. J'écris (pour une catégorie $\mathcal{M}$, et $\Phi, \Psi \subset \mathrm{Fl}(\mathcal{M})$)
\[
\Phi \Qo \Psi \qquad \iff \qquad (\Phi \subset \Psi^{*}) \qquad \iff \qquad (\Psi \subset \Phi_{*})\ ,
\tag*{(1.1)}
\]
i.e. $\Phi$ a la LLP par rapport à $\Psi$, ou encore $\Psi$ la RLP par rapport à $\Phi$. Je dis aussi que $(\Phi, \Psi)$ est un \textit{couple} Q-\textit{orthogonal.} Cette notion n'est pas symétrique, mais je me permets parfois (exceptionnellement) d'écrire
\[
\Psi \Qorev \Phi \qquad \text{comme synonyme de} \qquad \Phi \Qo \Psi\ .
\tag*{(1.2)}
\]
$\Phi$ est Q-\textit{orthogonal à gauche} de $\Psi$, $\Psi$ est Q-\textit{orthogonal à droite} à $\Phi$, et la Q-orthogonalité à gauche (resp. à droite) est notée par le demi-sommet de flèche pointant vers le bas (resp. vers le haut). J'écris
\[
\left\{
\begin{aligned}
&\Phi \quad \Qog \quad \Psi \quad \iff \quad (\Phi = \Psi^{*})\\
&\Phi \quad \Qod \quad \Psi \quad \iff \quad (\Psi = \Phi_{*})\\
&\Phi \quad \Qogd \quad \Psi \quad \iff \quad (\Phi = \Psi^{*} \ \text{ et } \ \Psi = \Phi_{*})\ .
\end{aligned}
\right.
\tag*{(1.3)}
\]
Ces trois cas de Q-orthogonalité renforcée sont donc ceux où on a de plus $\Phi = \Phi^{\tld}$
resp. $\Psi = \Psi^{\rcl}$ resp. $\Phi = \Phi^{\tld}$ et $\Psi = \Psi^{\rcl}$. On dit que $\Phi$ est Q-\textit{clos à gauche} si $\Phi = \Phi^{\tld}$, que % reading resolved: DER_XIX_corrections.csv
$\Psi$ est Q-\textit{clos à droite} si $\Psi = \Psi^{\rcl}$, que le couple Q-orthogonal $(\Phi, \Psi)$ est Q-\textit{clos à gauche} (resp. \textit{à droite}) si $\Phi$ (resp. $\Psi$) est Q-clos à gauche (resp. à droite). On dit qu'il est Q-\textit{clos} si les deux conditions % reading resolved: DER_XIX_corrections.csv

\msp{page 2}

sont satisfaites simultanément. Quand aucune confusion n'est à craindre, il m'arrivera de dire simplement que $\Phi$ est Q-clos, ou que $\Psi$ est Q-clos, sans spécifier si c'est à gauche où % [sic] où -> ou, conjunction "or" intended
à droite. Mais pour un couple Q-orthogonal $(\Phi, \Psi)$ (ou je préfère l'écrire $\Phi \Qo \Psi$), Q-clos % [sic] ou -> où, relative adverb intended
signifie toujours Q-clos à gauche et à droite.

\dgrblock{dgrb1158}{Orthogonality, closure, and factorization}
\begin{dcomment}{exposition}{p2:signifie toujours Q-clos à gauche et à droite.}{Orthogonality, closure, and factorization}
Grothendieck begins by separating the ingredients that will later have to survive passage to a category of diagrams. A lifting property is easy to transport along an adjunction. A factorization requires a construction. The purpose of the elaborate arrows is to prevent the existence of that construction from being smuggled into the lifting notation.

For classes of arrows $L,R$ in a locally small category, write $L\perp R$ when every commutative square with left side in $L$ and right side in $R$ has a diagonal filler. The filler need not be unique. In modern notation,
\[
 {}^{\square}R=\{i:i\perp R\},\qquad
 L^{\square}=\{p:L\perp p\}.
\]
These are the source's $R^*$ and $L_*$. They form an antitone Galois connection: $L\subseteq{}^{\square}R$ if and only if $R\subseteq L^{\square}$. Consequently
\[
 \overline L={}^{\square}(L^{\square}),\qquad
 \overline R=({}^{\square}R)^{\square}
\]
are extensive, monotone, idempotent operations. Grothendieck writes a small tilde for the first closure and an overbar for the second. An arrow with two complete heads means the stronger pair of identities $L={}^{\square}R$ and $R=L^{\square}$. A transverse stroke adds the assertion that every $f$ factors as $pi$ with $i\in L$ and $p\in R$. Together these are precisely the data of a weak factorization system, without a chosen functorial factorization.

There is a genuine logical slip in the sentence following (1.3). For an arbitrary orthogonal pair, separate closedness of $L$ and $R$ does not imply that they are one another's orthogonals. Indeed, in $\mathrm{Set}$ take both classes to be the isomorphisms. Each is closed under its appropriate double-orthogonal operation, because
\[
 {}^{\square}\mathrm{Fl}=\mathrm{Is},\quad
 \mathrm{Is}^{\square}=\mathrm{Fl},\quad
 {}^{\square}\mathrm{Is}=\mathrm{Fl},\quad
 \mathrm{Fl}^{\square}=\mathrm{Is}.
\]
For example, an arrow having the left lifting property against itself has a two-sided inverse, by using the square with identity horizontal sides; this proves the first identity. The other three follow in the same way or directly from invertibility. Thus $(\mathrm{Is},\mathrm{Is})$ is separately closed and orthogonal, but is not a mutually maximal orthogonal pair unless the category is a groupoid.

The repair is to retain the explicit equalities of (1.3) as the meaning of a closed \emph{pair}. There is no difficulty with the later terminology for a Quillen pair: once factorization is available, Proposition 1.3 shows that separate retract-closedness forces these equalities. The distinction matters precisely before factorization, where this chapter's questions lie. We therefore retain the printed sentence, record the defect in the audit, and use the mutually orthogonal meaning whenever a two-headed arrow occurs below.
\end{dcomment}


Quand un couple Q-orthogonal $\Phi \Qo \Psi$ a la propriété de factorisation de Quillen, je l'indique par la barre verticale croisant l'horizontale, ainsi
\[
\Phi \ \Qx{\leftharpoondown}{\rightharpoonup} \ \Psi\ ,
\tag*{(1.4)}
\]
et on définit de même le sens des graphismes
\[
\left\{
\begin{aligned}
&\Phi \quad \Qx{\leftharpoondown}{\rightharpoonup} \quad \Psi\\
&\Phi \quad \Qx{\leftarrow}{\rightharpoonup} \quad \Psi\\
&\Phi \quad \Qx{\leftharpoondown}{\rightarrow} \quad \Psi\\
&\Phi \quad \Qx{\leftarrow}{\rightarrow} \quad \Psi\ .
\end{aligned}
\right.
\tag*{(1.3)} % [sic] (1.3) -> (1.5), number already used above and (1.5) is missing
\]

\newpage

\printedpage{2}
% head: A. Grothendieck : Dérivateurs, Ch. XIX, Édition des Universités de Montpellier II et Paris VII
% foot: Version 11 février 2019

Un couple Q-\textit{orthogonal} Q-\textit{factorisant} $\Phi \Qx{\leftharpoondown}{\rightharpoonup} \Psi$ est appelé un \textit{précouple de Quillen.} S'il est de plus Q-clos, i.e. si on a $\Phi \Qx{\leftarrow}{\rightarrow} \Psi$, on dit que c'est un \textit{couple de Quillen.}

$\Phi$ est dit Q-\textit{saturé à gauche} s'il contient les isomorphismes et est stable par cochangement de base, par (petite) composition transfinie, et par facteurs directs. Dualement $\Psi$ est dit Q-\textit{saturé à droite} s'il

\msp{page 3}

contient les isomorphismes, est stable par changement de base, par composition transfinie, et par facteurs directs. Si $\Phi$ est Q-saturé à gauche, il est stable par (petites) sommes directes, i.e. si les $f_i : A_i \longrightarrow B_i$ ($i \in I$, $I$ petit) sont dans $\Phi$, $f = \coprod f_i : \coprod A_i \longrightarrow \coprod B_i$ aussi, et de même pour (petites) sommes amalgamées finies ou infinies (quand on se place dans un $A\backslash\mathcal{M}$ au lieu de $\mathcal{M}$). On a l'énoncé dual pour $\Psi$, quand $\Psi$ est Q-saturé à droite.

En fait, il y a des notions plus faibles de Q-saturation que celle que je viens de dire : Q-\textit{saturation simple} (à gauche : $\Phi$ contient les isomorphismes, stable par composition binaire, par cochangement de base), Q-\textit{saturation transfinie} (on exige de plus stabilité par composition transfinie), enfin Q-\textit{saturation totale,} ou Q-saturation tout court (quand on exige de plus la stabilité par facteurs directs). On a les implications
\[
\left\{
\begin{aligned}
&\text{Q-clos à gauche}\\
&\Downarrow\\
&\text{Q-saturé (total)}\\
&\Downarrow\\
&\text{Q-saturé transfini}\\
&\Downarrow\\
&\text{Q-saturé simple.}
\end{aligned}
\right.
\tag*{(1.6)}
\]
Pour $\Phi \subset \mathrm{Fl}(\mathcal{M})$, on a sa Q-clôture (à gauche) $\Phi^{\tld}$ (qui est la plus petite partie Q-close à gauche de $\mathrm{Fl}(\mathcal{M})$ contenant $\Phi$) et de même son Q-saturé (total) à gauche, son Q-saturé transfini, son Q-saturé simple

\msp{page 4}

avec les inclusions tautologiques
\[
\Phi \ \subset \
\underbrace{\Phi'}_{\substack{\text{Q-saturé}\\ \text{à gauche simple}}}
\ \subset \
\underbrace{\Phi''}_{\substack{\text{Q-saturé}\\ \text{à gauche transfini}}}
\ \subset \
\underbrace{\Phi'''}_{\substack{\text{Q-saturé}\\ \text{à gauche total}}}
\ \subset \
\underbrace{\Phi^{\tld}}_{\substack{\text{Q-clôture}\\ \text{à gauche}}}
\ ,
\tag*{(1.7)}
\]
et dualement pour $\Psi$. Entre $\Phi$ et $\Phi'$ il y a encore un $\Phi'_0$,
\[
\Phi \ \subset \ \Phi'_0 \ \subset \ \Phi'\ ,
\tag*{(1.8)}
\]
qui est le Q-\textit{saturé pour cochangement de base} (contenant aussi, par définition, tous les isomorphismes). Pour construire tous ces saturés, on procède ainsi :

\dgrblock{dgrb1159}{What cellular closure does and does not supply}
\begin{dcomment}{exposition}{p4:Pour construire tous ces saturés, on procède ainsi}{What cellular closure does and does not supply}
The four closures organize a constructive question: how much of a left lifting class can be built by attaching specified arrows? The elementary operations are chosen because every left lifting class survives them. For a pushout of $i$, a lifting square is equivalent, by the pushout property, to a lifting square against $i$. For a composite $i_2i_1$, solve the lifting problem first against $i_1$ and then against $i_2$. For a continuous ordinal sequence $(X_\alpha)_{\alpha\leq\lambda}$, compatible fillers are constructed at successor stages and are assembled at a limit stage by the colimit property. The resulting map $X_\lambda\to E$ is the required filler for the transfinite composite. Finally, a lifting problem against a retract of an arrow is a retract of a lifting problem against that arrow, so a filler restricts along the retraction.

These arguments prove the implications in (1.6), including the dual assertions, whenever the indicated colimits or limits exist. They also explain why small coproducts need not be added as an independent closure operation. Well-order the indexing set $S$ of a family $(A_s\to B_s)$. Starting with $\coprod_sA_s$, replace the summand $A_s$ by $B_s$ at stage $s$, using a pushout. At a limit stage take the colimit. The final map is $\coprod_sA_s\to\coprod_sB_s$. The same construction in an undercategory treats amalgamated coproducts. This use of a well-order belongs to the usual set-theoretic choice convention.

The formula for $\Phi''$ on the next page rests on a useful flattening argument. A pushout of a transfinite composite is the transfinite composite of the successive pushouts. A transfinite composite of transfinite composites can be indexed by the ordinal sum of the component ordinals. Thus the transfinite composites of pushouts of the original generators already form a class closed under pushout and transfinite composition. Identity maps allow stages at which nothing is attached. This proves the asserted description of $\Phi''$ without first adding an unspecified hierarchy of repeated cellular operations.

What remains at this point is the relation with retracts. Closure under retracts alone gives no factorization, and the notation $\Phi'''$ initially means the smallest class stable under \emph{all} the listed operations. Grothendieck will identify it with the retracts of $\Phi''$ under small-generation hypotheses, and simultaneously identify both with the entire lifting class. The first of these equalities in fact has a general proof, given in the next block; the second is what the small object argument adds. This distinction prevents a positive answer to the elementary closure question from being mistaken for a factorization theorem. A modern reference for the latter is \citet[Theorem~2.1.14 and Corollary~2.1.15]{Hovey1999}.
\end{dcomment}


\newpage

\printedpage{3}
% head: A. Grothendieck : Dérivateurs, Ch. XIX, Édition des Universités de Montpellier II et Paris VII
% foot: Version 11 février 2019

\begin{enumerate}[leftmargin=*]
\item[1)] $\Phi'_0 \quad = \quad$ ensemble des flèches $i : A \longrightarrow B$ s'insérant dans un carré cocartésien
\[
\begin{tikzcd}
A_0 \arrow[r, "i_0"] \arrow[d] & B_0 \arrow[d]\\
A \arrow[r, "i"] & B\ ,
\end{tikzcd}
\]
avec $i_0 \in \Phi$ ou $i_0$ une identité (auquel cas on trouve pour $i$ les isomorphismes).

\item[2)] $\Phi' \quad = \quad$ ensemble des flèches composées multiples
\[
i_n \, i_{n-1} \ \cdots \ i_1\ ,
\]
avec $i_\alpha \in \Phi'_0$.

\item[3)] $\Phi'' \quad = \quad$ ensemble des composés transfinis de flèches $\in \Phi'_0$.
\end{enumerate}

On a envie de poser
\[
\Phi''' \ = \ \Phi^{\prime\prime\natural} \ \overset{\text{déf}}{=} \ \text{ensemble des facteurs directs d'éléments de } \Phi'',
\]
et c'est un fait que $\Phi^{\prime\prime\natural} \subset \Phi'''$ trivialement, mais il n'est par clair en général qu'on ait % [sic] par -> pas, typo
\msp{page 5}
égalité, faute de pouvoir montrer que $\Phi^{\prime\prime\natural}$ est stable, non seulement par facteurs directs, mais aussi par cochangement de base et par composition transfinie. Cependant, rappelons :

\dgrblock{dgrb1160}{Retracts of cellular arrows are closed under transfinite composition}
\begin{dcomment}{later}{p5:égalité, faute de pouvoir montrer}{Retracts of cellular arrows are closed under transfinite composition}
The uncertainty expressed here has a positive answer without a small generating set. If $\mathcal M$ is cocomplete and $\mathcal C$ contains the isomorphisms and is stable under pushout and transfinite composition, then $\operatorname{Retr}(\mathcal C)$ has the same properties. This is \citet[Lemma~2.2]{MakkaiRosicky2014}. We give the explicit compatible-retraction construction, since it isolates exactly the issue Grothendieck raises.

First every retract of a map in $\mathcal C$ can be presented as a retract with its domain fixed. Suppose $f:A\to B$ is a retract of $g:X\to Y$, with maps $a:A\to X$, $b:X\to A$, $c:B\to Y$, $d:Y\to B$ satisfying
\[
ba=1_A,\quad dc=1_B,\quad ga=cf,\quad fb=dg.
\]
Push $g$ out along $b$, obtaining $k:A\to P=A\amalg_XY$ in $\mathcal C$. If $q:Y\to P$ is the other pushout map, put $s=qc:B\to P$. The maps $f:A\to B$ and $d:Y\to B$ induce $r:P\to B$. Then $sf=k$, $rk=f$, and $rs=1_B$. Thus $f$ is a retract of $k$ in $A\backslash\mathcal M$.

Now let $(X_\alpha)_{\alpha\leq\lambda}$ be a continuous chain with each $f_\alpha:X_\alpha\to X_{\alpha+1}$ in $\operatorname{Retr}(\mathcal C)$. Choose fixed-domain presentations
\[
 X_\alpha\xrightarrow{h_\alpha}Z_\alpha,
 \quad s_\alpha':X_{\alpha+1}\to Z_\alpha,
 \quad r_\alpha':Z_\alpha\to X_{\alpha+1},
\]
with $h_\alpha\in\mathcal C$, $s_\alpha'f_\alpha=h_\alpha$, $r_\alpha'h_\alpha=f_\alpha$, and $r_\alpha's_\alpha'=1$. Construct a second continuous chain $Y_\alpha$ and compatible retractions $X_\alpha\xrightarrow{s_\alpha}Y_\alpha\xrightarrow{r_\alpha}X_\alpha$, starting with $Y_0=X_0$. At a successor stage set
\[
Y_{\alpha+1}=Y_\alpha\amalg_{X_\alpha}Z_\alpha,
\]
using $s_\alpha$ and $h_\alpha$ for the pushout. Its structure map $g_\alpha:Y_\alpha\to Y_{\alpha+1}$ lies in $\mathcal C$. Define $s_{\alpha+1}$ by the inclusion of $Z_\alpha$ composed with $s_\alpha'$. Define $r_{\alpha+1}$ by $f_\alpha r_\alpha$ on $Y_\alpha$ and $r_\alpha'$ on $Z_\alpha$; these agree on $X_\alpha$. The equations
\[
r_{\alpha+1}s_{\alpha+1}=1,
\quad g_\alpha s_\alpha=s_{\alpha+1}f_\alpha,
\quad r_{\alpha+1}g_\alpha=f_\alpha r_\alpha
\]
follow directly. At a limit stage take colimits of both the chains and these equations. The composite $X_0\to X_\lambda$ is consequently a retract, with domain fixed, of $Y_0\to Y_\lambda\in\mathcal C$.

Pushout stability follows by pushing out a fixed-domain retraction. Stability under retracts follows by composing retraction diagrams. Taking $\mathcal C=\Phi''$ therefore proves $\Phi'''=\operatorname{Retr}(\Phi'')$ in every cocomplete category, even for a proper class of generators. It does not prove $\Phi'''={}^{\square}(\Phi^{\square})$: that last equality still needs a mechanism, such as Theorem 1.1, for constructing the appropriate factorizations.
\end{dcomment}


\textbf{Théorème 1.1.} \textit{Supposons $\mathcal{M}$ stable par petites $\varinjlim$ , et que ses éléments soient accessibles. Supposons de plus que $\Phi$ soit petit, plus généralement, qu'il existe une partie $\Phi_0 \subset \Phi$ petite telle que $\Phi \subset \Phi_0{}^{\tld}$, donc}
\[
\Psi \ \overset{\text{déf}}{=} \ \Phi_{*} \ = \ (\Phi_0)_{*} \qquad \text{\footnotemark}.
\tag*{(1.1.1)} % footnote mark printed as superscript (1)
\]
\footnotetext{\textbf{N.B.} On a donc $\Phi \Qod \Psi$, a fortiori $\Phi_0 \Qod \Psi$.}
\textit{Sous ces conditions, si $\Phi''_0$ désigne le \textup{Q}-saturé à gauche transfini de $\Phi_0$, alors on a $\Phi''_0 \Qx{\leftharpoondown}{\rightarrow} \Psi$, i.e. toute flèche $f$ de $\mathcal{M}$ se factorise en $f = pi$, avec $i \in \Phi''_0$, $p \in \Psi$\,\footnote{Il y a même factorisation \textit{fonctorielle} en $f$, propriété qui mérite une notation, p.ex. $\Phi''_0 \Qxx{\leftharpoondown}{\rightarrow} \Psi$.}. A fortiori, si $\Phi''$ désigne le \textup{Q}-saturé à gauche transfini de $\Phi$, on a} % footnote mark printed as superscript (2)
\[
\Phi'' \ \Qx{\leftharpoondown}{\rightarrow} \ \Psi\ .
\]

\dgrblock{dgrb1161}{The small object argument and the retract that finishes it}
\begin{dcomment}{exposition}{p5:A fortiori, si $\Phi''$ désigne}{The small object argument and the retract that finishes it}
Theorem 1.1 is the engine of the chapter. It recalls the transfinite construction developed in \citet[Chapter~XV, \S1, Theorem~1]{Grothendieck1990XV}. Its force is not merely that sufficiently many cells exist, but that a \emph{set} of generators admits one common bound on the stages needed to see all its lifting problems. Read ``accessible object'' here as an object whose maps into sufficiently filtered colimits are detected at a stage; for this argument the weaker condition of smallness relative to the relevant cellular sequences suffices. Local smallness is understood in the chosen universe. These hypotheses do not assert that the entire category is locally presentable.

Put $S=\Phi_0$ and let $f:X\to Y$. Choose an infinite regular cardinal $\kappa$ such that every domain of an arrow of $S$ is $\kappa$-small for the sequences being used. Such a common bound exists because $S$ is a set. Starting from $X_0=X$, form at stage $\alpha+1$ the pushout
\[
\begin{tikzcd}[column sep=large]
\displaystyle\coprod_{\sigma} A_\sigma\arrow[r]\arrow[d] & X_\alpha\arrow[d]\\
\displaystyle\coprod_{\sigma} B_\sigma\arrow[r] & X_{\alpha+1},
\end{tikzcd}
\]
where $\sigma$ runs over all commutative squares $s:A_\sigma\to B_\sigma$ against $X_\alpha\to Y$, with $s\in S$. There is a set of such squares, by local smallness and the smallness of $S$. Their lower maps induce $X_{\alpha+1}\to Y$. At a limit stage take the colimit. Choose the length $\lambda$ so that the sequence is sufficiently filtered for every domain in $S$, for example a regular cardinal larger than a common smallness bound.

The composite $i:X\to X_\lambda$ lies in $S''$: attaching a coproduct of cells can be serialized as on the preceding page. The induced $p:X_\lambda\to Y$ has the right lifting property against $S$. Indeed, in a square from $s:A\to B$ to $p$, the map $A\to X_\lambda$ factors through some $X_\alpha$. Together with the fixed map $B\to Y$, it is one of the squares attached at the next stage. The corresponding copy of $B$ in $X_{\alpha+1}$ gives the filler after composition with $X_{\alpha+1}\to X_\lambda$. The equation with $Y$ already holds at stage $\alpha$, since all maps are over $Y$. This proves $f=pi$. Attaching \emph{all} squares makes the construction functorial once colimit choices are fixed; a morphism of arrows sends each indexed square to an indexed square at the next stage.

The retract argument is independent of this construction and deserves to be explicit. Suppose $f:A\to B$ has the left lifting property against $p:E\to B$ and $f=pi$ for $i:A\to E$. A filler in
\[
\begin{tikzcd}
A\arrow[r,"i"]\arrow[d,"f"'] & E\arrow[d,"p"]\\
B\arrow[r,"1_B"']\arrow[ur,dashed,"s"] & B
\end{tikzcd}
\]
satisfies $sf=i$ and $ps=1_B$. Hence $(1_A,s):f\to i$ and $(1_A,p):i\to f$ exhibit $f$ as a retract of $i$ in the arrow category. In particular, for any factorizing orthogonal pair $(L,R)$,
\[
{}^{\square}R=\operatorname{Retr}(L),\qquad
L^{\square}=\operatorname{Retr}(R).
\]
The reverse inclusions follow from stability of lifting under retracts. The double-orthogonal closures equal these classes because orthogonality is unchanged by taking retracts. This proves Proposition 1.3 and, applied to $(S'',S^{\square})$, Corollary 1.2. It also proves Corollary 1.4: factor an arrow of $L$ as a map in $W$ followed by a map in $R$, then use the same square. If $W$ is retract-closed and $(L,R)$ itself factorizes, the converse follows from $L\subseteq W$.
\end{dcomment}


\bigskip

\textbf{Corollaire 1.2.} \textit{Sous les conditions du théorème, on a}
\[
\Phi''' \ \text{(Q-saturé total à gauche de $\Phi$)} \ = \
\underbrace{\Phi^{\prime\prime\natural}}_{\substack{\text{ensemble des}\\ \text{facteurs directs d'éléments}\\ \text{de $\Phi''$}}}
\ = \ \Phi^{\tld}\ .
\tag*{(1.2.1)}
\]
\textit{Plus généralement (en apparence seulement!), si $\Phi_1$ est une partie de \textup{$\mathrm{Fl}(\mathcal{M})$} qui est \textup{Q}-saturée à gauche transfinie, et telle que}
\[
\Phi_0 \ \subset \ \Phi_1 \ \subset \ \Phi^{\tld} \ = \ \Phi_0{}^{\tld}\ ,
\]

% ---- from der-ch19_p006-010.tex
\printedpage{4}
% head: A. Grothendieck : Dérivateurs, Ch. XIX, Édition des Universités de Montpellier II et Paris VII
% foot: Version 11 février 2019

\msp{page 6}
\textit{alors}
\begin{enumerate}[label={}]
\item[a)] \textit{on a $\Phi_1 \Qorthxfull \Psi$ et} % reading resolved: DER_XIX_corrections.csv
\item[b)] $\Phi_1^{\sim} = \Phi^{\sim} = \Phi_1^{\natural}$.
\end{enumerate}
Ce corollaire résulte du théorème en vertu de XV prop. 4.2, dont je rappelle ici l'essentiel :

\textbf{Proposition 1.3.} \textit{Soit $\Phi \Qorthx \Psi$. Alors $\Phi^{\sim} = \Phi^{\natural} = \Psi^{*}$, et dualement on a $\Psi^{-} = \Psi^{\natural} = \Phi_{*}$.}

Notons aussi, dans cet ordre d'idées :

\textbf{Corollaire 1.4.} \textit{Soit $\Phi \Qorth \Psi$, et soit $W$ une partie quelconque de $\mathrm{Fl}(\mathcal{M})$. Pour qu'on ait $\Phi \subset W^{\natural}$, il suffit que toute flèche de $\mathcal{M}$ se factorise en $f = pi$, avec $i \in W$, $p \in \Psi$. (Et cette condition est évidemment aussi nécessaire si $\Phi \Qorthx \Psi$, et $W$ stable par facteurs directs.)}

On dit qu'une partie $\Phi_0$ de $\Phi$ est Q-\textit{génératrice à gauche} de $\Phi$ si elle satisfait la condition
\[
\Phi \ \subset \Phi_0^{\sim}\ , \tag*{(1.9)}
\]
et on dit que $\Phi$ est une partie \textit{spéciale} si $\Phi$ admet une \textit{petite} partie Q-génératrice à gauche. Si $\Phi_0 \subset \Phi$ est petit, et si $\mathcal{M}$ satisfait les conditions du théorème 1.1, alors $\Phi_0$ est une partie Q-génératrice (à gauche) si et seulement si son Q-saturé total (à gauche) $\Phi_0'''$ contient $\Phi$ (condition en apparence plus forte que $\Phi \subset \Phi_0^{\sim}$, puisque en général on a $\Phi_0''' \subset \Phi_0^{\sim}$). Cela signifie donc aussi, via 1.2, que tout \msp{page 7} élément de $\Phi$ est facteur direct d'un composé transfini de flèches déduites de
\[
\Phi_0 \cup (\text{Identités de } \mathcal{M})
\]
par cochangement de base.

Je rappelle encore le sorite $\pm$ tautologique des foncteurs adjoints :

\textbf{Proposition 1.5.} \textit{Soit}
\[
\begin{tikzcd}
\mathcal{M}' \arrow[r, shift right=1, "v^{*}"'] & \mathcal{M} \arrow[l, shift right=1, "v_{!}"']
\end{tikzcd}
\tag*{(1.5.1)}
\]
\textit{un couple de foncteurs adjoints, soient}
\[
\Phi \ \subset\ \mathrm{Fl}(\mathcal{M})\ , \qquad \Psi' \ \subset\ \mathrm{Fl}(\mathcal{M}')\ ,
\]
\textit{alors on a}
\[
(v_{!}(\Phi) \Qorth \Psi') \qquad \Longleftrightarrow \qquad (\Phi \Qorth v^{*}(\Psi'))\ . \tag*{(1.5.2)}
\]

\textbf{Corollaire 1.6.} \textit{Pour toute partie $\Phi$ de} $\mathrm{Fl}(\mathcal{M})$,
\[
(v_{!}(\Phi))_{*} \ =\ (v^{*})^{-1}(\Phi_{*})\ . \tag*{(1.6.1)}
\]

\printedpage{5}
% head: A. Grothendieck : Dérivateurs, Ch. XIX, Édition des Universités de Montpellier II et Paris VII
% foot: Version 11 février 2019

En effet, appliquant 1.5 à $\Phi$ et à $\{f'\}$, où $f' \in \mathrm{Fl}(\mathcal{M}')$, on trouve
\[
f' \in (v_{!}(\Phi))_{*} \qquad \Longleftrightarrow \qquad v^{*}(f') \in \Phi_{*}\ ,
\]
ce qui équivaut à (1.6.1). Dualement :

\textbf{Corollaire 1.7.} \textit{Pour toute partie $\Psi'$ de $\mathcal{M}'$, on a}
\[
(v^{*}(\Psi'))^{*} \ =\ (v_{!})^{-1}(\Psi'^{*}). \tag*{(1.7.1)}
\]

En effet, si $f \in \mathrm{Fl}(\mathcal{M})$, on a par (1.5.2)
\[
f \in (v^{*}(\Psi'))^{*} \qquad \Longleftrightarrow \qquad v_{!}(f) \in \Psi^{*}\ . % [sic] \Psi^{*} -> \Psi'^{*}, cf. (1.7.1)
\]
\textsuperscript{(}\footnote[3]{donner ici 1.9 page 9.}\textsuperscript{)}.

\msp{page 8}
Un couple Q-orthogonal $\Phi \Qorth \Psi$ est dit \textit{spécial} s'il satisfait les conditions suivantes.
\[
\left\{
\begin{array}{ll}
\text{a)} & \parbox[t]{0.7\textwidth}{$\mathcal{M}$ est stable par petites $\varinjlim$ , les objets de $\mathcal{M}$ sont accessibles. (On dit alors que $\mathcal{M}$ est “spéciale”.)} \\[3ex]
\text{b)} & \text{Le couple est clos : } \Phi \clos \Psi. \\[1ex]
\text{c)} & \Phi \text{ est \textit{spéciale.}}
\end{array}
\right.
\tag*{(1.10)}
\]
Cela implique donc, par 1.1, que $\Phi \Qcouple \Psi$ est un couple de Quillen. Si $\mathcal{M}$ satisfait a), les couples Q-orthogonaux spéciaux dans $\mathcal{M}$, ou encore les couples de Quillen spéciaux, correspondent donc biunivoquement aux parties $\Phi$ de $\mathrm{Fl}(\mathcal{M})$ qui sont \textit{spéciales} et closes à gauche (ou ce qui revient au même, Q-saturées à gauche), en associant à une telle $\Phi$ le couple $\Phi \Qcouple \Phi_{*}$.

\dgrblock{dgrb1162}{An adjunction transports the whole lifting problem}
\begin{dcomment}{exposition}{p7}{An adjunction transports the whole lifting problem}
The next passage prepares the construction of diagrams by proving a statement more economical than a transfer theorem. If $L=v_!$ is left adjoint to $R=v^*$, a commutative square from $Li$ to $p'$ corresponds to a commutative square from $i$ to $Rp'$. This correspondence includes the diagonal fillers: a map from the codomain of $Li$ to the domain of $p'$ corresponds to a map from the codomain of $i$ to the domain of $Rp'$, and the two triangle equations correspond by naturality of the adjunction. Thus
\[
 Li\perp p'\quad\Longleftrightarrow\quad i\perp Rp'.
\]
There is no need for unique fillers, and no factorization hypothesis enters this equivalence.

Applying it one arrow at a time gives
\[
 (L\Phi)^{\square}=R^{-1}(\Phi^{\square}),\qquad
 {}^{\square}(R\Psi')=L^{-1}({}^{\square}\Psi').
\]
Here the inverse image means inverse image of a class of arrows, rather than an inverse functor. The unprimed $\Psi^*$ in the proof of Corollary 1.7 is a printed slip: its ambient category forces $\Psi'{}^*$.

The same calculation explains why a left adjoint preserves generation by a set even when it does not literally preserve a chosen list of cellular presentations. If $\Phi\subseteq{}^{\square}(\Phi_0^{\square})$, then $\Phi_0^{\square}\subseteq\Phi^{\square}$; inverse image under $R$ gives
\[
 (L\Phi_0)^{\square}\subseteq(L\Phi)^{\square},
\]
and taking left orthogonals gives $L\Phi\subseteq{}^{\square}((L\Phi_0)^{\square})$. This is Corollary 1.9. It needs neither accessibility nor cocompleteness. The extra hypotheses of Corollary 1.8 are used only after this calculation: the small object argument in the target produces the factorizations and identifies the left class with the retracts of the relative $L\Phi_0$-cell complexes.

This order of reasoning is helpful throughout the chapter. First identify the desired right class as an inverse image. Then identify a small family that defines it. Only then apply a factorization theorem in the target. In particular, the source category's ``specialness'' in Corollary 1.8 is stronger than is required for the transport of generators itself; Grothendieck explicitly observes this when extracting Corollary 1.9. The later transfer of a \emph{model structure} will require an additional condition on weak equivalences. That condition is absent here because a single weak factorization system has no class of weak equivalences to preserve.
\end{dcomment}


Ceci posé :

\marginpar{$\Vert$}% double vertical bar printed in the left margin beside Corollaire 1.8
\textbf{Corollaire 1.8.} \textit{Supposons que $\mathcal{M}$, $\mathcal{M}'$ satisfassent la condition} a) \textit{de} (1.10), \textit{et soit $(\Phi, \Psi)$ un couple de Quillen spécial dans $\mathcal{M}$. Alors $\Bigl(\Phi' = v_{!}(\Phi)^{\sim}, \Psi' = (v^{*})^{-1}(\Psi)\Bigr)$ est un couple de Quillen spécial dans $\mathcal{M}'$.}

En effet, par le corollaire 1.6 on a
\[
\Phi'_{*} \ =\ (v_{!}(\Phi))_{*} \ =\ (v^{*})^{-1}(\Psi) \ =\ \Psi'\ ,
\]
donc $\Phi' \clos \Psi'$. Il reste à prouver que $\Phi'$ est spécial. Soit $\Phi_0$ une partie petite de $\Phi$ telle que $\Phi \subset \Phi_0^{\sim}$, posons $\Phi'_0 \ =\ v_{!}(\Phi_0) \subset \Phi'$, \msp{page 9} il suffit de voir qu'on a
\[
v_{!}(\Phi) \ \subset\ v_{!}(\Phi_0)^{\sim}\ , \tag*{($*$)}
\]
d'où résultera qu'on a $\Phi' \overset{\text{déf}}{=} (v_{!}(\Phi))^{\sim} \subset v_{!}(\Phi_0)^{\sim}$. Mais on a
\[
v_{!}(\Phi_0)^{\sim} \ =\ (v_{!}(\Phi_0)_{*})^{*}\ ,
\]
\printedpage{6}
% head: A. Grothendieck : Dérivateurs, Ch. XIX, Édition des Universités de Montpellier II et Paris VII
% foot: Version 11 février 2019
et
\[
v_{!}(\Phi_0)_{*} \ =\ (v^{*})^{-1}(\underbrace{\Phi_{0\,*}}_{\Psi}) \ =\ (v^{*})^{-1}(\Psi) \ \overset{\text{déf}}{=}\ \Psi'\ ,
\]
donc
\[
v_{!}(\Phi_0)^{\sim} \ =\ \Psi'^{*} \ =\ \Phi' \ =\ v_{!}(\Phi)^{\sim}\ ,
\]
d'où ($*$). Cet argument ne fait pas appel à l'hypothèse “spéciale” sur $\mathcal{M}$, et peut s'expliciter en le

\textbf{Corollaire 1.9.} \textsuperscript{(}\footnote[4]{devrait passer avant 1.8.}\textsuperscript{)} \textit{Sous les conditions générales de} 1.5 \textit{($v_{!} : \mathcal{M} \fl \mathcal{M}'$ étant un foncteur admettant un adjoint à droite), soient $\Phi_0$, $\Phi$ deux parties de} $\mathrm{Fl}(\mathcal{M})$ \textit{telles que $\Phi \subset \Phi_0^{\sim}$ (ou ce qui revient au même, $\Phi^{\sim} \subset \Phi_0^{\sim}$). Alors $v_{!}(\Phi) \subset v_{!}(\Phi_0)^{\sim}$, i.e. $v_{!}(\Phi)^{\sim} \subset v_{!}(\Phi_0)^{\sim}$.}

En effet, cette relation équivaut aussi à $v_{!}(\Phi_0)_{*} \ \subset\ v_{!}(\Phi)_{*}$, c'est-à-dire (par 1.6) $(v^{*})^{-1}(\Phi_{0\,*}) \ \subset\ (v^{*})^{-1}(\Phi_{*})$, et est impliqué par $\Phi_{0\,*} \ \subset\ \Phi_{*}$. Cette dernière d'autre part équivaut à l'hypothèse $\Phi^{\sim} \subset \Phi_0^{\sim}$, q.e.d.

\msp{page 10}
Si $\Psi$ est une partie de $\mathrm{Fl}(\mathcal{M})$, on en déduit pour toute $I$ dans Cat une partie
\[
\left\{
\begin{array}{l}
\Psi(I) \subset \mathrm{Fl}(\mathcal{M}(I)), \quad \text{avec}\ \ \mathcal{M}(I) = \underline{\mathrm{Hom}}(I^{\circ}, \mathcal{M})\ , \\[1ex]
f \in \Psi(I) \iff f(i) \in \Psi \ \ \forall\, i \in \mathrm{Ob}\, I\ .
\end{array}
\right.
\tag*{(1.11)}
\]
Si on introduit la catégorie discrète $I_0$ sous-jacente à $\mathrm{Ob}\, I$, et le foncteur canonique
\[
v_I = v \ :\ I_0 \fl I\ , \tag*{(1.12)}
\]
on trouve un couple de foncteurs adjoints
\[
\begin{tikzcd}[row sep=small]
\mathcal{M}(I) \arrow[r, shift right=1, "v^{*}"'] & \mathcal{M}(I_0) \arrow[l, shift right=1, "v_{!}"'] \\
( & ) \arrow[l, dotted, "v_{*}"]
\end{tikzcd}
\ \simeq\ \mathcal{M}^{I_0}\ ,
\tag*{(1.13)}
\]
et on a
\[
\Psi(I) \ =\ (v^{*})^{-1}(\Psi(I_0))\ .
\]
Supposons maintenant qu'on ait
\[
\left\{
\begin{array}{rcl}
\Psi & = & \Phi_{*} \\[1ex]
\text{d'où}\ \ \Psi(I_0) & = & \Phi(I_0)_{*}
\end{array}
\right.
\tag*{(1.14)}
\]
(vérification immédiate), on conclut alors par 1.6
\[
\Psi(I) \ =\ (v_{!}(\Phi(I_0)))_{*} \ =\ (\Phi\langle I\rangle)_{*}\ , \tag*{(1.15)}
\]
avec
\[
\Phi\langle I\rangle \ \overset{\text{déf}}{=}\ v_{!}(\Phi(I_0))^{\sim} \ =\ \Psi(I)^{*}\ . \tag*{(1.16)}
\]
Si on a un foncteur
\[
u \ :\ I \fl I'\ , \tag*{(1.17)}
\]
d'où un triple de foncteurs adjoints
\printedpage{7}
% head: A. Grothendieck : Dérivateurs, Ch. XIX, Édition des Universités de Montpellier II et Paris VII
% foot: Version 11 février 2019
\msp{page 11}
\[
\left\{
\begin{tikzcd}[column sep=large]
\mathcal{M}(I) \arrow[r, shift left=3, "u_{!}"] \arrow[r, shift right=3, "u_{*}"] & \mathcal{M}(I') \arrow[l, "u^{*}"']
\end{tikzcd}
\right.
,
\tag*{(1.18)}
\]
on aura de façon $\pm$ tautologique
\[
\left\{
\begin{array}{rcl}
u_{!}(\Phi\langle I\rangle) & \subset & \Phi\langle I'\rangle \\
u^{*}(\Psi(I')) & \subset & \Psi(I)\ .
\end{array}
\right.
\tag*{(1.19)}
\]
Le corollaire 1.8 donne :

\dgrblock{dgrb1163}{Free diagrams and the projective construction}
\begin{dcomment}{exposition}{p11:on aura de façon $\pm$ tautologique}{Free diagrams and the projective construction}
The reason for passing through the discrete category $I_0$ is that independent lifting problems can then be solved component by component. The left Kan extension freely inserts the naturality required by the arrows of $I$. Throughout this chapter $\mathcal M(I)=\operatorname{Fun}(I^{\mathrm{op}},\mathcal M)$. For an object $i$ of $I$, the left adjoint to evaluation at $i$ is
\[
 i_!A(j)=\coprod_{I(j,i)} A.
\]
Thus the free diagram on a family $(A_i)$ has value $\coprod_i\coprod_{I(j,i)}A_i$ at $j$. This formula fixes the variance: it is the representable presheaf $I(-,i)$, not $I(i,-)$, that occurs.

Suppose $(\Phi,\Psi)$ is a mutually orthogonal pair and $S\subseteq\Phi$ generates it. The set
\[
 S_I=\{i_!s\mid i\in\operatorname{Ob}I,\ s\in S\}
\]
generates the projective pair. Indeed, $p$ has the right lifting property against $i_!s$ precisely when $p(i)$ has the right lifting property against $s$. Hence
\[
 S_I^{\square}=\Psi(I),\qquad
 {}^{\square}(S_I^{\square})=\Phi\langle I\rangle.
\]
This makes the smallness in Proposition 1.10 concrete. One need not use the much larger class of all objectwise $\Phi$-arrows on the discrete category. If the domain $A$ of $s$ is small, then $i_!A$ is small for pointwise colimits because
\[
\operatorname{Hom}(i_!A,X)=\operatorname{Hom}(A,X(i)).
\]
The same common cardinal as in the original small object argument therefore suffices for these generators.

There is a small scope error in the preliminary equation (1.14). Its equality $\Psi(I_0)=\Phi(I_0)^{\square}$ is valid for the closed left classes used in Proposition 1.10: identities in $\Phi$ allow one to isolate a chosen coordinate, using identity arrows on the domains of the right-hand maps at all other coordinates. It is not valid for an arbitrary $\Phi$ merely satisfying $\Psi=\Phi^{\square}$. For instance, let $\Phi=\{\varnothing\to1\}$ in $\mathrm{Set}$ and let $I_0$ have two objects. The arrow $(\varnothing\to\varnothing,\varnothing\to1)$ has the right lifting property against the sole member $(\varnothing\to1,\varnothing\to1)$ of $\Phi(I_0)$, because no square can be formed in the first coordinate. Its second component is not surjective. Adding identities to $\Phi$ repairs this preliminary assertion; it is automatic in all the subsequent applications to a closed pair.

The source also asserts that the diagram category itself is special. Cocompleteness is pointwise. To check its object-smallness condition, fix $X:I^{\mathrm{op}}\to\mathcal M$ and choose one regular cardinal dominating the ranks of its set of values and the cardinality of the set of arrows of $I$. A natural transformation out of $X$ is a family of component maps subject to one equation for each arrow of $I$. In a sufficiently filtered colimit, all these component maps occur at one stage, and all their naturality equations hold together at a later stage. Equality of two transformations is detected similarly. This proves smallness of $X$. If one uses the stronger filtered-presentability convention for accessible objects, the same argument is expressed by the equalizer of the two products computing the set of natural transformations.

Restriction plainly preserves the objectwise class $\Psi(I')$. Its left adjoint consequently preserves $\Phi\langle I\rangle$, by the preceding lifting correspondence. The dual construction uses objectwise left arrows $\Phi(I)$ and defines $\Psi[I]=\Phi(I)^{\square}$. This is the injective pair; its factorization is a different question. Any displayed right Kan extension is understood where its required limits exist. Cocompleteness and smallness of individual objects alone do not imply completeness: the ordered class of ordinals, with arrows $\alpha\to\beta$ for $\alpha\leq\beta$, has all small suprema and small individual objects, but no terminal object. This observation keeps the existence of the right adjoints separate from the lifting calculations that use them.
\end{dcomment}


\textbf{Proposition 1.10.} \textit{Supposons que $\Phi \Qcouple \Psi$ soit un couple de Quillen} spécial \textit{dans $\mathcal{M}$, i.e. de la forme $(\Phi, \Phi_{*})$, avec $\Phi$ partie Q-close à gauche spéciale dans $\mathcal{M}$ ($\mathcal{M}$ satisfaisant la condition} (1.10. a))\textit{. Alors pour toute $I$ dans} Cat, \textit{$(\Phi\langle I\rangle, \Psi(I))$ est un couple spécial dans $\mathcal{M}(I)$, en particulier il satisfait au théorème de factorisation.}

Considérons la construction duale, qui nous fait introduire, à partir d'un couple $\Phi \clos \Psi$ dans $\mathcal{M}$, le couple
\[
\Phi(I) \ \clos\ \Psi[I]\ , \tag*{(1.20)}
\]
où $\Phi(I)$ est défini par la formule générale (1.11) (où $\Psi$ est remplacé par $\Phi$), et où $\Psi[I]$ est défini soit comme $\Phi(I)_{*}$, soit comme
\[
\Psi[I] \ =\ v_{*}(\Psi(I_0))^{-} \ =\ \Phi(I)_{*}\ . \tag*{(1.21)}
\]

\msp{page 12}
On a alors, pour $u : I \fl I'$
\[
\left\{
\begin{array}{rcl}
u^{*}(\Phi(I')) & \subset & \Phi(I) \\
u_{*}(\Psi[I]) & \subset & \Psi[I']\ .
\end{array}
\right.
\tag*{(1.22)}
\]
Il se pose encore la question si le couple $(\Phi(I), \Psi[I])$ satisfait au théorème de factorisation, sous des hypothèses convenables sur $(\Phi, \Psi)$. Je peux seulement hasarder la

\textbf{Question-conjecture 1.11.} \textit{Supposons que $(\Phi, \Psi)$ soit un couple de Quillen} spécial, \textit{en est-il de même de $(\Phi(I), \Psi[I])$ pour toute $I$ dans} (Cat) \textit{? En d'autres termes, si $\mathcal{M}$ est spéciale, condition} (1.10. a))\textit{, et si $\Phi$ est une partie Q-close à gauche} spéciale \textit{de $\mathcal{M}$, i.e. admettant une petite partie Q-génératrice, est-il vrai que pour toute $I$ dans} Cat, \textit{$\Phi(I)$ est également spéciale ? Je conjecture que oui, en faisant au besoin des hypothèses supplémentaires sur $\mathcal{M}$ (existence d'une filtration cardinale ?). P.ex. si $\Phi_0$ est Q-génératrice dans $\Phi$, et petite, alors $\Phi_0(I)$ est-elle Q-génératrice dans $\Phi(I)$ ?}

\dgrblock{dgrb1164}{The injective question under local presentability}
\begin{dcomment}{later}{p12:Je conjecture que oui, en faisant au besoin}{The injective question under local presentability}
Grothendieck's first question in 1.11 has a precise affirmative answer after the additional hypothesis of local presentability. In a locally presentable category, a weakly saturated class generated by a set remains generated by a set when imposed objectwise on a small diagram category. This is exactly \citet[Lemma~A.2.8.3]{Lurie2009HTT}; its proof is supplied in the more general Lemma~A.3.3.3. The corresponding statement for combinatorial categories, without introducing weak equivalences, is \citet[Corollary~3.4]{MakkaiRosicky2014}, using the numbering of the 2013 arXiv version. These statements apply directly to a cofibrantly generated weak factorization system and prove that $(\Phi(I),\Psi[I])$ is again cofibrantly generated.

The additional hypothesis is substantive. A category being cocomplete with individually accessible objects does not, by itself, specify a set of presentable objects generating the whole category under filtered colimits. The cited theorems therefore settle the locally presentable form of the question, not its entire literal generality. Our literature check did not establish a theorem with exactly that weaker generality. The final sentence of 1.11 asks for more than existence of \emph{some} generating set: it proposes the particular set $\Phi_0(I)$. That proposal fails, as the following calculation shows.
\end{dcomment}


\dgrblock{dgrb1165}{Fixed points detect the missing generators}
\begin{dcomment}{examples}{p12:Je conjecture que oui, en faisant au besoin}{Fixed points detect the missing generators}
The difficulty with the proposed generators can be seen without homotopy theory. Work in $\mathrm{Set}$ with its usual choice convention. The single map $\varnothing\to 1$ generates the weak factorization system of injections and surjections: its right lifting property is surjectivity, injections lift against surjections by choosing preimages on their complements, and the standard factorization is
\[
 X\longrightarrow X\amalg Y\xrightarrow{[f,1_Y]}Y.
\]
Take $I=BG$ for a nontrivial group $G$, say $G=C_2$. The diagram category consists of right $G$-sets. The class $\Phi_0(I)$ contains only the map from the empty $G$-set to the one-point trivial $G$-set. A $G$-map $p:X\to Y$ has its right lifting property if and only if $X^G\to Y^G$ is surjective: a map from the trivial point chooses a fixed point.

Now consider $p:\varnothing\to G$, where $G$ acts on itself freely. Both fixed-point sets are empty, so $p$ passes the proposed test. But the objectwise monomorphism $\varnothing\to G$ does not lift against $p$: in the square whose lower horizontal map is $1_G$, a filler would be a map $G\to\varnothing$. Thus $\Phi_0(I)$ does not Q-generate $\Phi(I)$.

The correct generating family in this example is
\[
 \{\varnothing\to G/H\mid H\leq G\},
\]
with right-coset or left-coset notation adjusted to the chosen action convention. An equivariant injection is obtained by attaching the orbits of its invariant complement; conversely, such attachments are injections, and injections are preserved under retracts. The right class consists of maps surjective on $H$-fixed points for every subgroup $H$. To see this directly, a map from a transitive orbit to $Y$ is a point fixed by the stabilizer of that orbit, and a lift is a fixed preimage. Choosing such lifts orbit by orbit solves an arbitrary equivariant lifting problem against an injection. The single proposed test sees only $H=G$ and misses the free orbit $H=\{1\}$. This fully accounts for the failure; it does not contradict existence of the injective pair.
\end{dcomment}


\dgrblock{dgrb1166}{Exactly when the original singleton still generates}
\begin{dcomment}{development}{p12:Je conjecture que oui, en faisant au besoin}{Exactly when the original singleton still generates}
The preceding example suggests a sharper answer to Grothendieck's last sentence: one can characterize all diagram shapes for which his proposed singleton works. The following elementary result is our development of that sentence. We checked the generator constructions of \citet[Lemma~A.2.8.3]{Lurie2009HTT} and \citet[Corollary~3.4]{MakkaiRosicky2014}; we did not locate this precise characterization there. It is a calculation with constant presheaves, not a claim of priority.

\textbf{Proposition (the constant-generator test).}\ours
Let $I$ be a small category, let $\Delta$ denote the constant-presheaf functor, and let $\Phi_0=\{\varnothing\to1\}$ in $\mathrm{Set}$. Then $\Phi_0(I)$ Q-generates all monomorphisms in $\mathrm{Set}^{I^{\mathrm{op}}}$ if and only if $I$ is empty or equivalent to the terminal category.

\textit{Proof.} The class $\Phi_0(I)$ consists of the single map $\varnothing\to\Delta1$: all its component domains and codomains are prescribed, and their transition maps are forced. A pushout of a coproduct of copies of this map adds a constant summand. Consequently every relative cell complex starting from the initial presheaf is $\varnothing\to\Delta S$ for some set $S$. This includes transfinite cell complexes, since colimits of the successive constant summands are computed pointwise. The small object argument applies, since the initial presheaf is small. Therefore every object $P$ for which $\varnothing\to P$ belongs to the generated left class is a retract of some $\Delta S$.

Suppose now that this left class contains every monomorphism and that $I$ is nonempty. For every $i\in I$, the representable $h_i=I(-,i)$ is then a retract of a constant presheaf. Choose natural transformations
\[
 h_i\xrightarrow{a}\Delta S_i\xrightarrow{b}h_i,
 \qquad ba=1_{h_i}.
\]
For any arrow $\alpha:j\to k$ of $I$, naturality gives
\[
 a_jh_i(\alpha)=a_k,\qquad
 h_i(\alpha)b_k=b_j.
\]
It follows that $h_i(\alpha)$ has inverse $b_ka_j$: composing in either order and using the two displayed equations yields the appropriate identity. Since this holds for every $i$, every map
\[
 I(k,i)\longrightarrow I(j,i),\qquad t\longmapsto t\alpha
\]
is bijective. This forces $\alpha$ to be invertible. Explicitly, surjectivity with $i=j$ gives $\beta:k\to j$ with $\beta\alpha=1_j$; injectivity with $i=k$ then gives $\alpha\beta=1_k$, since $(\alpha\beta)\alpha=\alpha=1_k\alpha$. Thus $I$ is a groupoid.

Moreover, $S_i$ is nonempty because $h_i(i)$ contains $1_i$ and maps into $S_i$. For every $j$, the map $b_j:S_i\to I(j,i)$ therefore shows that $I(j,i)$ is nonempty. The groupoid is connected. Finally, an endomorphism $\gamma:i\to i$ acts trivially on $h_i(i)$: naturality gives $a_ih_i(\gamma)=a_i$, and $a_i$ is injective because it has a left inverse. Applying this identity to $1_i$ gives $\gamma=1_i$. Every isotropy group is therefore trivial. A connected groupoid with trivial isotropy has exactly one arrow between any two objects, so it is equivalent to the terminal category.

Conversely, if $I$ is equivalent to the terminal category, evaluation at any object is an equivalence $\mathrm{Set}^{I^{\mathrm{op}}}\simeq\mathrm{Set}$ and takes $\varnothing\to\Delta1$ to $\varnothing\to1$. Equivalences preserve and reflect lifting properties and monomorphisms, so the stated generation follows from the set calculation. If $I$ is empty, the presheaf category has one object and one arrow; both the generated left class and its class of monomorphisms consist of that arrow. This proves all cases.\hfill$\square$

The statement concerns the proposed \emph{unchanged} generator. Even a discrete category with two objects fails the test: it asks one attachment to add a point simultaneously in both coordinates. The affirmative injective-generation theorem supplies a different set, large enough to express the independent and compatible attachments required by the chosen shape.
\end{dcomment}


\textbf{Proposition 1.12} \textsuperscript{(}\footnote[5]{cf. XVII 6.4 (p. 117), mais où les notations pour $\Phi\langle I\rangle$, $\Psi(I)$ sont inverses.}\textsuperscript{)}. \textit{Soit $\Phi \clos \Psi$ dans $\mathcal{M}$, soit $u : I \fl I'$ une flèche dans} Cat, \textit{d'où}
\[
\left\{
\begin{tikzcd}[column sep=large]
\mathcal{M}(I) \arrow[r, shift left=3, "u_{!}"] \arrow[r, shift right=3, "u_{*}"] & \mathcal{M}(I') \arrow[l, "u^{*}"']
\end{tikzcd}
\right.
.
\tag*{(1.12.1)}
\]
\printedpage{8}
% head: A. Grothendieck : Dérivateurs, Ch. XIX, Édition des Universités de Montpellier II et Paris VII
% foot: Version 11 février 2019
\textit{On a alors} (1.19), \textit{i.e.}
\[
\left\{
\begin{array}{rcl}
u_{!}(\Phi\langle I\rangle) & \subset & \Phi\langle I'\rangle \\
u^{*}(\Psi(I')) & \subset & \Psi(I)\ ,
\end{array}
\right.
\tag*{(1.13.1)}
\]
\textit{et de plus l'équivalence}
\[
u^{*}(\Phi(I')) \subset \Phi\langle I'\rangle \qquad \Longleftrightarrow \qquad u_{*}(\Psi(I)) \subset \Psi(I')\ . \tag*{(1.13.2)}
\]
\textit{Ces relations équivalentes sont satisfaites dans chacun des trois cas}~\textsuperscript{(}\footnote[6]{\textbf{N.B.} On a $1^{\circ} \Longrightarrow 3^{\circ}$, $2^{\circ} \Longrightarrow 3^{\circ}$.}\textsuperscript{)} \textit{:}
\begin{enumerate}[label={}]
\item[1°)] \textit{$u$ est coétale.}
\item[2°)] \textit{$I$ est discrète.}
\item[3°)] $\forall\, i' \in \mathrm{Ob}\, I'$\textit{, toute composante connexe de $I/i'$ admet un objet final.}
\end{enumerate}
\msp{page 13}

\dgrblock{dgrb1167}{The comma-category test for preserving projective arrows}
\begin{dcomment}{exposition}{p12:Ces relations équivalentes sont satisfaites}{The comma-category test for preserving projective arrows}
Proposition 1.12 identifies shapes for which restriction preserves the projective left class. This is a useful exception to the usual asymmetry: restriction always preserves objectwise right arrows, but it need not preserve the left arrows defined by a lifting property against them.

The printed formula (1.13.2) contains two slips that have to be corrected in interpreting the statement. Its left side should read
\[
 u^*\bigl(\Phi\langle I'\rangle\bigr)
       \subseteq\Phi\langle I\rangle,
 \tag{*}
\]
not $u^*(\Phi(I'))\subseteq\Phi\langle I'\rangle$. The latter has the wrong ambient category on the right and the wrong, objectwise, class on the left. The intended equivalence follows by applying the lifting correspondence to the adjunction $u^*\dashv u_*$:
\[
\begin{aligned}
u^*\Phi\langle I'\rangle\subseteq{}^{\square}\Psi(I)
&\Longleftrightarrow
\Phi\langle I'\rangle\perp u_*\Psi(I)\\
&\Longleftrightarrow
u_*\Psi(I)\subseteq\Psi(I').
\end{aligned}
\]
The source's established mutual-orthogonality identities are exactly what the last step needs.

To prove the sufficient condition, let $K_{i'}=(u\downarrow i')$: its objects are pairs $(i,u(i)\to i')$. For presheaves, the pointwise right Kan extension is
\[
 (u_*X)(i')=\lim_{K_{i'}^{\mathrm{op}}}X(i).
\]
Suppose each connected component $c$ of $K_{i'}$ has a terminal object $t_c=(i_c,u(i_c)\to i')$. Then $t_c$ is initial in the opposite category, and a limit over that component is evaluation at $i_c$. Therefore
\[
 (u_*X)(i')\cong\prod_cX(i_c).
\]
For an objectwise arrow in $\Psi(I)$, the induced arrow is a product of arrows in $\Psi$. Any right lifting class is closed under products: project a lifting square to each factor, choose its fillers, and combine them by the product property. The empty product gives an identity of a terminal object. Thus $u_*$ preserves $\Psi(I)$ and (*) follows. This supplies the proof left implicit in the statement.

If $I$ is discrete, every $K_{i'}$ is discrete and each nonempty component is a singleton. For a coétale functor in the source's convention, lifting a map into $i'$ identifies a terminal lift in each component of $K_{i'}$. These are the two advertised special cases of the third condition. All formulas use the relevant products or right Kan extensions where they exist. No factorization theorem and no accessibility hypothesis are needed for this argument; those enter only when one wants the orthogonal pairs themselves to give replacement factorizations.
\end{dcomment}


\section*{2\quad Rappels sur les triples de Quillen clos et les triples de Quillen spéciaux}

Soit $(\mathcal{M}, W)$ une \textit{catégorie de modèles,} i.e. $\mathcal{M}$ munie d'un localiseur $W \subset \mathrm{Fl}(\mathcal{M})$, i.e. une partie satisfaisant
\[
\parbox{0.75\textwidth}{$W$ contient les isomorphismes, et si $f$, $g$ sont des flèches composables, si deux parmi $f$, $g$, $gf$ sont dans $W$, la troisième aussi \textit{(saturation faible).}}
\tag*{(2.1)}
\]
Je rappelle que la donnée d'un \textit{triple de Quillen clos} $(W, C, F)$\textit{,} ou d'une structure de Quillen close, sur $\mathcal{M}$, ayant $W$ comme ensemble de quasi-isomorphismes, revient à la donnée de deux couples de Quillen
\[
C \Qcouple TF\ , \qquad TC \ \Qcouple\ F\ , \tag*{(2.2)}
\]
satisfaisant
\[
TC \ \subset\ C \cap W\ , \qquad TF \ \subset\ F \cap W\ , \tag*{(2.3)}
\]
plus une des trois conditions équivalentes :
\[
\left\{
\begin{array}{lrcll}
\text{a)} & TC & = & C \cap W & \qquad (\text{i.e. } C \cap W \subset TC) \\[1ex]
\text{b)} & TF & = & F \cap W & \qquad (\text{i.e. } F \cap W \subset TF) \\[1ex]
\text{c)} & W & = & TF \circ TC\ . &
\end{array}
\right.
\tag*{(2.4)}
\]
Rappelons aussi que ceci implique que $W$ \textit{est fortement saturée,} et a fortiori que $W = W^{\natural}$, i.e. $W$ stable par facteurs directs.

Rappelons aussi qu'en vertu de (1.4), on a :

% ---- from der-ch19_p011-015.tex
\printedpage{9}
% head: A. Grothendieck : Dérivateurs, Ch. XIX, Édition des Universités de Montpellier II et Paris VII
% foot: Version 11 février 2019

\msp{page 14}

\dgrblock{dgrb1168}{How the two factorizations recognize weak equivalences}
\begin{dcomment}{exposition}{p13}{How the two factorizations recognize weak equivalences}
The purpose of the second section is to add compatibility with a localizer to the machinery already established for one weak factorization system. Grothendieck's ``category of models'' at (2.1) means a category with a specified two-out-of-three class. It should not be silently identified with a modern model category. The model structure enters through the two factorizing pairs
\[
 (C,TF),\qquad (TC,F),
\]
with $TC\subseteq C$ and $TF\subseteq F$. The letters $T$ stand for the arrows intended to be trivial, and that intention has to be checked against $W$. For the equivalence in (2.4), the source explicitly assumes $TC\subseteq C\cap W$ and $TF\subseteq F\cap W$ before asking whether these inclusions are equalities.

The equivalence of the three conditions (2.4) is a short but instructive use of both factorizations. Suppose first that $TC=C\cap W$. If $f\in W$, factor $f=pi$ using $(C,TF)$. Since $p\in W$, two-out-of-three gives $i\in C\cap W=TC$, hence $W\subseteq TF\circ TC$. The reverse inclusion follows because both factors belong to $W$. Conversely, if $W=TF\circ TC$ and $f\in C\cap W$, choose such a factorization. Since $f\perp p$, the retract argument makes $f$ a retract of $i\in TC$, so $f\in TC$. This proves equivalence of (a) and (c). Using the other factorization and the dual retract argument proves equivalence with (b).

Formula (2.5) is another application of the same idea. To test $TC\subseteq W$ when $W$ is retract-closed, it suffices to factor each arrow of $TC$ as a $W$-arrow followed by an $F$-arrow; the left lifting property makes the original arrow a retract of the first factor. If $TC\subseteq W$ already holds, the $(TC,F)$-factorization supplies this for all arrows. The dual test is a factorization $f=pi$ with $i\in C$ and $p\in W$. The printed last clause ``$i\in W$, $f\in C$'' is incorrect and is retained as an edition slip, with its intended dual stated here.

The ``special'' condition adds small generators for both left classes, together with the ambient smallness hypotheses that make their small object arguments work. It is closely related to cofibrant generation in contemporary terminology, but stronger ambient hypotheses are explicitly included here. It is not synonymous with ``combinatorial'': local presentability is not part of (2.6). The recalled strong saturation of $W$ uses the ambient model-category convention: \citet[Chapter~XIII, \S3, manuscript pp.~27--29]{Grothendieck1990XIII} explicitly assumes M0, existence of finite limits and finite colimits, and invokes Quillen's saturation result. Strong saturation means that precisely the arrows of $W$ become invertible in $\mathcal M[W^{-1}]$; retract-closedness follows since a retract of an isomorphism is an isomorphism. A modern reference is \citet[Theorem~1.2.10(iv)]{Hovey1999}. The audit records this as a recalled theorem with its inherited ambient hypotheses, not as a gap in the present proof.
\end{dcomment}


\begin{itemize}[leftmargin=4em]
\item[(2.5)] Étant donnés les deux couples (2.2), avec $TC \subset C$ (équivalent à $TF \subset F$), pour qu'on ait la première relation (2.3), i.e. $TC \subset W$, il faut, et il suffit si $W = W^{\natural}$, que toute flèche $f$ de $\M$ (ou seulement toute flèche de $TC$) se factorise en $f = pi$, avec $i \in W$, $p \in F$. Dualement, pour qu'on ait $TF \subset W$, il faut et il suffit si $W = W^{\natural}$, que toute flèche $f$ de $\M$ (ou seulement toute flèche de $TF$) se factorise en $f = pi$, avec $i \in W$, $f \in C$. % [sic] i \in W, f \in C -> i \in C, p \in W (?), dual of the preceding factorisation with p \in F
\end{itemize}

On dit que la structure de Quillen $(W,C,F)$ est \textit{spéciale} si les deux couples de Quillen (2.3) sont spéciaux, i.e. si les conditions suivantes sont satisfaites.

\begin{itemize}[leftmargin=4em]
\item[(2.6)]
\begin{itemize}
\item[a)] $\M$ est stable par petites $\varinjlim$, et ses éléments sont accessibles ($\M$ “spéciale”).
\item[b)] Il existe des \textit{petites} parties $C_0 \subset C$, $TC_0 \subset TC$ qui soient Q-génératrices à gauche, i.e. telles que l'on ait
\[
\left\{
\begin{array}{rclcrclcrcl}
C & \subset & C_0^{\sim}\ , & \text{i.e.} & C & = & C_0^{\sim}\ , & \text{i.e.} & TF & = & C_{0*}\ ;\\
TC & \subset & (TC_0)^{\sim}\ , & \text{i.e.} & TC & = & (TC_0)^{\sim}\ , & \text{i.e.} & F & = & (TC_0)_*\ .
\end{array}\right. % reading resolved: DER_XIX_corrections.csv
\]
\end{itemize}
\end{itemize}

Ainsi, la donnée d'un tel triple de Quillen \textit{spécial} dans $\M$, quand $\M$ satisfait (2.6. a), revient à la donnée des deux parties
\msp{page 15}
\[
TC \ \subset\ C
\tag*{(2.7)}
\]
de $\Fl(\M)$, satisfaisant les conditions :
\[
\left\{
\begin{array}{ll}
C,\ TC\ \text{“spéciales”} &
\left\{\begin{array}{ll}
\text{a)} & \text{$C$, $TC$ sont Q-closes à gauche.}\\[1ex]
\text{b)} & \parbox[t]{0.45\textwidth}{$C$, $TC$ sont Q-engendrées à gauche par des \textit{petites} parties $C_0$, $TC_0$ de $C$, $TC$ respectivement.}
\end{array}\right.\\[5ex]
 & \quad\ \begin{array}{ll}
\text{c)} & TC \subset W,\ TF \overset{\text{déf}}{=} C_* \subset W.\\[1ex]
\text{d)} & \parbox[t]{0.45\textwidth}{On a l'une des trois conditions équivalentes supplémentaires de (2.4), i.e. :}\\[5ex]
 & \begin{array}{lrcllrcl}
\alpha) & TC & = & C \cap W\ , & \text{i.e.} & C \cap W & \subset & TC\ .\\
\beta) & TF & = & F \cap W\ , & \text{i.e.} & F \cap W & \subset & TF\ .\\
\gamma) & W & = & TF \circ TC\ . & & & &
\end{array}
\end{array}
\end{array}\right.
\tag*{(2.8)}
\]

\dgrblock{dgrb1169}{The extra condition in transferring a model structure}
\begin{dcomment}{exposition}{p15:de $\Fl(\M)$, satisfaisant les conditions}{The extra condition in transferring a model structure}
Before Proposition 2.1, it helps to identify what can go wrong. The adjunction transfers the two lifting relations formally, and small generation can provide both factorizations. Nevertheless, a map built from the transferred trivial generators might cease to be a weak equivalence after applying the right adjoint. The condition $TC'\subseteq W'$ rules out exactly that possibility.

Write $L=v_!$ and $R=v^*$. With
\[
 W'=R^{-1}W,\quad F'=R^{-1}F,\quad
 TF'=R^{-1}TF,
\]
the adjunction identifies
\[
 C'={}^{\square}TF',\qquad TC'={}^{\square}F'.
\]
Thus both pairs are already mutually orthogonal, although not yet known to factorize. Since $TC\subseteq C$, monotonicity gives $TC'\subseteq C'$, and inverse image gives $TF'\subseteq F'$. The class $W'$ satisfies two-out-of-three because $R$ is a functor. Once the two factorizations exist, the equation $TF'=F'\cap W'$ is built into the definitions; the previous block therefore reduces the entire model-structure question to $TC'\subseteq W'$.

The equivalent condition (2.1.5) has a useful practical interpretation: factor every target arrow into a weak equivalence followed by a member of $F'$, without insisting in advance that the first arrow is in $TC'$. If the arrow being factored already belongs to $TC'$, its left lifting property makes it a retract of that weak equivalence. Since $W'$ is retract-closed, it belongs to $W'$. Conversely, a $(TC',F')$-factorization works once $TC'\subseteq W'$ is known.

The convenient stronger condition is $RL(TC)\subseteq TC$. It tests the images of the original trivial left arrows before taking their left closure. Cases A and B show two mechanisms by which that condition passes to the closure. Case C, completed on manuscript page 19, instead transports a factorization of $Rf'$ through a fully faithful adjunction. This is why the three cases have different hypotheses. A right adjoint generally does not preserve pushouts, so the hypothesis of Case B must not be suppressed in a purported general transfer theorem.
\end{dcomment}


\printedpage{10}
% head: A. Grothendieck : Dérivateurs, Ch. XIX, Édition des Universités de Montpellier II et Paris VII
% foot: Version 11 février 2019

\textbf{Proposition 2.1.} \textit{Soit}
\[
\begin{tikzcd}
\M' \arrow[r, shift right, "v^*"'] & \M \arrow[l, shift right, "v_!"']
\end{tikzcd}
\tag*{(2.11)} % [sic] (2.11) -> (2.1.1), numbering sequence (2.1.2), (2.1.3), ...
\]
\textit{un couple de foncteurs adjoints, $\M$ étant muni d'une structure de Quillen close} $(W,C,F)$. \textit{Posons}
\[
\left\{
\begin{array}{rcl}
W' & = & (v^*)^{-1}(W)\\
F' & = & (v^*)^{-1}(F)\\
TF' & = & F' \cap W' \ = \ (v^*)^{-1}(TF)\\
C' & = & v_!(C)^{\sim}\\
TC' & = & v_!(TC)^{\sim}\ .
\end{array}\right.
\tag*{(2.1.2)}
\]

\begin{itemize}
\item[1°)] \textit{On a alors}
\[
\left\{
\begin{array}{rclcrcl}
C' & \longleftrightarrow & TF'\ , & & TC' & \longleftrightarrow & F'\ ,\\
TC' & \subset & C'\ , & & TF' & \subset & F'\ .
\end{array}\right.
\tag*{(2.1.3)}
\]

\textit{Pour que les données} (2.1.2) \textit{dans $\M'$ y correspondent à un triple de Quillen $(W',C',F')$, il faut et il suffit que les conditions suivantes soient satisfaites :}
\msp{page 16}
\[
\left\{
\begin{array}{ll}
\text{a)} & \parbox[t]{0.7\textwidth}{\textit{Factorisation pour $C' \longleftrightarrow TF'$ et $TC' \longleftrightarrow F'$, i.e. les relations $C' \Qcouple TF'$ et $TC' \Qcouple F'$,}}\\[4ex]
\text{b)} & TC' \subset W',
\end{array}\right. % reading resolved: DER_XIX_corrections.csv
\tag*{(2.1.4)}
\]
\textit{ou encore, au lieu de} b)\textit{, la condition équivalente} (en présence de a)) :
\[
\left\{
\begin{array}{ll}
\text{b$'$)} & \parbox[t]{0.7\textwidth}{\textit{Toute flèche $f'$ dans $\M'$ (non seulement $f' \in TC'$) se factorise en $f' = p'i'$, avec $i' \in W'$, $p' \in F'$.}}
\end{array}\right.
\tag*{(2.1.5)}
\]
\item[2°)] \textit{Supposons} (2.1.4. a) \textit{satisfaite (factorisation). Alors, pour que $(W',C',F')$ soit un triple de Quillen clos, il} suffit \textit{que l'on ait}
\[
\text{b$''$)}\qquad \boxed{\ v^*v_!(TC)\ \subset\ TC\ }\ ,
\tag*{(2.1.6)}
\]
\textit{et que de plus on soit dans l'un des trois cas suivants :}
\begin{itemize}
\item[A)] \textit{$v^*$ admet un adjoint à droite $v_*$.}
\item[B)] \textit{$v^*$ commute aux petites} $\varinjlim$\textit{, $(W,C,F)$ est spécial, et $\M'$ est spéciale, i.e. satisfait à la condition} (2.6. a).
\item[C)] \textit{$v^*$ est pleinement fidèle, et $v^*v_!(F) \subset F$.}
\end{itemize}
\end{itemize}
Notons tout de suite le

\textbf{Corollaire 2.2.} \textit{Supposons que $(W,C,F)$ soit un triple de Quillen spécial, que $\M'$ soit spéciale, que $v^*$ commute aux petites} $\varinjlim$\textit{, enfin que l'on ait} (2.1.6)\textit{, i.e. $v^*v_!(TC) \subset TC$. Alors $(W',C',F')$ est un triple de Quillen spécial.}

\printedpage{11}
% head: A. Grothendieck : Dérivateurs, Ch. XIX, Édition des Universités de Montpellier II et Paris VII
% foot: Version 11 février 2019

\msp{page 17}
En effet, en vertu du 2.1. 2°, cas B, pour voir que $(W',C',F')$ soit un triple de Quillen clos, il est à vérifier (2.1.6. b$''$). Mais cela résulte de 1.8, appliqué aux deux couples de Quillen spéciaux $(C,TF)$, $(TC,F)$ dans $\M$ : les couples $(C',TC')$ et $(TC',F')$ sont des couples de Quillen spéciaux, ce qui prouve en même temps que le triple $(W',C',F')$ est spécial, q.e.d. % [sic] (C',TC') -> (C',TF'), image of the couple (C,TF)

Démontrons 2.1. Les premières deux relations (2.1.3) sont contenues dans 1.6, les relations $TC' \subset C'$, $TF' \subset F'$ sont tautologiques. Le critère (2.1.4) pour que $(W',C',F')$ soit un triple de Quillen clos, est contenu dans le rappel de la page 13, compte tenu que l'on sait déjà que $TF' = F' \cap W'$ (relation (2.4. b)). Que la condition (2.1.4. b), en présence de (2.1.4. a), soit équivalente à la condition (2.1.5. b$'$), est contenu dans la première assertion (2.5), en notant que l'on a $W' = W'^{\natural}$ (ce qui résulte de la définition de $W'$ et de $W = W^{\natural}$).

Il reste à prouver 2°), i.e. que la condition b) est satisfaite moyennant a) et b$''$) (2.1.6) dans chacun des trois cas A), B), C).

\textit{Cas A.} Par définition de $W' = (v^*)^{-1}(W)$, pour prouver que $TC' \subset W'$, il suffit de prouver
\[
v^*(TC')\ \subset\ W'\ ,  % [sic] W' -> W, v^*(TC') lies in M and W' = (v^*)^{-1}(W)
\]
et a fortiori, il suffit de prouver
\[
v^*(TC')\ \subset\ TC\ ,\qquad\quad \text{i.e. } v^*(TC') \harp F\ .
\tag*{(2.1.7)} % reading resolved: DER_XIX_corrections.csv
\]
Or par 1.5, appliqué au couple de foncteurs
\msp{page 18}
adjoints $(v^*,v_*)$, cette relation équivaut à
\[
TC' \harp v_*(F)\ .
\]
Comme $TC' \overset{\text{déf}}{=} v_!(TC)^{\sim}$, ceci équivaut à
\[
v_!(TC) \harp v_*(F)\ ,
\]
ou encore (en réappliquant 1.5 à $(v_!(TC),F)$ au lieu de $(TC',F)$), à
\[
v^*v_!(TC) \harp F\ ,\qquad\quad \text{i.e. } v^*v_!(TC) \subset \underbrace{TC}_{=\,F^*}\ ,
\]
or c'est là l'hypothèse b$''$).

\textit{Cas B.} Il suffit encore de prouver 2.1.7. Or, du fait que $v_!(TC)$ admet un petit ensemble Q-générateur à gauche, $v_!(TC)^{\sim}$ est la clôture de $v_!(TC)$ pour les opérations de cochangement de base, de composition transfinie, de facteurs directs. Comme $v^*$ commute à toutes ces opérations (comme il commute aux petites $\varinjlim$), et comme $TC$ dans $\Fl(\M)$ est stable par ces opérations, on voit donc que pour montrer que $v^*$ applique $v_!(TC)^{\sim}$ dans $TC$, il suffit de voir qu'il applique $v_!(TC)$ dans $TC$, ce qui est encore l'hypothèse b$''$).

\dgrblock{dgrb1170}{Two ways of passing from generators to their closure}
\begin{dcomment}{exposition}{p18:Comme $v^*$ commute à toutes ces opérations}{Two ways of passing from generators to their closure}
The proofs of Cases A and B establish the stronger conclusion $R(TC')\subseteq TC$. They are worth comparing because one uses only lifting properties, whereas the other uses the cellular description supplied by small generation.

In Case A, $R$ has a right adjoint $V=v_*$. For each $p\in F$, the condition $RL(TC)\perp p$ is equivalent to $L(TC)\perp Vp$. Taking the left double-orthogonal closure of $L(TC)$ does not change which arrows have the right lifting property against it; hence $TC'\perp Vp$. Applying the adjunction once more gives $R(TC')\perp p$. Since $TC={}^{\square}F$, this proves the desired containment. This argument does not require $R$ to preserve any colimit and does not require the ambient categories to be locally presentable.

In Case B, the small object argument identifies $TC'$ with the retracts of cellular composites of a small family inside $L(TC)$. The functor $R$ preserves pushouts and transfinite composites by the stated colimit hypothesis; every functor preserves retract diagrams. Since $TC$ is closed under these operations and $RL(TC)\subseteq TC$, it follows that $R(TC')\subseteq TC$. One may equally use the full weakly saturated closure of $L(TC)$: Corollaries 1.8--1.9 show that its small generating family has the same left closure.

Corollary 2.2 combines this argument with Corollary 1.8, which gives both target factorizations and their small generators. Its printed proof contains two local slips. The first transferred pair is $(C',TF')$, not $(C',TC')$. Moreover, the remaining condition supplied by Corollary 1.8 is factorization (2.1.4.a), not (2.1.6.b$''$), which is already an explicit hypothesis of Corollary 2.2. In the subsequent Case A proof, $R(TC')$ is a class in $\mathcal M$, so the intermediate target must be $W$, not $W'$. With these corrections of interpretation, the proof establishes the corollary exactly as stated.
\end{dcomment}


\textit{Cas C.} On va vérifier b) sous la forme équivalente b$'$). Soit donc $f' \in \Fl(\M')$, $f' : X' \longrightarrow Y'$,
\printedpage{12}
% head: A. Grothendieck : Dérivateurs, Ch. XIX, Édition des Universités de Montpellier II et Paris VII
% foot: Version 11 février 2019
\msp{page 19}
posons
\[
f\ =\ v^*(f'),
\]
et factorisons $f$ en $pi$, avec $i \in TC$, $p \in F$, d'où
\[
v_!v^*(f') = v_!(f) = v_!(p)v_!(i)\ .
\]
Mais $v^*$ étant pleinement fidèle, $v_!v^* \xrightarrow{\ \sim\ } \mathrm{id}$, donc on trouve une factorisation de $f'$ en $v_!(p)v_!(i)$. Il suffit donc de voir que
\[
v_!(p)\ \in\ F'\ ,\qquad\quad v_!(i)\ \in\ W'\ ,
\]
ce qui équivaut à
\[
v^*v_!(p) \in F\ ,\qquad\quad v^*v_!(i) \in W\ ,
\]
et résulte des hypothèses que $v^*v_!(F) \subset F$ et $v^*v_!(TC) \subset TC \subset W$, q.e.d.

\textbf{N.B.} J'ai noté les cas A) et C) par acquis de conscience, sans application précise en vue. Le cas le plus utile sera justiciable du cas B), et mérite d'être réexplicité en un

\dgrblock{dgrb1171}{Free diagrams and the projective model structure}
\begin{dcomment}{exposition}{p19:J'ai noté les cas A) et C) par acquis de conscience}{Free diagrams and the projective model structure}
The concluding Case~C can first be made completely explicit. Write $L=v_!$, $R=v^*$ and $\varepsilon:LR\to1$ for the counit. Since $R$ is fully faithful, $\varepsilon$ is an isomorphism. For $f':X'\to Y'$, choose in $\mathcal M$ a factorization
\[
 R X'\xrightarrow{i}Z\xrightarrow{p}R Y',
 \qquad i\in TC,\quad p\in F.
\]
Define
\[
 i'=L(i)\varepsilon_{X'}^{-1}:X'\longrightarrow LZ,
 \qquad p'=\varepsilon_{Y'}L(p):LZ\longrightarrow Y'.
\]
Naturality of the counit gives $p'i'=f'$. Moreover $R(i')=RL(i)R(\varepsilon_{X'}^{-1})$ lies in $TC$, since $RL(TC)\subseteq TC$ and $TC$ is stable under composition with isomorphisms. Likewise $R(p')=R(\varepsilon_{Y'})RL(p)$ lies in $F$, using the other hypothesis $RL(F)\subseteq F$. Thus $i'\in W'=R^{-1}W$ and $p'\in F'=R^{-1}F$, proving the factorization condition (2.1.5). The two factorizing pairs required separately in Proposition~2.1 remain part of its hypotheses; this argument supplies their compatibility with weak equivalences.

Grothendieck now puts the transfer argument to its principal use: a homotopy theory of coefficients should supply a homotopy theory of diagrams of every small shape. The weak equivalences must be tested at the objects of the indexing category, since this is what allows a diagram to retain its homotopy coherence while its individual values are replaced. The freedom lies in choosing the auxiliary cofibrations and fibrations. For homotopy left Kan extensions, the useful choice fixes the fibrations componentwise and lets lifting determine the cofibrations.

Write $\mathcal M(I)=\operatorname{Fun}(I^{\mathrm{op}},\mathcal M)$, as in the source, and use $L^{\square}$ and ${}^{\square}R$ for right and left lifting classes. The intended formulas are
\[
 C_{\mathrm{proj}}(I)={}^{\square}(F(I)\cap W(I)),
 \qquad
 TC_{\mathrm{proj}}(I)={}^{\square}F(I).
\]
The first equality in Theorem 2.3 as printed has $F(I)$ where $TF(I)$ is required. This is visible in the edition and is therefore retained in the transcription. Its correction follows already from (1.16), applied to $(C,TF)$. With the printed reading, the alleged cofibrations would be precisely the trivial cofibrations, and even the case $I=e$ would give the wrong class.

The free-diagram adjunction makes the construction concrete. For $i\in I$, evaluation at $i$ has a left adjoint
\[
 (i_!A)(k)=\coprod_{I(k,i)}A.
\]
A lifting problem against $i_!a$ is exactly a lifting problem against $a$ at $i$. If $S$ and $T$ generate the two lifting pairs in $\mathcal M$, the small families
\[
 \{i_!s\mid i\in\operatorname{Ob}I,\ s\in S\},
 \qquad
 \{i_!t\mid i\in\operatorname{Ob}I,\ t\in T\}
\]
therefore detect the componentwise trivial fibrations and fibrations. The required smallness is inherited because evaluation preserves the colimits used in the small object argument. Thus a domain small for a sufficiently long cellular construction in $\mathcal M$ remains small after applying $i_!$. The common bound may depend on $I$ and on the generating families; no bound independent of every small category is needed.

To see why the second family produces weak equivalences, evaluate a cell attachment at $k$. It becomes a coproduct of pushouts of maps in $T$, indexed by the sets $I(k,i)$. Trivial cofibrations in $\mathcal M$ are closed under coproducts, pushouts, transfinite composites and retracts. The evaluated relative cell map is consequently a trivial cofibration. The transfer criterion applies. This is the projective model structure of present terminology; the modern combinatorial statement appears in \citet[Theorem~1.28]{Groth2013}. The word ``special'' in this chapter includes small generating families and the stated accessibility of objects, but it does not by itself assert that the underlying category is locally presentable.

The rest of Theorem 2.3 is formal once these factorizations exist, but one step deserves to be made explicit. Suppose $L$ preserves cofibrations and trivial cofibrations and that $f:X\to Y$ is a weak equivalence between cofibrant objects. Factor
\[
 (f,1_Y):X\amalg Y\xrightarrow{k}Z\xrightarrow{p}Y
\]
with $k$ a cofibration and $p$ a trivial fibration. Let $i=k|_X$ and $j=k|_Y$. Both are cofibrations because $X$ and $Y$ are cofibrant. Since $pi=f$ and $pj=1_Y$, two-out-of-three makes $i$ and $j$ trivial cofibrations. Now $L(p)L(j)=1$, while $L(j)$ is weak; hence $L(p)$ is weak, and so is $L(f)=L(p)L(i)$. This proves exactly the homotopy invariance needed to define $u_!^W$ on cofibrant diagrams. It also explains why taking the ordinary Kan extension of an arbitrary representative is insufficient.
\end{dcomment}


\textbf{Théorème 2.3}\footnote[7]{cf. scholie 2.10, qui devrait précéder le théorème 2.3.}. \textit{Soit $(W,C,F)$ un triple de Quillen spécial dans $\M$. Pour toute $I$ dans} Cat\textit{, considérons le triple $(W(I),C\langle I\rangle,F(I))$ dans $\M(I) = \underline{\mathrm{Hom}}(I^{\circ},\M)$, où les notations sont celles de page 10, donc $W(I)$, $F(I)$ et $TF(I)$ se définissent fibre par fibre, et $C\langle I\rangle = (F(I))^* = (v_{I\,!}(C(I_0)))^{\sim}$}\footnote[8]{où $v_I : I_0 \longrightarrow I$ ($I_0$ catégorie discrète ayant mêmes objets que I).}\textit{.}
\begin{itemize}
\item[1°)] \textit{Pour toute $I$ dans} Cat\textit{, $(W(I),C\langle I\rangle,F(I))$ est}
\msp{page 20}
\textit{un triple de Quillen spécial dans $\M(I)$.}
\item[2°)] \textit{Si $u : I \longrightarrow I'$ est une flèche dans} Cat\textit{, d'où}
\[
\begin{tikzcd}
\M(I) \arrow[r, shift left, "u_!"] & \M(I'), \arrow[l, shift left, "u^*"]
\end{tikzcd}
\tag*{(2.3.1)}
\]
\textit{alors}
\[
\left\{
\begin{array}{l}
\left.\begin{array}{rcl}
u^*(W(I')) & \subset & W(I)\\
u^*(F(I')) & \subset & F(I)\\
u^*(TF(I')) & \subset & TF(I)
\end{array}\right\}\textit{tautologie}\\
\begin{array}{rcl}
u_!(C\langle I\rangle) & \subset & C\langle I'\rangle\\
u_!(TC\langle I\rangle) & \subset & TC\langle I'\rangle\ ,
\end{array}
\end{array}\right.
\tag*{(2.3.2)}
\]
\textit{enfin, $u_!$ induit un foncteur}
\[
\M(I)_{\cc} \longrightarrow \M(I')_{\cc} \qquad\quad \text{\footnotemark[9]}
\tag*{(2.3.3)}
\]
\footnotetext[9]{Sous-catégories des objets cofibrants.}
compatible aux localiseurs $W(I)_{\cc} = W(I)|\M(I)_{\cc}$ \textit{et} $W(I')_{\cc} = W(I')|\M(I')_{\cc}$\textit{, donc induit un foncteur}
\[
u_!^W\ :\ \Ho_W(I) \longrightarrow \Ho_W(I')\ .
\tag*{(2.3.4)}
\]
\printedpage{13}
% head: A. Grothendieck : Dérivateurs, Ch. XIX, Édition des Universités de Montpellier II et Paris VII
% foot: Version 11 février 2019
\textit{Ce foncteur est un adjoint à gauche du foncteur $u^*_W$,}
\[
u^*_W\ :\ \Ho_W(I') \longrightarrow \Ho_W(I)\ ,
\tag*{(2.3.5)}
\]
\textit{induit par $u^*$} (3.3.1). % [sic] (3.3.1) -> (2.3.1), the adjunction displayed in (2.3.1)
\end{itemize}
\textsc{Démonstration de 2.3.}

\textit{Prouvons} 1°). On considère comme dans 1.10 le foncteur
\[
v = v_I\ :\ I_0 \longrightarrow I
\]
et le couple de foncteurs adjoints
\[
\begin{tikzcd}
\M(I) \arrow[r, shift right, "v^*"'] & \M(I_0)\ , \arrow[l, shift right, "v_!"']
\end{tikzcd}
\tag*{(2.3.6)}
\]
auquel on peut appliquer 2.2. Il reste
\msp{page 21}
à vérifier seulement
\[
v^*v_!(TC)\ \subset\ TC\ ,
\tag*{(2.3.7)}
\]
or on a, pour $X = (X_i)_{i \in I_0} \in \M(I_0)$,
\[
\begin{aligned}
v_!(X)(i) &\ \simeq\ \textstyle\coprod_{j \in I_0} v_{j\,!}(X_j)(i)\\
&\ \simeq\ \textstyle\coprod_{j \in I_0}\ \ \underbrace{X_j^{\mathrm{Hom}(i,j)}}_{\substack{\text{somme directe}\\ \text{de }\mathrm{Hom}(i,j)\text{ copies}}}\qquad ,
\end{aligned}
\]
et la relation (2.3.7) résulte du fait que $TC$, étant Q-clos à gauche, est stable par petites sommes directes.

\dgrblock{dgrb1172}{A constant diagram whose homotopy colimit is a classifying space}
\begin{dcomment}{examples}{p21:et la relation (2.3.7) résulte}{A constant diagram whose homotopy colimit is a classifying space}
The free objects in the preceding proof have a geometric meaning: they introduce the room needed to carry the relations in a diagram through a replacement. A diagram may have cofibrant values and still fail to be cofibrant as a diagram. Group actions give a complete calculation of the difference.

Let $G$ be a nontrivial group and let $I=BG$ be the category with one object and automorphism group $G$. A contravariant diagram of simplicial sets on $BG$ is a simplicial set with a right $G$-action. Give simplicial sets their usual Kan model structure and give these action diagrams the projective structure. The ordinary colimit is the orbit simplicial set. In particular the colimit of the point with trivial action is a point. We calculate its derived colimit directly.

Let $\mathcal E_G$ be the category whose objects are the elements of $G$ and which has exactly one morphism between any ordered pair of objects. Right translation gives a free action on this category. Its nerve $EG$ has
\[
 (EG)_n=G^{n+1},
\]
with faces deleting an entry and degeneracies repeating an entry; $G$ acts diagonally on the right. The category $\mathcal E_G$ has an initial object, indeed every object is initial, so its nerve is contractible. An explicit contraction is obtained by choosing the identity element and inserting it as an extra first vertex. The augmentation $EG\to *$ is therefore a weak equivalence of the underlying simplicial sets. Since $\mathcal E_G$ is a groupoid, its nerve is Kan: a missing face in a horn is supplied by composition and inversion. Thus the augmentation is also a fibration.

It is also a projective cofibrant replacement. To verify cofibrancy without assuming an equivariant model theorem, filter $EG$ by its skeleta. The action on its nondegenerate simplices is free: an element fixing the tuple $(g_0,\ldots,g_n)$ must fix $g_0$, hence is the identity. After choosing one representative in every orbit of nondegenerate $n$-simplices, the $n$-skeleton is obtained from the preceding skeleton by attaching copies of
\[
 G\times\partial\Delta^n\longrightarrow G\times\Delta^n.
\]
These are the free diagrams on the generating cofibrations of simplicial sets. Starting from the empty diagram proves that $EG$ is projectively cofibrant.

Its orbit simplicial set is the ordinary bar construction $BG$. There is one vertex, an edge for every group element, and a two-simplex for every composable pair. Its fundamental group has generators $[g]$, $g\in G$, with relations
\[
 [1]=1,\qquad [g][h]=[gh],
\]
so the evident map identifies it with $G$. Thus the derived colimit of the constant point is
\[
 \operatorname{hocolim}_{BG^{\mathrm{op}}}*\simeq EG/G\simeq BG,
\]
which cannot be contractible for nontrivial $G$. For $G=C_2$ the first homology group is $\mathbb Z/2$, whereas the first homology of the ordinary orbit point vanishes. The calculation also proves that the trivial-action point was not projectively cofibrant: otherwise the weak equivalence $EG\to *$ between cofibrant diagrams would remain a weak equivalence after taking orbits, contrary to the calculation.

Every individual simplicial set in this example is cofibrant. What the replacement corrects is therefore the action, including its coherent relations, rather than the values of the diagram. This is the concrete content of the angle brackets in $C\langle I\rangle$. The same distinction will reappear with arrows reversed when injectively fibrant diagrams are used to compute right Kan extensions. The general diagram conventions and their Kan adjunctions are recorded in \citet[Definitions~A.2.8.1 and Proposition~A.2.8.7]{Lurie2009HTT}; the calculation above gives the elementary free resolution behind this instance.
\end{dcomment}


\textit{Prouvons} 2°). Les formules (2.3.2) sont essentiellement tautologiques, les deux dernières résultent de la troisième et de la deuxième, en vertu de 1.5 appliqué au couple de foncteurs adjoints $(u_!,u^*)$. La relation $u_!(C\langle I\rangle) \subset C\langle I'\rangle$, et le fait que $u_!$ transforme objet initial en objet initial, impliquent (2.3.3). Que $u_!$ restreint à $\M(I)_{\cc}$ soit compatible aux localiseurs provient de $u_!(TC\langle I\rangle) \subset TC\langle I'\rangle$, et du fait que toute flèche $f$ dans $\M(I)_{\cc}$ qui est dans $W(I)$, se décompose en $f = pi$, avec $i \in TC$ et $p$ inverse à gauche d'une $j \in TC\langle I\rangle$. Ceci vu, joint aux deux premières formules (2.3.2) montre qu'on est sous les conditions
\msp{page 22}
du théorème d'adjonction de Quillen, d'où l'adjonction de $u_!^W$ et de $u^*_W$, q.e.d.

\textbf{Corollaire 2.4.} \textit{Soit $W$ un localiseur dans $\M$, tel qu'il existe un triple de Quillen spécial $(W,C,F)$ incluant $W$. Alors le prédérivateur} $\Ho_W$ \textit{sur} Cat \textit{satisfait aux axiomes Dér 1 (décomposition), Dér 2 (localisation), Dér 3 (objet final), Der 4 bis (existence des $u_!^W$), Dér 5 bis (calcul des fibres des $u_!^W$), Dér 6 bis (exactitude à gauche, i.e. exactitude à gauche des carrés $W$-cocartésiens dans les} $\Ho_W(I)$\textit{), Dér 6 (exactitude à droite des carrés $W$-cartésiens).} % [sic] Der 4 bis -> Dér 4 bis, accent missing as printed

On vient de prouver Dér 4 bis, et Dér 5 bis est pratiquement contenu dans le mode d'obtenir de $u_!$, compte tenu de proposition 1.12, qui implique que pour $X$ dans $\M(I)_{\cc}$, et $i'$ dans $I'$, l'image inverse de $X$ dans $\M(i'\backslash I)$ est dans $\M(i'\backslash I)_{\cc}$ \dots

% ---- from der-ch19_p016-020.tex
\printedpage{14}
% head: A. Grothendieck : Dérivateurs, Ch. XIX, Édition des Universités de Montpellier II et Paris VII
% foot: Version 11 février 2019

\msp{page 23}
L'axiome de décomposition est satisfait pour tout prédérivateur de la forme $\mathbf{D}_{\mathcal{M},W}$. L'axiome d'exactitude à gauche Dér 5 bis est valable, pourvu que les $(\mathcal{M}(I),W(I))$ s'insèrent dans une structure à cofibrations, ce qui est le cas. De même pour l'axiome Dér 2 de localisation, pourvu que de plus $W$ soit fortement saturé, ce qui est également le cas. Enfin, l'axiome de l'objet final dit que si $I$ admet un objet final, alors
\[ \Ho_W(e) \longrightarrow \Ho_W(I) \]
est pleinement fidèle, ou encore que le morphisme d'adjonction
\[ u_!u^*(\xi) \longrightarrow \xi \]
est un isomorphisme. Mais en vertu de 1.12, cas 3°, dans ce cas $u^*$ applique $\mathcal{M}(e)_{\mathrm{c}} = \mathcal{M}_{\mathrm{c}}$ dans $\mathcal{M}(I)_{\mathrm{c}}$, donc le calcul du couple de foncteurs adjoints $u_!^W$, $u_W^*$ peut se faire en utilisant les restrictions de $u_!$, $u^*$ à $\mathcal{M}(I)_{\mathrm{c}}$ et $\mathcal{M}_{\mathrm{c}}$ respectivement. Comme le morphisme d'adjonction $u_!^{\mathrm{c}}u_{\mathrm{c}}^* \longrightarrow \mathrm{id}$ est un isomorphisme, il en est donc de même de $u_!^W u_W^* \longrightarrow \mathrm{id}$, qui s'en déduit en localisant, q.e.d.

\dgrblock{dgrb1173}{The comma formula and the first derivator axioms}
\begin{dcomment}{exposition}{p23:Comme le morphisme d'adjonction}{The comma formula and the first derivator axioms}
Grothendieck wants more than an adjoint on each localized diagram category. He wants its values to be calculable from the part of the diagram lying over a chosen object. This local calculation is what makes the construction a derivator and permits the subsequent base-change arguments. The opposites in his convention must be retained throughout the calculation.

Let $u:I\to J$, $j\in J$, and let $q:j\downarrow u\to I$ forget the arrow $j\to u(i)$. For a contravariant diagram $X$, the ordinary formula is
\[
 (u_!X)(j)=\operatorname{colim}_{(j\downarrow u)^{\mathrm{op}}}q^*X.
\]
To derive this equality, one must know that the same representatives compute both sides. Put $E=j\downarrow u$. The functor $q:E\to I$ is a discrete opfibration: an arrow $i\to i'$ carries $(i,a:j\to u(i))$ to $(i',u(i\to i')a)$, uniquely. For every $i\in I$ and object $A\in\mathcal M$, this unique-lifting property gives
\[
 q^*(i_!A)\cong\coprod_{e\in q^{-1}(i)} e_!A.
\]
Indeed, at $e'\in E$, the indexing sets on the two sides are identified by the bijection
\[
 \coprod_{e\in q^{-1}(i)} E(e',e)
       \xrightarrow{\ \sim\ } I(q(e'),i).
\]
Each arrow on the right has exactly one lift with source $e'$, whose target is in the indicated fibre. Uniqueness also makes the bijection compatible with precomposition, so it defines an isomorphism of diagrams, natural in $A$. The same formula holds for arrows $a$ in place of objects $A$.

Apply it to the two generating families in the preceding block. Restriction sends each generating projective cofibration to a coproduct of generating projective cofibrations in $\mathcal M(E)$, and likewise for trivial cofibrations. It preserves colimits, computed pointwise, and every functor preserves retract diagrams. It consequently preserves the cellular and retract closures of both families, hence all projective cofibrations and trivial cofibrations. It also preserves the initial object, so $q^*QX$ is projectively cofibrant whenever $QX$ is. This proof uses precisely the cocompleteness and generation hypotheses already available here.

Apply the ordinary formula to a projectively cofibrant replacement $QX\to X$. On the left it computes $u_!^WX$; on the right it computes the homotopy colimit of $q^*X$. Hence
\[
 (u_!^WX)_j\cong
 \operatorname{hocolim}_{(j\downarrow u)^{\mathrm{op}}}q^*X.
\]
The isomorphism is the canonical comparison obtained from the comma square, since the ordinary comparison was used before localization. This last point matters: an abstract isomorphism between the two objects would not establish the base-change axiom.

The argument for a terminal object in $I$ is similarly precise. If $t$ is terminal, then $I(i,t)$ is a singleton for each $i$, so $t_!A$ is the constant diagram on $A$. Thus the constant-diagram functor takes cofibrant objects to projectively cofibrant diagrams. Since $I$ is nonempty and connected, the colimit of a constant diagram is $A$, and the counit is an isomorphism already before localization. Its derived counit is consequently an isomorphism. This establishes full faithfulness of constant coefficients in the form claimed by Dér~3.

There is one overstatement in the paragraph being commented on. Decomposition for arbitrary small disjoint unions does not follow for an entirely arbitrary relative category merely from the formula $\mathbf D_{\mathcal M,W}(I)=\mathcal M(I)[W(I)^{-1}]$. Localization need not commute with infinite products: a single morphism in a product localization has a finite zigzag presentation whose length is uniform over all components. The commentary on Chapter~VI, manuscript p.~131, addresses that obstruction, and should be consulted here. In the present model-category situation the conclusion is valid. Componentwise replacements and the description of maps between cofibrant-fibrant objects by homotopy classes supply uniform representatives and uniform homotopies. The same description proves that evaluation detects isomorphisms: a representative with weak-equivalence components is a weak equivalence in the diagram model structure.

Finally, the two exactness axioms concern finite homotopy pushouts and pullbacks. They do not require that arbitrary right Kan extensions have already been constructed. The cofibration and fibration structures provide the required finite calculus. The reference to ``Dér~5 bis'' for left exactness on manuscript page~23 should read ``Dér~6 bis''; the statement on page~22 has the right numbering. The author supplies the omitted right-exactness argument himself on page~24. These are local repairs in the exposition, not failures of Corollary~2.4.
\end{dcomment}


\textbf{2.5.} Pour achever de construire $\Ho_W$ comme un dérivateur sur Cat tout entier, il faut \msp{page 24} donc établir Dér 4, Dér 5 : existence des $u_*^W$, et calcul des fibres de $u_*^W$. Quant à Dér 6, il est acquis automatiquement, du fait que chaque $(\mathcal{M}(I),W(I))$ s'insère dans une structure de fibration $(W(I),F(I))$. (Oubli dans la démonstration de 2.4.) L'idée naturelle est d'utiliser à présent la structure de triple
\[ (W(I),C(I),\ \underbrace{F[I]}_{\deq TC(I)_*}\ )\ , \tag*{(2.5.1)} \]
mais la difficulté est de prouver la factorisation pour un couple
\[ C(I) \longleftrightarrow TF[I]\ , \qquad \underbrace{TC(I)}_{\deq C(I)\cap W(I)} \quad \longleftrightarrow F[I]\ , \tag*{(2.5.2)} \]
faute de pouvoir montrer, dans les conditions générales de 2.3 seulement (triple de Quillen spécial) que $C(I)$ et $TC(I)$ sont également spéciaux. (Cf. question 1.11.) Mais moyennant des conditions convenables sur $\mathcal{M}$, on arrivera à s'en tirer si $C = \mathit{mono}$. L'énoncé pertinent, un peu technique, se trouve dans XV cor. 7.4 (p. 98-99), que je vais rappeler ici.

\dgrblock{dgrb1174}{Why monomorphisms enter the construction}
\begin{dcomment}{context}{p24:faute de pouvoir montrer}{Why monomorphisms enter the construction}
The obstruction to right Kan extensions is now a problem about generating cofibrations. The source already knows how to generate the projective classes by free diagrams. The injective classes are prescribed componentwise, and a small generating family for them is not produced by simply taking the old generators componentwise. Grothendieck returns to monomorphisms because subobjects provide a way to control the size of successive attachments. This is the same geometry of unions and injective objects that underlies the older theory of sheaves and abelian categories.

The six conditions in Theorem~2.6 separate this size control from the exactness needed to glue subobjects. Cocompleteness and accessibility of individual objects supply the small object argument. The cardinal filtration adds a supply of objects bounded in size; it must not be silently absorbed into the word ``special.'' Left exactness of filtered colimits allows intersections to pass through increasing unions. Condition~(e) says that the union of two subobjects can be formed by their pushout over the intersection. Condition~(f) ensures that attaching such a piece preserves monicity even when the attaching map is not itself monic. These hypotheses explain why the subsequent proof can remain inside the category of subobjects of a fixed target.

The direction of the limit in (2.6.1d) requires an editorial correction. The edition prints a projective-limit arrow. The cited antecedent, XV, Theorem~6.1(d), manuscript page~76, has inductive limits, and XVIII, Proposition~4.6.7, manuscript page~157, likewise gives closure of the cardinal stages under inductive limits of bounded size \citep{Grothendieck1990XV,Grothendieck1990XVIII}. This is also what the intended examples require. For the ordinary filtration of sets by cardinality at most $\pi$, a coproduct of at most $\pi$ sets of that size still has size at most $\pi$. By contrast, the product of $\pi$ copies of a two-element set has size $2^\pi>\pi$. We preserve the printed sign and use the intended inductive-limit condition when explaining the theorem.

The examples in the footnotes give a useful test of the exactness assumptions. In an abelian category the pushout of subobjects $A_1,A_2\subset B$ over their intersection is their sum $A_1+A_2$, which embeds in $B$. If $i:A_0\hookrightarrow A_1$ is monic and $f:A_0\to A_2$ is arbitrary, the pushout is
\[
 (A_1\oplus A_2)/\operatorname{im}(i,-f).
\]
The map from $A_2$ into this quotient is monic: an element of the intersection of $0\oplus A_2$ with $\operatorname{im}(i,-f)$ must come from an element killed by $i$, hence from zero. The same proof can be phrased with kernels and images in any abelian category. Thus these conditions extend beyond topoi.

For groups the effective-union condition fails visibly. Inside $C_2\times C_2$, take the two coordinate subgroups. Their intersection is trivial, while their pushout in groups is $C_2*C_2$. The canonical map to $C_2\times C_2$ kills the commutator of the two nonidentity generators. That commutator is a nonempty reduced word in the free product, so the map is not injective. The appeal to the absence of enough injectives in groups is therefore not an arbitrary warning: the precise gluing mechanism used by the theorem has already broken. The cited SGA passage establishes that there are groups which embed in no injective group \citep[Exposé~I, Example~8.12.9]{SGA4I1972}.

The tentative footnote about removing condition~(d) should be read alongside Grothendieck's own refinement in XVIII, 4.6.7. There he replaces the excessive demand on all quotients by factorization into an effective epimorphism followed by a monomorphism, recovering the weaker quotient condition actually used in XV. This is a specified replacement of hypotheses; it does not justify the projective-limit sign printed here, nor does it establish a small generating set for every accessible class of monomorphisms.
\end{dcomment}


\printedpage{15}
% head: A. Grothendieck : Dérivateurs, Ch. XIX, Édition des Universités de Montpellier II et Paris VII
% foot: Version 11 février 2019

\msp{page 25}
\textbf{Théorème 2.6.} \textit{Soit $\mathcal{M}$ une catégorie satisfaisant les conditions suivantes.}
\[ \left\{ \begin{minipage}{0.72\textwidth}\itshape
\begin{itemize}[leftmargin=2.5em]
\item[\upshape a)] $\mathcal{M}$ est “spéciale”, i.e. stable par petites $\varinjlim$ , et ses éléments sont accessibles.
\item[\upshape b)] Les limites inductives filtrantes dans $\mathcal{M}$ sont exactes à gauche (et dans $\mathcal{M}$ les $\varprojlim$ finies sont représentables – résulte de \textup{a} et \textup{c}).
\item[\upshape c)] $\mathcal{M}$ admet une filtration cardinale $(\Filt^\pi)_{\pi\geq\pi_0}$.
\item[\upshape d)] Les $\Filt^\pi$ sont stables par $\varprojlim$ indexées par des catégories $I$ telles que $\operatorname{card}\Fl(I) \leq \pi$, et par quotients ($Z \in \Filt^\pi$ et $Z \longrightarrow Z'$ épimorphisme implique $Z' \in \Filt^\pi$).
\item[\upshape e)] Si $A_1$, $A_2$ sont deux sous-objets d'un objet $B$ de $\mathcal{M}$, et $A_0 = A_1 \cap A_2$, alors $A_1 \amalg_{A_0} A_2 \longrightarrow B$ est un monomorphisme.
\item[\upshape f)] Si on a un carré cocartésien
\begin{center}
\begin{tikzcd}
A_0 \arrow[r, "i"] \arrow[d] & A_1 \arrow[d] \\
A_2 \arrow[r, "i'"] & A_3
\end{tikzcd}
\end{center}
avec $i$ monomorphisme, alors $i'$ est un monomorphisme.
\end{itemize}
\end{minipage}\right. \tag*{(2.6.1)} \]
$^{(10,11,12,13)}$\footnotetext[10]{\textbf{N.B.} a + c signifie aussi que $\mathcal{M}$ est une “catégorie $\mathfrak{U}$-paratopos”, et implique l'existence d'une filtration cardinale satisfaisant d) (sauf l'assertion concernant les quotients $Z'$ de $Z$). Donc il s'agit d'un $\mathfrak{U}$-paratopos satisfaisant de plus b) e) f). (Est-ce alors un topos\,? \textit{Non,} car e), f) sont O.K. dans les \textit{catégories abéliennes\,!})}\footnotetext[11]{\textbf{N.B.} La condition d) est sans doute inutile, du moins si les épimorphismes sont effectifs, car d) sera satisfaite en tous cas pour $\pi = 2^c$, avec $c \geq \pi_0$.}\footnotetext[12]{\textbf{N.B.} La condition e) n'est \textit{pas} satisfaite pour la catégorie (Gr), et le théorème est en défaut pour cette catégorie, laquelle n'admet \textit{pas} “assez d'injectifs”, cf. SGA 4 I 8.12.9. Par contre, e) et f) sont satisfaites par les catégories \textit{abéliennes.}}\footnotetext[13]{Cf. XVIII 4.6.7 (p. 157) pour une formulation plus sympathique de l'ensemble des conditions a) à f) pour $\mathcal{M}$.}. \textit{(On dira alors que $\mathcal{M}$ est} “très spéciale”\textit{). Soit de plus}
\[ W \ \subset \Fl(\mathcal{M}) \tag*{(2.6.2)} \]
\textit{un localiseur dans $\mathcal{M}$, qui satisfait les conditions}
\[ \left\{ \begin{minipage}{0.72\textwidth}\itshape
\begin{itemize}[leftmargin=2.5em]
\item[\upshape a)] $W$ est stable par $\varinjlim$ filtrantes $^{\textup{(14)}}$.
\item[\upshape b)] $W$ est de plus $\mathrm{L}_1$-accessible, i.e. il existe une \textup{petite} partie $W_0 \subset W$ telle que $W$ soit la $\mathrm{L}_1$-enveloppe de $W$, i.e. que $W$ soit la plus petite partie de $\Fl(\mathcal{M})$ contenant $W_0$, et stable par $\varinjlim$ filtrantes. % [sic] la L_1-enveloppe de W -> la L_1-enveloppe de W_0, W_0 being the small generating part just introduced
\end{itemize}
\end{minipage}\right. \tag*{(2.6.3)} \]
\footnotetext[14]{\textbf{N.B.} a) implique déjà que $W = W^\natural$.}

\printedpage{16}
% head: A. Grothendieck : Dérivateurs, Ch. XIX, Édition des Universités de Montpellier II et Paris VII
% foot: Version 11 février 2019

\msp{page 26}
\textit{Posons alors}
\[ C \ = \ \Mono(\mathcal{M}), \qquad TC \ = \ C\cap W\ , \qquad TF \ = \ C_*\ , \qquad F = (TC)_*\ . \tag*{(2.6.4)} \]
\textit{Sous ces conditions :}
\begin{itemize}[leftmargin=2.5em]
\item[1°)] \textit{$C \Qcouple TF$, plus précisément $C \qpair TF$ est un couple de Quillen spécial.} % reading resolved: DER_XIX_corrections.csv
\item[2°)] \textit{Pour que $TC \qpair F$ soit un couple de Quillen, il faut et il suffit que $TC$ soit stable par cochangement de base, et alors c'est un couple de Quillen spécial.} % reading resolved: DER_XIX_corrections.csv
\item[3°)] \textit{Pour que $(W,C,F)$ soit un triple de Quillen (lequel sera nécessairement spécial par \textup{1°, 2°}, et en particulier clos), il faut et il suffit que les conditions suivantes soient satisfaites.}
\[ \left\{ \begin{array}{ll}
\text{a)} & TC\ (\deq W\cap\Mono)\ \text{\itshape est stable par cochangement de base.}\\[1ex]
\text{b)} & TF\ (\deq (\Mono)_*) \ = \ \text{\itshape “applications de type injectif”}\ \subset W.
\end{array}\right. \tag*{(2.6.5)~$^{(15)}$} \]
\footnotetext[15]{\textbf{N.B.} Ces conditions (2.6.5) équivalent à dire que $(W,C)$ est une structure de catégorie à cofibrations sur $\mathcal{M}$.}
\textit{Et la condition} b) \textit{équivaut encore à :}
\[ \begin{minipage}{0.7\textwidth}
\begin{itemize}[leftmargin=2.5em]
\item[b$'$)] \textit{Toute flèche $f$ de $\mathcal{M}$ (ou seulement de $TF$) se factorise en $f = pi$, avec $i$ monomorphisme, $p \in W$.}
\end{itemize}
\end{minipage} \tag*{(2.6.6)} \]
\end{itemize}

\textbf{Corollaire 2.7.} \textit{$\mathcal{M}$ et $W$ étant comme dans} (2.6.1), (2.6.3), \textit{pour que $(W,C,F)$ soit un triple de Quillen} propre à gauche \textit{(i.e. $C \subset \Cof_W$), il faut et il suffit que $W$ satisfasse la condition} (2.6.5. b) \textit{(i.e. $W$ contient les flèches de type injectif), et de plus}
\msp{page 27}

\begin{itemize}[leftmargin=7em]
\item[] $C \subset \Cof_W$, \textit{i.e. la condition :}
\item[a$'$)] \textit{Pour tout diagramme cocartésien}
\[ \begin{tikzcd}
A_0 \arrow[r, "i"] \arrow[d, "f"'] & A_1 \arrow[d, "f'"'] \\
A_2 \arrow[r, "i'"] & A_3
\end{tikzcd} \tag*{(2.7.7)} \]
\textit{avec $i$ monomorphisme, si $f \in W$, alors $f' \in W$. (Stabilité de $W$ par cochangement de base par monomorphismes.)}
\end{itemize}

\textit{Sous ces conditions, ce triple de Quillen est spécial, et a fortiori clos.}

Ce corollaire résulte de 2.6. 3°), puisque $C \subset \Cof_W$ implique $TC \deq C\cap W = C\cap W^{\mathrm{univ}}$ (car $\Cof_W \cap W \ = \ W^{\mathrm{univ}}$), donc que $TC$ est stable par cochangement de base, d'où la condition (2.6.5. a).

\dgrblock{dgrb1175}{The recognition argument and its essential cancellation}
\begin{dcomment}{exposition}{p27:Ce corollaire résulte de 2.6. 3°)}{The recognition argument and its essential cancellation}
Grothendieck's criterion is designed so that the difficult set-theoretic work produces two factorizations, after which the model axioms follow from a short argument. Keeping these stages distinct also prevents the recalled theorem from being confused with a broader assertion about arbitrary accessible classes of monomorphisms. The latter is not what Theorem~2.6 needs.

Set $C=\operatorname{Mono}$, $TC=C\cap W$, $TF=C^{\square}$ and $F=(TC)^{\square}$. Suppose that the two asserted weak factorization systems have been obtained and that $TF\subset W$. Then $TF\subset F$ because $TC\subset C$. We verify the remaining equality $F\cap W=TF$. Given $f\in F\cap W$, use the first factorization system to write
\[
 f=pi,\qquad i\in C,\quad p\in TF.
\]
Since $p,f\in W$, two-out-of-three puts $i$ in $TC$. Since $f$ has the right lifting property against $i$, the usual diagonal makes $f$ a retract of $p$ in the arrow category. Right lifting classes are retract-closed, so $f\in TF$. The reverse inclusion has already been observed. Hence the two weak factorization systems have exactly the intersections required for a closed Quillen triple.

The alternative condition (2.6.6) is the same retract argument in a different position. If $TF\subset W$, factorization with respect to $(C,TF)$ supplies a monomorphism followed by a weak equivalence. Conversely suppose each $f\in TF$ factors as $f=pi$ with $i$ monic and $p\in W$. The right lifting property of $f$ against $i$ makes $f$ a retract of $p$, so $f\in W$, provided $W$ is retract-closed. That last property is already forced by (2.6.3a): if $f$ is a retract of $g\in W$, split the associated idempotent $e:g\to g$ by the chain
\[
 g\xrightarrow e g\xrightarrow e g\xrightarrow e\cdots .
\]
Its colimit in the arrow category is $f$. Filtered-colimit closure therefore gives $f\in W$.

The historical proof of the generating-set assertion should be located accurately. XV, Theorem~6.1 carries the author's warning that its proof is invalid. The surviving argument is the separate Lemma~7.3 and its application in~7.1--7.4. Starting with $A\hookrightarrow B$ in $TC$, one chooses a bounded subobject $B_0\subset B$ for which $A\cap B_0\hookrightarrow B_0$ lies in $TC$ and adds it to $A$. The density argument permits $B_0$ to contain a prescribed bounded piece of $B$; when $A\ne B$, choose that piece not contained in $A$. Thus the union strictly enlarges $A$. Effective unions keep $A\cup B_0\to B$ monic; pushout stability makes $A\to A\cup B_0$ weak. Since $A\to B$ is weak, cancellation makes $A\cup B_0\to B$ weak. Repeat this construction, taking unions at limit stages. Filtered-colimit closure keeps the resulting map to $B$ in $TC$ at such stages. Since $B$ has only a set of subobjects, strict enlargement cannot continue indefinitely: the process reaches $B$. The cancellation step is the extra ingredient missing for an arbitrary accessible mono class \citep[§7, especially Lemma~7.3 and Corollary~7.4]{Grothendieck1990XV}.

A later independent formulation is Smith's recognition theorem, proved in \citet[Theorem~1.7]{Beke2000}: with a locally presentable underlying category and a chosen set generating the cofibrations, accessibility supplies the solution-set condition for weak equivalences, and pushout/transfinite closure of the trivial cofibrations completes the construction. Beke's Proposition~1.12 identifies hypotheses ensuring that all monomorphisms admit such generators. This provides a modern route to the intended very-special case without appealing to the withdrawn general claim.

Corollary~2.7 records an additional condition, left properness. Its pushout square tests whether a weak equivalence remains weak after attachment along a monomorphism. When the attached map itself is both a monomorphism and a weak equivalence, the two-out-of-three argument in the source identifies it with a universally weak map, whose every cobase change is weak. Hence left properness supplies the pushout stability of $TC$ required by Theorem~2.6. The recognition theorem needs that stability; the stronger properness property is singled out for applications rather than inserted into the definition of a model structure.
\end{dcomment}


\printedpage{17}
% head: A. Grothendieck : Dérivateurs, Ch. XIX, Édition des Universités de Montpellier II et Paris VII
% foot: Version 11 février 2019

\textbf{2.8.} On dit qu'une \textit{structure de Quillen} $(W,C,F)$ sur $\mathcal{M}$ est \textit{très spéciale,} si elle est obtenue comme au théorème 2.1., i.e. si on a les conditions :
\[ \left\{ \begin{minipage}{0.72\textwidth}
\begin{itemize}[leftmargin=2.5em]
\item[a)] $\mathcal{M}$ est très spéciale, i.e. satisfait aux conditions (2.6.1).
\item[b)] $W$ est stable par petites $\varinjlim$ filtrantes, et est accessible (conditions (2.6.3)).
\item[c)] $C = \Mono$.
\end{itemize}
\end{minipage}\right. \tag*{(2.8.1)} \]
Donc pour qu'un localiseur $W$ dans $\mathcal{M}$ corresponde à une structure de Quillen très spéciale \msp{page 28} dans $\mathcal{M}$ (laquelle sera forcément unique), il faut et il suffit que $\mathcal{M}$ et $W$ satisfassent les conditions (2.8.1. a)\,b)), et que de plus on ait :
\[ \left\{ \begin{minipage}{0.65\textwidth}
\begin{itemize}[leftmargin=2.5em]
\item[a)] $TC \deq \Mono\cap W$ est stable par cochangement de base.
\item[b)] $TF \deq (\Mono)_* \subset W$, i.e. $W$ contient les flèches de type injectif dans $\mathcal{M}$.
\end{itemize}
\end{minipage}\right. \tag*{(2.8.2)} \]
\footnote[16]{On dit aussi que le localiseur $W$ est très spécial, s'il définit une structure de Quillen très spéciale.}.

\dgrblock{dgrb1176}{Injective structures and the role of accessibility}
\begin{dcomment}{later}{p28:laquelle sera forcément unique}{Injective structures and the role of accessibility}
Grothendieck has reached a class of coefficient categories for which the second half of the diagram theory can be constructed. Theorem~2.9 uses componentwise cofibrations, so restriction along every functor preserves both cofibrations and trivial cofibrations. The right Kan adjoint is then a right Quillen functor. This is the injective model structure; the square brackets in $F[I]$ indicate that its fibrations have been determined by a lifting property on whole diagrams.

There are two sorts of permanence hidden in the phrase ``par AQT.'' The exactness conditions are evaluated at objects of $I$: limits, colimits, intersections, pushouts and the maps from unions to ambient objects are computed there. Monicity is also componentwise. For the converse direction of that assertion, apply the defining monomorphism test to free diagrams $i_!A$: monicity of $X\to Y$ in the diagram category says that $\mathcal M(A,X_i)\to\mathcal M(A,Y_i)$ is injective for every $A$, hence $X_i\to Y_i$ is monic. The existence of the free adjoint is what makes this argument valid in the stated generality.

The size statements require a common bound. One chooses a cardinal larger than the arrow set of $I$ and enlarges the bounds used for the coefficients and for $W$. A small diagram then involves only a bounded family of objects, structure maps and equations. A sufficiently large accessible presentation can carry all these data at one stage. Equalities of composites are imposed at a later common stage, and the filtered system of such choices recovers the original diagram. Thus the bounds are allowed to grow with the shape. The source's cardinal-filtration hypothesis is doing real work here; accessibility of each individual object alone would not justify this conclusion. A precise antecedent for the diagram-filtration step is \citet[Exposé~I, Corollary~9.24]{SGA4I1972}, with the indexing category held fixed and the cardinal enlarged as prescribed there.

In the later combinatorial formulation, \citet[Lemma~A.2.8.3]{Lurie2009HTT} proves that the componentwise class associated to a set-generated weakly saturated class is again set-generated in a small diagram category. Proposition~A.2.8.2 then gives both projective and injective model structures. This settles the obstacle of page~24 when local presentability is imposed; it should not be quoted as a theorem about every historical ``special'' category with that hypothesis omitted. For presheaves with an accessible localizer and monomorphism cofibrations, the corresponding statement is \citet[Corollary~1.4.22]{Cisinski2006}. These references supply exact later versions of the passage from coefficients to diagrams.

The derived right adjunction itself can be checked without new size arguments. Since $u^*$ preserves the componentwise trivial cofibrations, adjunction shows that $u_*$ preserves injective fibrations. Applying the same reasoning to all componentwise cofibrations gives preservation of trivial fibrations. It preserves the terminal object, hence fibrant objects. The dual of the Brown argument above shows that it preserves weak equivalences between those objects. Localizing its restriction to fibrant objects defines $u_*^W$, and the derived unit and counit give
\[
 u_W^*\dashv u_*^W.
\]
The replacement needed for the pointwise formula can also be specified. Let $q:E=(u\downarrow j)\to I$ forget the arrow $u(i)\to j$. This projection is a discrete fibration: an arrow $i'\to i$ uniquely lifts into an object $(i,u(i)\to j)$ by precomposition. Consequently its left Kan extension on contravariant diagrams is
\[
 (q_!Y)(i)\cong\coprod_{e\in q^{-1}(i)}Y(e).
\]
Coproducts preserve cofibrations and trivial cofibrations, so $q_!$ preserves both componentwise left classes. The adjunction $q_!\dashv q^*$ therefore makes $q^*$ a right Quillen functor for the injective structures. In particular, if $X\to RX$ is an injective fibrant replacement, then $q^*RX$ is injectively fibrant and $q^*X\to q^*RX$ is a weak equivalence. Applying the ordinary right Kan formula to $RX$ now gives
\[
 (u_*^WX)_j\cong
 \operatorname{holim}_{(u\downarrow j)^{\mathrm{op}}}q^*X.
\]
Here the comma category has objects $u(i)\to j$, opposite to the comma category used for the left extension. This accounts for the two different fibre formulas in the derivator axioms.

The two model structures have the same weak equivalences and therefore produce the same localized diagram category. Their purpose is computational: one makes the left adjoints homotopical on cofibrant objects, the other makes the right adjoints homotopical on fibrant objects. The special monomorphism hypotheses guarantee these choices in the source's proof. They are not an invariant property that every presentation of the resulting derivator must share.
\end{dcomment}


On trouve maintenant le

\textbf{Théorème 2.9.} \textit{Soit} $(W, C = \Mono, F = (TC)_* = (W\cap\Mono)_*)$ \textit{un triple de Quillen très spécial dans une catégorie $\mathcal{M}$} (2.8). \textit{Alors :}
\begin{itemize}[leftmargin=2.5em]
\item[1°)] \textit{Pour toute $I$ dans} Cat, $(W(I),C(I),F[I])$ \textit{est un triple de Quillen très spécial dans $\mathcal{M}(I)$, et on a en particulier}
\[ C(I) \ = \ \Mono(\mathcal{M}(I))\ , \tag*{(2.9.1)} \]
\textit{et de plus}
\[ F[I]\cap W(I) \ = \ (TF)[I]\ . \tag*{(2.9.2)} \]
\item[2°)] \textit{Pour toute flèche $u : I \longrightarrow I'$ dans} Cat, \textit{d'où le couple de foncteurs adjoints}
\[ \begin{tikzcd}
\mathcal{M}(I) \arrow[r, shift right, "u_*"'] & \mathcal{M}(I') \arrow[l, shift right, "u^*"']
\end{tikzcd}\ , \tag*{(2.9.3)} \]
\textit{on a}
\[ \left\{ \begin{array}{rcl}
u^*(W(I')) & \subset & W(I)\\
u^*(C(I')) & \subset & C(I)\\
u^*(TC(I')) & \subset & TC(I)\\
u_*(F[I]) & \subset & F[I']\\
u_*(TF[I]) & \subset & TF[I']\ ,
\end{array}\right. \tag*{(2.9.4)} \]
\textit{enfin $u_*$ induit un foncteur}
\printedpage{18}
% head: A. Grothendieck : Dérivateurs, Ch. XIX, Édition des Universités de Montpellier II et Paris VII
% foot: Version 11 février 2019
\msp{page 29}
\[ u_*^{\mathrm{f}} \ :\ \mathcal{M}(I)_{\mathrm{f}} \longrightarrow \mathcal{M}(I')_{\mathrm{f}} \qquad\quad ^{(17)} \tag*{(2.9.5)} \]
\footnotetext[17]{sous-catégories pleines des objets fibrants.}
compatible aux localiseurs $W(I)_{\mathrm{f}}$ \textit{et} $W(I')_{\mathrm{f}}$ \textit{induits par $W(I)$, $W(I')$ respectivement, donc induit un foncteur}
\[ u_*^W \ :\ \Ho_W(I) \longrightarrow \mathit{Ho}_W(I')\ . \tag*{(2.9.6)} \] % [sic] italic Ho_W(I') -> upright Ho_W(I'), inconsistent with the upright Ho_W elsewhere
\textit{Ce foncteur est un} adjoint à droite \textit{du foncteur}
\[ u_W^* \ :\ \mathit{Ho}_W(I') \longrightarrow \Ho_W(I) \tag*{(2.9.7)} \] % [sic] italic Ho_W(I') -> upright Ho_W(I'), inconsistent with the upright Ho_W elsewhere
\textit{induit par $u^*$} (2.9.3).
\end{itemize}
\textsc{Démonstration de 2.9.} Pour prouver 1°), il faut vérifier d'abord que les conditions préliminaires (2.8.1. a)\,b)) sur $\mathcal{M}$, $W$ sont stables par passage de $\mathcal{M}$, $W$ à $\mathcal{M}(I)$, $W(I)$, ce qui se fait par AQT \texttt{[âne qui trotte]}. D'ailleurs on vérifie tout de suite que
\[ C(I) \ = \ \Mono(\mathcal{M}(I))\ , \]
et le théorème 2.6 (1°, 2°) implique alors que
\[ \begin{array}{ccccc}
C(I) & \qpair & C(I)_* & \deq & TF[I]\\
\underbrace{TC(I)}_{\deq C(I)\cap W(I)} & \qpair & (TC(I))_* & \deq & F[I]
\end{array} \] % reading resolved: DER_XIX_corrections.csv
sont des couples de Quillen spéciaux. (En notant pour le deuxième, que $TC(I)$ est bien \msp{page 30} stable par cochangement de base, $TC$ l'étant.) Pour vérifier que l'on a un triple de Quillen très spécial, il reste en vertu de 2.8 à vérifier que
\[ TF[I] \ \subset\ W(I)\ , \]
ce qui s'écrit
\[ v^*(TF[I]) \ \subset\ W(I_0)\ , \]
où
\[ v \ :\ I_0 \longrightarrow I \]
est défini comme dans (1.12) (p. 10). En fait, on aura même
\[ v^*(TF[I]) \ = \ v^*v_*(TF(I_0)) \ \subset\ \underbrace{TF(I_0)}_{= F(I_0)\cap W(I_0)}\quad , \]
et la démonstration est duale de celle de
\[ v^*v_!(TC(I_0)) \ \subset\ TC(I_0) \]
(cf. (2.3.7), p. 21), elle ne dépend pas des questions de factorisation dans $\mathcal{M}(I)$, et provient ici de la stabilité de $TF$ par petits produits. Cela prouve 1°.

\dgrblock{dgrb1177}{The cofree image is smaller than the lifting class}
\begin{dcomment}{development}{p30:elle ne dépend pas des questions de factorisation}{The cofree image is smaller than the lifting class}
Grothendieck needs the inclusion $TF[I]\subset W(I)$ to finish the injective model structure. The intended conclusion is correct, but the equality involving cofree diagrams on this page is stronger than the adjunction permits. We spell out both the direct repair and a finite example distinguishing the two classes.

First there is a short proof of the necessary inclusion. For $v:I_0\to I$, the free-diagram functor $v_!$ takes componentwise monomorphisms to monomorphisms. Indeed, at $i$ it is a coproduct of the given component maps, and these coproducts are monic under the hypotheses in use. One can also see the coproduct assertion by expressing a coproduct of monomorphisms as a transfinite succession of pushouts, using closure of monomorphisms under filtered colimits at limit stages. If $p\in TF[I]=\operatorname{Mono}(\mathcal M(I))^{\square}$, adjunction now gives
\[
 v^*p\in\operatorname{Mono}(\mathcal M(I_0))^{\square}=TF(I_0).
\]
Since $TF\subset W$, this proves the required inclusion. No factorization of $p$ and no assertion that it is cofree have been used.

\textbf{Finite separation of the two image classes.}\ours
In $\mathrm{Set}$, with weak equivalences all maps and cofibrations the injections, take $I=(0\to1)$. Then
\[
 v^*(TF[I])\ne v^*v_*(TF(I_0)).
\]
Here the model structure is generated by the usual injection--surjection weak factorization system; it satisfies the intended, inductive-limit version of the very-special hypotheses.

\emph{Proof.} A contravariant $I$-diagram is a map $X_1\to X_0$. Let $X_0$ have two elements and $X_1$ three, and choose a surjection $r:X_1\twoheadrightarrow X_0$. Let $p:X\to1$ be the map to the terminal diagram. We show directly that it has the right lifting property against every componentwise injection $A\hookrightarrow B$. Regard $A_i$ as a subset of $B_i$. Given a map $A\to X$, first extend $A_0\to X_0$ to $h_0:B_0\to X_0$ by sending all new elements to one fixed element of $X_0$. Choose a section $s:X_0\to X_1$ of $r$. On $B_1\setminus A_1$, define
\[
 h_1(b)=s\bigl(h_0(r_B(b))\bigr),
\]
and on $A_1$ retain the specified map. The two definitions have disjoint domains and agree with the required compatibility on $A_1$. Thus $rh_1=h_0r_B$, giving a lift of diagrams. Hence $p\in TF[I]$.

For a family $(E_0,E_1)$, the cofree diagram is
\[
 (v_*E)_0=E_0,\qquad (v_*E)_1=E_0\times E_1,
\]
with the projection as its structure map. Consequently the restriction of a cofree arrow has component source sizes $|E_0|$ and $|E_0||E_1|$. If it were isomorphic, after restriction, to $p$, these sizes would be $2$ and $3$. The equation $2|E_1|=3$ has no solution. Therefore $v^*p$ is not in the cofree image class, even up to isomorphism. This proves the assertion.\hfill$\square$

The product argument in the source does establish
\[
 v^*v_*(TF(I_0))\subset TF(I_0),
\]
because right lifting classes are closed under products. It does not replace the direct lifting proof for the larger class $TF[I]$. In a closure formulation one must retain the appropriate right-orthogonal saturation; deleting it creates exactly the mistaken equality.

We checked the diagram-model discussions of \citet[§A.2.8]{Lurie2009HTT} and \citet[Corollary~1.4.22]{Cisinski2006}; both use lifting to establish the relevant preservation statements. We did not find this particular finite illustration there. The example is recorded as a local development of the printed argument, with no claim of priority for the general distinction between cofree objects and a right lifting class.
\end{dcomment}


La démonstration de 2° est maintenant la duale de celle de (2.3. 2°). Comme cela nous fait répéter deux fois la même démonstration, cela montre qu'on n'a pas “monté le sorite” tout à fait avec assez de soin. Donc à titre de rattrapage, j'ajoute un

% ---- from der-ch19_p021-025.tex
\printedpage{19}
% head: A. Grothendieck : Dérivateurs, Ch. XIX, Édition des Universités de Montpellier II et Paris VII
% foot: Version 11 février 2019

\msp{page 31}
\textbf{Scholie 2.10.} \textit{Soit $(W,C,F)$ un triple de Quillen clos dans $\M$. Pour $I$ dans \textup{Cat}, pour que $(W(I),C\langle I\rangle,F(I))$ soit un triple de Quillen clos pour $\M(I)$, il faut et il suffit que les couples $C\langle I\rangle\longleftrightarrow TF(I)$ et $TC\langle I\rangle\longleftrightarrow F(I)$ soient factorisants. Supposons que cette condition soit satisfaite pour $I$ et pour $I'$, et soit $u : I\longrightarrow I'$ une flèche de $I$ dans $I'$. Alors les assertions de \textup{2.3. 2°} sont valables, en particulier on trouve un foncteur $u^{W}_{!} : \Ho_W(I)\longrightarrow\Ho_W(I')$, adjoint à gauche de $u^{*}_{W}$\footnote[18]{Dire aussi que le calcul des fibres de $u^{W}_{!}(X)$, prévu dans Dér 5 bis, est OK.}. Ces résultats se dualisent aussitôt (par $\M\longmapsto\M^{\circ}$, $I\longmapsto I^{\circ}$) pour les triples $(W(I),C(I),F[I])$, et pour la construction de $u^{W}_{*}$, adjoint à droite de $u^{*}_{W}$. En particulier, pour que le triple précédent soit un triple de Quillen (forcément clos), il faut et il suffit que les deux couples $C(I)\longleftrightarrow TF[I]$ et $TC(I)\longleftrightarrow F[I]$ soient factorisants ; et quand cette condition est satisfaite pour $I$ et pour $I'$ (pour une flèche $u : I\longrightarrow I'$ dans \textup{Cat} donnée), alors $u^{W}_{*}$ existe et se construit dualement de la construction de $u^{W}_{!}$ explicitée dans \textup{2.3. 2°}.}

\msp{page 32}
% a vertical bar is printed in the left margin alongside Corollaire 2.11
\textbf{Corollaire 2.11.} \textit{Le prédérivateur $\mathbf{D}_{\M,W}$ défini pour un localiseur \textup{très spécial} $W$ dans $\M$ est un dérivateur en plein sens du terme, i.e. il satisfait à tous les axiomes Dér 1 à Dér 6 (y inclus Dér 4 bis et Dér 5 bis et Dér 6 bis) des dérivateurs.}

\textbf{N.B.} À ma propre surprise, je n'ai pas été capable de trouver un cas d'un localiseur $W$ qui définisse un dérivateur sur Cat tout entier (avec existence de $u^{W}_{!}$, $u^{W}_{*}$), \textit{que} pour $W$ un localiseur très spécial, donc s'insérant dans un triple de Quillen très spécial -- ou quand on peut s'y ramener de la façon standard. (Par exemple dans le cas où $\M=\Cat$, $W$ un localiseur fondamental satisfaisant Loc 6 a) et Loc 8, 8 bis)\footnote[19]{axiome des carrés cocartésiens (version directe), axiome des limites inductives, accessibilité.} -- les monomorphismes dans $W$ se refusent formellement à être stable par cochangement de base !) J'ignore totalement si le résultat mirobolant ``prouvé'' par \textsc{Anderson} est vrai ou faux. Ça a l'air presque trop beau pour être vrai \ldots{}D'ailleurs, je ne connais pas de structure de Quillen spéciale qui ne soit très spéciale -- désolé ! % [sic] stable -> stables, agreement with « les monomorphismes »

\dgrblock{dgrb1178}{The existence question after Chapter XIX}
\begin{dcomment}{status}{p32:J'ignore totalement si le résultat mirobolant}{The existence question after Chapter XIX}
The final note measures the distance between the construction just completed and the general statement Grothendieck hoped for. Theorem~2.11 provides both Kan extensions by imposing enough smallness and exactness to construct two model structures on every diagram category. The subsequent uncertainty concerns how much of this apparatus is necessary. That existence question now has a definitive answer for ordinary complete and cocomplete Quillen model categories.

\citet[Theorem~6.11]{Cisinski2003} proves the existence and pointwise calculation of homotopy right Kan extensions for a Quillen model category admitting all small limits; applying the dual result supplies the left extensions when all small colimits exist. This gives the associated derivator on all small categories, without requiring the original category to be combinatorial or the diagram categories to possess the two projective and injective model structures used here. The same general statement is recorded in \citet[Theorem~1.32]{Groth2013}. Groth's Proposition~1.30 gives a shorter proof for combinatorial models and establishes the usual strength axiom in that setting. These later results settle the mathematical question; they do not certify the details of the historical proof attributed to Anderson.

There is already a very small distinction between ``special'' and ``very special.'' On $\mathrm{Set}$ take
\[
 W=\operatorname{Iso},\qquad C=F=\operatorname{Fl}(\mathrm{Set}).
\]
The two weak factorization systems are $(\operatorname{Fl},\operatorname{Iso})$ and $(\operatorname{Iso},\operatorname{Fl})$. The first is generated by $\varnothing\to1$ and $2\to1$: right lifting against the first means surjectivity, and right lifting against the second means injectivity. The empty generating family gives the second system. All sets are accessible at a suitable cardinal and all small colimits exist, so this is special in the chapter's sense. It is not very special because $2\to1$ is a cofibration but not a monomorphism. Its associated derivator is the represented diagram theory of sets. This elementary example separates a chosen model presentation from the special construction with all monomorphisms as cofibrations.

The warning about $\mathrm{Cat}$ points to a different, genuine obstruction. Here is a complete finite construction for the ordinary weak-homotopy localizer. Let $P$ be the poset of nonempty proper faces of a tetrahedron. Its nerve is the barycentric subdivision of the boundary sphere, so $|NP|$ is homeomorphic to $S^2$. For each nonidentity arrow $a$ of $P$, take a copy of the walking arrow $[1]$ and its inclusion in the groupoid $J$ with two objects and one isomorphism between them. Form the disjoint unions
\[
 A=\coprod_a[1],\qquad B=\coprod_a J,
\]
and map $A\to P$ by the chosen arrows. The map $A\to B$ is a monomorphism of categories: it is injective on objects and faithful. It is also a weak homotopy equivalence, since it is a disjoint union of maps between contractible nerves.

The pushout $P\amalg_A B$ is the ordinary groupoid completion of $P$, by the universal property of freely adjoining inverses to all arrows. The fundamental groupoid of the nerve of a category is its groupoid completion: its generators are the one-simplices, and its two-simplex relations are exactly composition. Since $S^2$ is connected and simply connected, this groupoid is equivalent to the terminal groupoid. Thus the pushout map
\[
 P\longrightarrow P\amalg_A B
\]
induces $S^2\to *$ on homotopy types and is not a weak equivalence, as its effect on $H_2(-;\mathbb Z)$ already shows. All categories used to specify the pushout are finite. The claim on page~32 that weak monomorphisms in $\mathrm{Cat}$ need not survive pushout therefore has a concrete content; one cannot obtain the usual homotopy theory of categories by simply declaring every monomorphism to be a cofibration.

Scholie~2.10 isolates what the direct proof actually requires: the two lifting pairs must factor at the diagram shapes being used. For its pointwise assertion one must also justify the corresponding comma-shape calculation, as done above when all the projective structures are available. Cisinski's theorem supplies the general alternative when that direct model-structure construction is unavailable. The categorical limits and colimits remain essential hypotheses; this is the distinction the examples in the following section must continue to respect.
\end{dcomment}


\msp{page 33}

\section*{3\quad Cas des catégories à cofibrations, des catégories à fibrations}

Je rappelle les axiomes d'une \textit{structure à cofibrations} $(W,C)$ sur une catégorie $\M$.

\begin{itemize}[leftmargin=4em,labelwidth=3.5em,align=left]
\item[Cof$_0$ :] $W$ est un localiseur.
\item[Cof$_1$ :] $C$ contient les isomorphismes, est stable par cochangement de base, par composition.
\item[Cof$_2$ :] $C\cap W$ est stable par cochangement de base (et aussi par composition, et contient les isomorphismes, en vertu de Cof$_0$, Cof$_1$).
\item[Cof$_3$ :] Toute flèche $f$ dans $\M$ se factorise en $f=pi$, avec $i\in C$, $p\in W$.
\end{itemize}

(Axiomes de K. S. \textsc{Brown}, légèrement généralisés par \textsc{Anderson}). Il y a un ``axiome
\printedpage{20}
% head: A. Grothendieck : Dérivateurs, Ch. XIX, Édition des Universités de Montpellier II et Paris VII
% foot: Version 11 février 2019
de continuité'' ou \textit{``axiome transfini''} facultatif, qui est le suivant :

\begin{itemize}[leftmargin=4em,labelwidth=3.5em,align=left]
\item[Cof$_4$ :] $\M$ stable par petites $\varinjlim$ filtrantes, et $C$ et $TC = C\cap W$ sont stables par composition transfinie.
\end{itemize}

Ainsi, Cof$_1$, Cof$_2$, Cof$_4$ peuvent se résumer en disant que $C$ et $TC$ sont Q-saturés transfiniment à gauche.

La structure de cofibration est dite \textit{spéciale} si elle satisfait Cof$_4$, si $\M$ est spéciale, et si $C$ admet une \textit{petite} partie $C_0$ qui soit Q-génératrice à gauche, i.e. telle que
\[
C \ \subset\ C_0\Qs\,, \qquad \text{i.e.}\quad \tilde{C}=C_0\Qs. \tag*{(3.1)}
\]

\msp{page 34}
Alors, par le théorème de factorisation 1.1, on a
\[
C \ \fact\ TF \ \deq\ C_* \ ,
\]
et l'axiome de factorisation Cof$_3$ implique (cf. 1.4) $TF\subset W^{\natural}$. Plus généralement, on a :

\textbf{Proposition 3.1.} \textit{Supposons que $\M$ soit spéciale (stable par petites $\varinjlim$ , et éléments accessibles), que $(W,C)$ satisfasse à Cof$_0$, Cof$_1$, et que $C$ soit stable par composition transfinie (ce qui, joint à Cof$_1$, s'énonce en disant que $C$ est Q-saturé transfini à gauche), enfin que $C$ admette une petite partie $C_0$ Q-génératrice à gauche, de sorte que $C\fact TF\deq C_*$. Alors pour que l'axiome de factorisation Cof$_3$ soit satisfait, il suffit qu'on ait}
\[
TF \ \subset\ W \ , \tag*{(3.1.1)}
\]
\textit{et il faut qu'on ait $TF\subset W^{\natural}$, donc si $W=W^{\natural}$, l'axiome Cof$_3$ équivaut à} (3.1.1).

\dgrblock{dgrb1179}{What one factorization can accomplish}
\begin{dcomment}{exposition}{p34:l'axiome Cof$_3$ équivaut}{What one factorization can accomplish}
Grothendieck now removes the fibrations from the starting data. His purpose is to isolate what is actually needed to construct homotopy colimits: the ability to replace a map by a cofibration, together with the invariance of gluing along an acyclic cofibration. The resulting structure is deliberately asymmetric. It controls homotopy cofibres without committing to a second factorization, and therefore without assuming that every lifting class already belongs to a model structure.

Write $TC=C\cap W$ and, in modern notation, write $S^\perp$ for the maps having the right lifting property against $S$. The symbol $TF=C_* =C^\perp$ initially names a lifting class; its membership in $W$ is a separate assertion. Proposition 3.1 keeps this distinction correctly. If the small object argument gives a factorization $f=pi$ with $i\in C$ and $p\in C^\perp$, then $C^\perp\subset W$ immediately supplies Cof$_3$. Conversely, apply Cof$_3$ to a map $f\in C^\perp$, obtaining $f=pi$, $i\in C$, $p\in W$. The square with left side $i$, right side $f$, and identity on the common source has a diagonal $s$ satisfying $si=1$ and $fs=p$. In the arrow category, $(i,1)$ and $(s,1)$ exhibit $f$ as a retract of $p$. This proves exactly $C^\perp\subset W^{\natural}$; replacing the latter class by $W$ requires retract closure. No additional lifting axiom has entered the argument.

The two closure requirements on $C$ and $TC$ have different tasks. Closure under cobase change gives the elementary gluing operation. Closure under transfinite composition permits a diagram to be assembled through its successive stages. Closure under retracts then turns a factorization into a characterization by lifting. In the terminology of these pages, Q-closed means equality with ${}^\perp(S^\perp)$, whereas Q-saturated describes specified operations on arrows. Theorem 1.1 identifies them under a small generation hypothesis; section 3 ends by asking how far that hypothesis can be removed.

The historical starting point is Brown's theory of categories of fibrant objects, used here in its dual form \citep{Brown1973}. Theorem-Scholie~2 of Chapter~VI, manuscript p.~108, already places this asymmetry in the development of the manuscript. What is new here is its systematic application to categories of diagrams. Grothendieck's assertion that a special cofibration structure has functorial factorizations uses the functorial version of the small object construction from section 1. It is a conclusion of that construction, not part of the definition of an arbitrary cofibration category.
\end{dcomment}


La structure de cofibration est dite \textit{semi-close} si $C$ est Q-clos à gauche, i.e. si on a $C=C\Qs$ (ce qui implique que $C$ est stable par composition transfinie), \textit{close} si de plus $TC=(TC)\Qs$ (ce qui implique l'axiome transfini Cof$_4$).

\msp{page 35}
Je ne reviens pas ici sur la théorie de K. S. \textsc{Brown} des catégories à cofibrations. C'est à ma connaissance l'ensemble d'axiomes le plus économique connu, pour assurer d'une part une bonne prise sur $\Ho_W$ (par une notion correspondante d'homotopie d'applications, et un calcul de fractions dans $\Pi\M_{\mathrm{c}}$), et d'autre part l'existence de la suite \textit{exacte} de cofibrations, donc une théorie de cofibres $W$-homotopiques. (\textbf{N.B.} Quand la structure est \textit{propre,} i.e. $C\subset\Cof_W$, alors l'existence de cette structure revient à dire simplement que ``$W$ admet assez de $W$-cofibrations'', i.e. que toute $f\in\Fl(\M)$ se factorise en $pi$, avec $i\in\Cof_W$, $p\in W$ ; et on peut donc prendre $C=\Cof_W$. Dans le cas où $W$ est stable par petites $\varinjlim$ filtrantes, $C=\Cof_W$ et $TC=W^{\mathrm{univ}}$ satisfont de plus l'axiome transfini Cof$_4$.) Enfin, K. S. \textsc{Brown} montre que toute flèche $i : X\longrightarrow Y$ dans $C$, \textit{avec $X$, $Y$ cofibrants} (i.e. $\emptyset_{\M}\longrightarrow X$, $\emptyset_{\M}\longrightarrow Y$ dans $C$) est dans $\Cof_W$ (``Lemme de Brown'').

\dgrblock{dgrb1180}{Brown excision and the scope of properness}
\begin{dcomment}{examples}{p35:avec $X$, $Y$ cofibrants}{Brown excision and the scope of properness}
Grothendieck wants gluing to depend only on the weak equivalences, which explains the intrinsic class $\Cof_W$. One qualification is necessary in the final sentence on this page. Brown excision gives homotopy invariance of a pushout along a cofibration when the weak-equivalence comparison is between cofibrant objects. It does not say that a cofibration with cofibrant endpoints belongs to the ambient, universally defined class $\Cof_W$.

This is a distinction in the manuscript's own terminology. Chapter XIII, ms.~p.~8, defines $\Cof_W$ by two consecutive pushout squares: after any initial cobase change, pushing out a weak equivalence must again give a weak equivalence \citep[Ch.~XIII, ms.~p.~8]{Grothendieck1990XIII}. Thus it requires invariance on the entire undercategory. The safe form of Brown's assertion is obtained by restricting the coefficient category and the comparison maps to cofibrant objects. Alternatively, one may assume $C\subset\Cof_W$, the properness hypothesis already distinguished in the preceding parenthesis.

A concrete obstruction makes the difference visible. In simplicial associative unital $\mathbb Z$-algebras, let $A=\mathbb Z/p$ be constant and choose a simplicial free algebra resolution $B\longrightarrow A$. Adjoining a free noncommuting generator has underlying simplicial abelian group
\[
 B\langle x\rangle\cong\bigoplus_{n\geq1}B^{\otimes_{\mathbb Z} n}.
\]
The summand for $n=2$ has first homotopy group
$\operatorname{Tor}^{\mathbb Z}_1(A,A)\cong\mathbb Z/p$, whereas
$A\langle x\rangle$ is constant. Hence
$B\langle x\rangle\longrightarrow A\langle x\rangle$ is not a weak equivalence. But $\mathbb Z\longrightarrow\mathbb Z\langle x\rangle$ is a free cofibration with cofibrant source and target. Its pushout functor is precisely adjoining $x$, so it is not a $W$-cofibration. This is the mechanism of \citet[Example 2.11, arXiv version]{Rezk2002Proper}. The error concerns the unrestricted formulation of the sentence, not Brown's excision theorem and not the finite-diagram constructions that follow.
\end{dcomment}


Je m'intéresse ici à la question d'étendre la structure à cofibrations à une catégorie de la forme $\M(I)$, et d'utiliser cette extension pour construire
\msp{page 36}
les foncteurs $u^{W}_{!}$. Notons déjà que la théorie de K. S. \textsc{Brown}, avec les compléments d'\textsc{Anderson}\footnote[20]{revus et arrangés par mes soins.}, donne essentiellement ceci :

\printedpage{21}
% head: A. Grothendieck : Dérivateurs, Ch. XIX, Édition des Universités de Montpellier II et Paris VII
% foot: Version 11 février 2019

\textbf{Proposition 3.2.} \textit{Soit $W\subset\Fl(\M)$ un localiseur associé à une structure à cofibrations $(C,W)$ dans $\M$.}
\begin{itemize}[leftmargin=3em,labelwidth=2.5em,align=right]
\item[1°)] \textit{Sur la catégorie \textup{$\Diag_0$} des ensembles ordonnés \textup{finis}, ce prédérivateur satisfait aux axiomes Dér 1 (décomposition), Dér 2 (localisation), Dér 3 (objet final), Dér 4 bis (existence de $u^{W}_{!}$), Dér 5 bis (calcul des fibres de $u^{W}_{!}$), Dér 6 bis (exactitude à gauche).}
\item[2°)] \textit{C'est encore vrai sur la sous-catégorie \textup{$\Diag_1$} dans \textup{Cat} plus grande, formée des ensembles ordonnés $I$ \textup{noethériens} tels que l'espace topologique associé à $I^{\circ}$ soit noethérien (toute suite croissante d'objets de $I$, ou d'ouverts de $I^{\circ}$, est stationnaire).}
\item[3°)] \textit{Enfin, si l'axiome transfini Cof$_4$ est satisfait, i.e. si $C$ et $TC$ sont stables par composition transfinie, alors la conclusion de \textup{1°} reste valable sur la sous-catégorie pleine \textup{$\Diag_2$} de \textup{Cat} formée des ensembles ordonnés qui sont noethériens sans plus.}
\end{itemize}

\textbf{Corollaire 3.3.} \textit{Supposons que l'on ait Cof$_3$ sous la forme plus forte d'une factorisation
\msp{page 37}
$f=pi$ \textup{fonctorielle} en $f$, ce qui est le cas en particulier si la structure à cofibrations est spéciale. Alors pour toute $I$ dans \textup{Cat}, $\M(I)$ muni de $(W(I),C(I))$ est une catégorie à cofibrations, laquelle satisfait l'axiome transfini Cof$_4$ si $(W,C)$ y satisfait. Donc on peut lui appliquer les résultats de \textup{3.2}, qui impliquent notamment que les axiomes Dér 1, 2, 3 ( ?) sont valables sur \textup{Cat} tout entier, et que $u^{W}_{!}$ existe quand $u$ est de la forme $\id_I\times v$, avec $v : J\longrightarrow K$ une flèche dans \textup{$\Diag_1$} (resp. dans \textup{$\Diag_2$}, si l'axiome transfini est satisfait).}

\textbf{3.4.} La clef de la théorie est dans l'extension de la structure à cofibrations $(W,C)$ de $\M$ à une catégorie $\M(I)$, où $I$ est un ensemble ordonné, en $(W(I),C_I)$. Les faits principaux sont énoncés dans XVII, dans les théorèmes p. 26-27, p. 50, p. 60. C'est quand $I$ est \textit{noethérien,} et qu'on suppose soit que $\mathrm{Esp}(I^{\circ})$ est également noethérien (p.ex. $I$ fini), soit que $(W,C)$ satisfait l'axiome transfini Cof$_4$, qu'on prouve que $(W(I),C_I)$ est également une structure à cofibrations sur $\M(I)$.

\dgrblock{dgrb1181}{Attaching a diagram in the order of its dependencies}
\begin{dcomment}{exposition}{p37:Les faits principaux sont énoncés dans XVII}{Attaching a diagram in the order of its dependencies}
Grothendieck's noetherian hypothesis makes an inductive construction possible: the value at an object can be attached after the values on which it depends have been constructed. With his convention $\M(I)=\operatorname{Fun}(I^\circ,\M)$, it is $I^\circ$ that serves as a direct indexing category. In a finite example this is the familiar Reedy construction, although the manuscript allows transfinite well-founded orders as well.

For a direct category $J$ and a diagram $X:J\longrightarrow\M$, let
\[
 L_jX=\mathop{\rm colim}_{(k\to j),\ k\ne j}X_k
\]
be the latching object, when the indicated colimit exists. A map $f:X\longrightarrow Y$ is a Reedy cofibration when
\[
 X_j\amalg_{L_jX}L_jY\longrightarrow Y_j
\]
belongs to $C$ for every $j$. The relative latching object incorporates all the compatibility already imposed at preceding objects. Merely requiring $X_j\longrightarrow Y_j$ to be a cofibration would omit those compatibility conditions.

For the arrow $0\to1$, the two tests are $X_0\to Y_0\in C$ and
$X_1\amalg_{X_0}Y_0\to Y_1\in C$. Given a map of arrows $X\to Y$, first factor its component at $0$ as $X_0\to Z_0\to Y_0$. Next factor the induced map $X_1\amalg_{X_0}Z_0\to Y_1$ as a cofibration followed by a weak equivalence. Its intermediate object defines $Z_1$, and the pushout supplies $Z_0\to Z_1$. The resulting $X\to Z$ is a Reedy cofibration and $Z\to Y$ is objectwise a weak equivalence. The same construction applies successively at every vertex of a finite direct category. Transfinite stages require colimits and the continuity condition Cof$_4$.

The identity $C_I\cap W(I)=(TC)_I$ records a second induction. In the arrow example, if the component at $0$ is acyclic, its pushout $X_1\to X_1\amalg_{X_0}Y_0$ is acyclic. Consequently the relative latching map is a weak equivalence exactly when $X_1\to Y_1$ is one, by two-out-of-three. This proves both directions of the identity at the second vertex. The general induction repeats this argument with the latching diagram; at a limit stage Cof$_4$ supplies the required continuity. This is why the equality, rather than the bare definition of $C_I$, is central to homotopy left Kan extension.

Corollary 3.3 uses a different construction. A functorial factorization in $\M$ can be applied to all components and all transition maps of an arbitrary diagram at once. Cof$_0$, Cof$_1$, and Cof$_2$ then hold objectwise, and functoriality supplies Cof$_3$ in $\M(I)$. This gives finite homotopy-colimit operations inside every shifted diagram category. It does not itself make the ordinary colimit functor preserve objectwise acyclic cofibrations. That additional issue is exactly what the projective classes $C\langle I\rangle$ are meant to address.
\end{dcomment}


\msp{page 38}
Un fait important pour la construction des foncteurs $u^{W}_{!}$, pour $u : I\longrightarrow I'$ une flèche entre catégories ordonnées noethériennes, est la relation
\[
C_I\cap W(I) \ = \ (TC)_I \ , \tag*{(3.4.1)}
\]
d'où résulte que l'on a, pour $u$ comme dessus,
\[
u_!(C_I\cap W(I)) \ \subset\ C_{I'}\cap W(I') \ , \tag*{(3.4.2)}
\]
ingrédient essentiel de la construction du foncteur $u^{W}_{!}$. Pour étendre cette construction au cas d'une flèche $\pm$ quelconque dans Cat, il se pose donc le problème d'étendre la définition de $C_I$, et les formules (3.4.1) et (3.4.2), au cas de catégories $\pm$ quelconques dans Cat. Or on a prouvé :

\textbf{Proposition 3.5.} \textit{Si $\Phi=\Psi^*$ est une partie \textup{Q-close} à gauche de \textup{$\Fl(\M)$}, et $I$ un ensemble ordonné noethérien, $\Phi_I$ l'extension de $\Phi$ en une partie de \textup{$\Fl(\M(I))$} définie par \textup{\textsc{Anderson}}, alors on a}
\[
\Phi_I \ = \ \Psi(I)^* \ = \ v_!(\Phi(I_0))\Qs \ , \tag*{(3.5.1)}
\]
\textit{où $v : I_0\longrightarrow I$ est le foncteur canonique de la catégorie discrète définie par \textup{$\mathrm{Ob}\,I$}.}

\printedpage{22}
% head: A. Grothendieck : Dérivateurs, Ch. XIX, Édition des Universités de Montpellier II et Paris VII
% foot: Version 11 février 2019

C'est (avec des notations différentes) le lemme 6.2.5 (p. 107) de XVII, dont je ne reprends pas ici la démonstration.

Ceci donne la clef d'une définition de $\Phi_I$ comme $\Phi\langle I\rangle$ (avec les notations des paragraphes 1, 2),
\msp{page 39}
pour $I$ quelconque dans Cat,
\[
\Phi_I \ \deq\ \Phi\langle I\rangle \ \deq\ \Psi(I)^* \ = \ v_!(\Phi(I_0))\Qs \ , \tag*{(3.5.2)}
\]
où $\Psi=\Phi_*$, et où on suppose $\Phi$ Q-clos à gauche, i.e. $\Phi=\Psi^*$. Dans le cas où on suppose seulement $\Phi$ Q-saturé transfini sans plus, on aurait envie de poser
\[
\Phi_I \ \deq\ \text{Q-saturé transfini de } v_!(\Phi(I_0)) \ , \tag*{(3.5.3)}
\]
et je présume que pour $I$ ordonné noethérien, cela coïncide avec la définition d'\textsc{Anderson}. Mais je vais me borner pour la suite au cas où $C=\Phi$ est Q-clos à gauche.

Revenons alors au cas d'une structure à cofibrations $(W,C)$ dans $\M$, en supposant $C$ et $TC=C\cap W$ Q-clos à gauche. Définissons pour toute $I$ dans Cat $C_I=C\langle I\rangle$ comme dans (3.5.2). On a alors, pour $u : I\longrightarrow I'$ une flèche dans Cat,
\[
u_!(C\langle I\rangle) \ \subset\ C\langle I'\rangle \ , \tag*{(3.5.4)}
\]
mais on n'y a guère de raison d'espérer avoir aussi (3.4.2), sans disposer de la formule
\[
C\langle I\rangle\cap W(I) \ \overset{?}{=}\ (TC)\langle I\rangle \qquad\qquad \text{\footnote[21]{\textbf{N.B.} On a l'inclusion $\supset$, car $v^*v_!TC(I_0)\subset TC(I_0)$.}}. \tag*{(3.5.5)}
\]
On a prouvé que cette formule est valable si $I$ est ordonné noethérien (XVII th. 2.1, p. 26-27). Cela
\msp{page 40}
nous amène à la

\textbf{Question 3.6.} \textit{Soit $\M$ une catégorie stable par petites limites, et soit $(W,C)$ une structure à cofibrations dans $\M$ telle que $C$ et $TC=C\cap W$ soient Q-clos à gauche. Soit $I$ un objet de \textup{Cat}, a-t-on alors}
\[
C\langle I\rangle\cap W(I) \ = \ (TC)\langle I\rangle \text{ ?} \tag*{(3.6.1)}
\]

\dgrblock{dgrb1182}{Lifting properties retain the coherence}
\begin{dcomment}{exposition}{p40:telle que $C$ et $TC=C\cap W$ soient Q-clos à gauche}{Lifting properties retain the coherence}
Grothendieck replaces a geometric attachment rule by a characterization that makes sense for every small category. Put $TF=C^\perp$, $F=(TC)^\perp$, and
\[
 C\langle I\rangle={}^\perp TF(I),\qquad
 (TC)\langle I\rangle={}^\perp F(I).
\]
The lifting problem here is a problem in the category of diagrams. A diagonal must commute with every arrow of $I$; a collection of independently chosen componentwise diagonals is insufficient. This is the coherence retained by the projective lifting class.

Let $v:I_0\to I$ include the discrete category of objects. The adjunction $v_!\dashv v^*$ gives
\[
 (v_!\Phi(I_0))^\perp
   =\{p:v^*p\in\Phi(I_0)^\perp\}
   =\Psi(I),
\]
where $\Psi=\Phi^\perp$ and $\Phi$ is Q-closed. The last equality follows because a lifting problem in a discrete family is solved component by component; equivalently, one tests a single component and puts identity maps in the others. Taking left orthogonals proves (3.5.2). The same adjunction calculation for $u_!\dashv u^*$ gives (3.5.4), since restriction preserves the objectwise right class. These arguments require existence of the ordinary left Kan extension, not a factorization system on diagrams.

The specification of $\Psi$ matters, just as in the arbitrary-test-class defect in (1.14) discussed earlier. If Proposition 3.5 is read as allowing an arbitrary presentation $\Phi={}^\perp\Psi$, its first equality needs qualification. In sets, take $\Psi=\{\emptyset\to1\}$. Then $\Phi$ consists of the maps with nonempty source together with $1_{\emptyset}$. On the direct arrow category $J=[1]$, the diagram map
\[
 f:(\emptyset\longrightarrow1)\longrightarrow(1\xrightarrow{1}1)
\]
fails the relative latching condition at $0$. Yet it lifts vacuously against $\Psi(J)$: that class contains only the map from the constant empty diagram to the constant singleton diagram, and the source of $f$ has no map to the former. Thus ${}^\perp\Psi(J)$ is too large for this nonclosed $\Psi$. The canonical choice $\Psi=\Phi^\perp$, explicitly imposed in (3.5.2), removes this issue. It leaves the intended construction and the subsequent propositions intact.

For a well-founded direct indexing order, the latching description and this lifting description agree by induction on objects: at stage $j$, the already constructed diagonals determine a map out of the relative latching object, and its lifting property supplies the next component. The converse can be checked explicitly for the posets in Proposition 3.5. Given a lifting problem from the relative latching map at $j$ to $p:A\to B$, use the diagram map which is $1_A$ at strict predecessors of $j$, is $p$ at $j$, and is the identity of the terminal object elsewhere. The map $L_jY\to A$ supplies its predecessors. A coherent diagonal then gives exactly the requested relative lift. The canonical right class $\Phi^\perp$ contains all identities, as this test requires. This proves the comparison under the stated existence assumptions.

There is also a direct cellular verification of the conjecture (3.5.3) whenever the small colimits needed for its transfinite construction exist. Assume that $\Phi$ contains isomorphisms and is closed under pushouts and small transfinite composition. The following is the direct-poset version of the standard Reedy cellular construction; compare \citet[Proposition~6.3 and Corollary~6.7]{RiehlVerity2014Reedy}. Let $J=I^\circ$ be the well-founded poset and well order its objects so that every strict predecessor is earlier. Write $j_!$ for the left adjoint of evaluation, so $(j_!A)_k=A$ if $j\leq k$ and is initial otherwise. A map $j_!a$, $a\in\Phi$, has relative latching map $a$ at $j$ and an identity elsewhere.

The class of maps whose relative latching maps lie in $\Phi$ is closed under pushouts, since forming a relative latching map takes a pushout of diagram maps to a pushout of that relative map. For transfinite closure, let $(X^\alpha)_{\alpha\leq\lambda}$ be a continuous chain of such maps. For fixed $j$ put
\[
 Q_\alpha=X^\alpha_j\amalg_{L_jX^\alpha}L_jX^\lambda.
\]
Each $Q_\alpha\to Q_{\alpha+1}$ is a pushout of the relative latching map of $X^\alpha\to X^{\alpha+1}$. The $Q$-chain is continuous because latching objects are colimits. Its starting object is the relative-latching domain of $X^0\to X^\lambda$, and $Q_\lambda=X^\lambda_j$. Transfinite closure of $\Phi$ therefore proves the assertion. Consequently this relative-latching class contains the transfinite saturation of all the maps $j_!a$.

For the reverse inclusion, take $f:X\to Y$ with all relative latching maps in $\Phi$. Starting from $X$, successively attach the value $Y_j$ at each object $j$. When $j$ is reached, all strict predecessors already equal their values in $Y$, and the current value at $j$ is
\[
 X_j\amalg_{L_jX}L_jY.
\]
Push out the free diagram map $j_!(X_j\amalg_{L_jX}L_jY\to Y_j)$ along its evaluation-adjoint map to the current diagram. This replaces the value at $j$ by $Y_j$ and leaves every earlier value fixed. At limit ordinals take colimits. The resulting composite ends at $Y$, since every component is eventually fixed to $Y_j$. Each step is a pushout of $j_!a$ for some $a\in\Phi$. Finally these one-object free maps belong to $v_!\Phi(I_0)$ by putting identity maps of the initial object in the other components; conversely a whole family is a coproduct of such one-object maps, formed by successive pushouts. This proves precisely the proposed transfinite-saturation description, without adding retract closure. It is a presentation by a class of cells, not an assertion that a small set of them suffices. If the required colimits have not been assumed, that existence qualification must remain explicit.

The easy inclusion in (3.5.5) is structural. Under the small-completeness hypothesis of 3.6, evaluation at an object $j$ has a right adjoint whose values are products indexed by the arrows into $j$ in $I^\circ$. Right lifting classes are closed under products, so this right adjoint sends a map of $F$ to an objectwise $F$-map. If $f\in{}^\perp F(I)$, adjunction therefore implies $f_j\in{}^\perp F=TC$. Hence projective $TC$-cofibrations are objectwise in $TC$, and therefore in $W(I)$; also $TC\subset C$ gives $(TC)\langle I\rangle\subset C\langle I\rangle$. The difficult inclusion says that an objectwise weak equivalence which already satisfies one coherent lifting condition satisfies the stronger one. Existence of abstract homotopy Kan extensions, established later, does not by itself identify these two specifically prescribed classes of strict diagrams.
\end{dcomment}


On a vu au paragraphe précédent qu'il en est ainsi si $(W,C)$ s'insère dans un triple de Quillen clos $(W,C,F=(TC)_*)$, et quand on suppose de plus que les couples
\[
C\langle I\rangle \ \longleftrightarrow\ TF(I) \ , \qquad TC\langle I\rangle \ \longleftrightarrow\ F(I)
\]
sont factorisant, auquel cas $(W(I),C\langle I\rangle,F(I))$ est une structure de Quillen close dans $\M(I)$. D'ailleurs, $(W,C,F)$ est une structure de Quillen (nécessairement close) si et seulement si les couples $C\longleftrightarrow \underbrace{TF}_{\deq C_*}$, $TC\longleftrightarrow \underbrace{\ F\ }_{\deq (TC)_*}$ eux-mêmes sont factorisants, et si % [sic] factorisant -> factorisants, agreement with « les couples »
\printedpage{23}
% head: A. Grothendieck : Dérivateurs, Ch. XIX, Édition des Universités de Montpellier II et Paris VII
% foot: Version 11 février 2019
de plus $TF\subset W$ (ce qui est automatique si $W=W^{\natural}$, car on aura en tous cas, par Cof$_4$, $TF\subset W^{\natural}$).

Mais peut-on s'en tirer avec moins, pour pouvoir conclure à (3.6.1) ? Ne suffirait-il pas de savoir que $TC\langle I\rangle\longleftrightarrow F(I)$ est factorisant (ce qui n'a guère de chances d'être vrai pour une $I$ $\pm$ générale, si on n'a déjà $TC\fact F$, factorisation dans $\M$ lui-même). C'est bien le cas :

\msp{page 41}
\textbf{Proposition 3.7.} \textit{Sous les conditions générales de \textup{(3.6)} sur $\M$, $W$, $C$, $I$, supposons que le couple $TC\langle I\rangle\longleftrightarrow F(I)$ (où $F\deq (TC)_*$) soit factorisant, i.e. $TC\langle I\rangle\fact F(I)$. Alors la formule \textup{(3.6.1)} est vraie.}

En effet, si $f\in C\langle I\rangle\cap W(I)$, on factorise $f$ en $pi$, avec $i\in TC\langle I\rangle$, $p\in F(I)$. Comme $f,i\in W(I)$, on aura $p\in W(I)$, i.e. $p\in TF(I)$, et à cause de $\underbrace{C\langle I\rangle}_{\ni\, f}\harp\underbrace{TF(I)}_{\ni\, p}$, il existe $p : Y\longrightarrow X'$ telle que $pq=\id_Y$, $qf=i$, % reading resolved: DER_XIX_corrections.csv; % [sic] p : Y -> q : Y, the arrow Y -> X' is q (see the diagram and pq = id_Y)
\[
\begin{tikzcd}[column sep=small, row sep=large]
X \arrow[rr, "f\in C\langle I\rangle\cap W(I)"] \arrow[dr, "i\in TC\langle I\rangle"'] & & Y \arrow[dl, dotted, shift right=1ex, "q"'] \\
& X' \arrow[ur, "p\in F(I)"'] \ , &
\end{tikzcd}
\]
donc $f$ est facteur direct de $i\in TC\langle I\rangle$. Comme $TC\langle I\rangle$ est stable par facteurs directs, on gagne.

\textbf{Corollaire 3.8.} \textit{Soit $(W,C,F)$ une structure de Quillen close dans $\M$, et supposons que $TC=C\cap W$ soit une partie Q-spéciale de $\M$, i.e. le couple $TC\fact F$ un couple spécial (i.e. $\M$ est \textup{spéciale}, et il existe une petite partie $TC_0$ Q-génératrice à gauche pour $TC$). Alors pour toute $I$ dans \textup{Cat}, on a}
\[
C\langle I\rangle\cap W(I) \ = \ TC\langle I\rangle \ . \tag*{(3.8.1)}
\]

\dgrblock{dgrb1183}{The lifting lemma inside the retract argument}
\begin{dcomment}{exposition}{p41:Q-génératrice à gauche pour $TC$}{The lifting lemma inside the retract argument}
Grothendieck has found that only one diagram factorization is needed for (3.6.1). The proof is sound, but the line identifying a weak equivalence in $F$ with a map lifting against $C$ uses a short lemma that is useful beyond the present proposition.

Suppose $(W,C)$ satisfies Cof$_0$--Cof$_3$, and set $F=(C\cap W)^\perp$. Then
\[
 F\cap W\subset C^\perp.
\]
To prove this, consider a square from $a:A\to B$ in $C$ to $p:X\to Y$ in $F\cap W$, with top map $h:A\to X$ and bottom map $k:B\to Y$. Form $Z=X\amalg_A B$, write $j:X\to Z$, and let $g:Z\to Y$ be induced by $p$ and $k$. Factor $g=q\ell$ by Cof$_3$, where $\ell:Z\to T$ is in $C$ and $q:T\to Y$ is in $W$. The composite $\ell j$ belongs to $C$, and $q\ell j=p$ shows by two-out-of-three that $\ell j\in TC$. Since $p\in F$, the square with left side $\ell j$, right side $p$, top $1_X$, and bottom $q$ admits $s:T\to X$ satisfying $s\ell j=1_X$ and $ps=q$. If $b:B\to Z$ is the other pushout map, then $s\ell b$ is the desired lift. This proves the inclusion without a second factorization in $\M$, and without retract closure of $W$.

Now take $f:X\to Y$ in $C\langle I\rangle\cap W(I)$ and factor it as $f=pi$ with $i\in(TC)\langle I\rangle$ and $p\in F(I)$. The already established easy inclusion makes $i$ an objectwise weak equivalence, hence $p\in F(I)\cap W(I)$. Applying the preceding lemma at every object gives $p\in TF(I)=C^\perp(I)$. The defining coherent lifting property of $f$ produces a natural map $q:Y\to X'$ with $qf=i$ and $pq=1_Y$. The arrow-category maps $(1_X,q)$ and $(1_X,p)$ exhibit $f$ as a retract of $i$. Left lifting classes are retract closed, so $f\in(TC)\langle I\rangle$.

The printed prose writes $p:Y\to X'$ for the new diagonal; the diagram and both identities show that $q$ is intended. The earlier reference on ms.~p.~40 to Cof$_4$ in the assertion $TF\subset W^{\natural}$ is likewise a reference slip: the relevant axiom is Cof$_3$, as in Proposition 3.1. Neither slip changes the argument. Corollary 3.8 is then exactly the small object theorem applied to $TC$ in diagram categories. It proves the equality without assuming that $C$ has a small generating set; the separate difficulty is obtaining the other factorization needed for a complete diagram model structure.
\end{dcomment}


\msp{page 42}
En apparence, on n'a pas eu à supposer que $C$ soit elle aussi une partie Q-spéciale de $\Fl(\M)$. Mais comme on a besoin de la factorisation pour $C\longleftrightarrow TF$ (pour avoir au moins une structure de cofibrations sur $\M$), on n'a guère le moyen de l'établir sans supposer $C$ spéciale. Sans compter que pour pouvoir expliciter (3.8.1) pour la construction de $u^{W}_{!}$, on a besoin de savoir que $(W(I),C\langle I\rangle)$ définit bien une structure à cofibrations sur $\M(I)$, ce qui ne peut guère se faire (il me semble) sans avoir la factorisation pour $C\langle I\rangle\longleftrightarrow TF(I)$ également ; chose que je n'arrive à établir pour un $I$ général dans Cat, que moyennant l'hypothèse de spécialité sur $C$ également.

\textbf{3.9.} Pour en avoir le cœur net, il faudrait quelques exemples concrets et illuminants, du type suivant :

a) \texttt{[Il n'y a pas de b).]} Structure de Quillen close $(W,C,F)$ sur $\M$, telle que pour le prédérivateur associé, $u^{W}_{!}$ n'existe pas toujours, ou $u^{W}_{*}$ n'existe par toujours. Je pense notamment au cas d'une flèche $u : I\longrightarrow I'$ avec $I'=e$ (présumant que si dans ce cas, $u^{W}_{!}$ resp. $u^{W}_{*}$ existe, alors il existe dans tous les cas). Un tel exemple serait encore plus % [sic] n'existe par toujours -> n'existe pas toujours, typing error

% ---- from der-ch19_p026-028.tex
\printedpage{24}
% head: A. Grothendieck : Dérivateurs, Ch. XIX, Édition des Universités de Montpellier II et Paris VII
% foot: Version 11 février 2019
\msp{page 43}
instructif, bien sûr, si $\mathcal{M}$ est stable par petites $\limind$ , que ses objets soient accessibles etc. et que $W$ soit L-accessible (i.e. stable par petites $\limind$ filtrantes, et de plus accessible \dots). Le premier candidat pour $I$ auquel je pense serait un $B_G$ ($G$ un groupe, p.ex. $G=\pm1$), c'est en effet l'exemple aux antipodes des ensembles ordonnés, lesquels, eux, ne semblent pas faire trop de problème (moyennant l'hypothèse noethérienne tout au moins).

Les exemples suivants ne sont pas aussi convaincants que ça, hélas!

\textbf{Exemple 3.10.} Soit $\mathcal{M}$ une catégorie quelconque, alors les couples $(\Fl(\mathcal{M}),\Is(\mathcal{M}))$ et $(\Is(\mathcal{M}),\Fl(\mathcal{M}))$ sont tous deux des couples de Quillen,
\[
\Fl(\mathcal{M}) \Qcouple \Is(\mathcal{M}) \ , \qquad \Is(\mathcal{M}) \Qcouple \Fl(\mathcal{M}) \ . \tag*{(3.10.1)} % reading resolved: DER_XIX_corrections.csv
\]
Il s'ensuit que si on pose
\[
W \ = \ \Is(\mathcal{M}) \ , \qquad C \ = \ F \ = \ \Fl(\mathcal{M}) \ , \tag*{(3.10.2)}
\]
$(W,C,F)$ est un triple de Quillen clos \textsuperscript{(}\footnote[22]{au moins si on admet que $\mathcal{M}$ est stable par $\limind$ et $\limproj$ \textit{finies.}}\textsuperscript{)}. D'ailleurs pour toute $I$ dans Cat, on a $W(I)=\Is(\mathcal{M}(I))$, donc
\[
\Ho_W(I) \ = \ \mathcal{M}(I) \ . \tag*{(3.10.3)}
\]
Donc la question de l'existence des foncteurs $u_!^W$, $u_*^W$ se réduit à celle des foncteurs $u_!$, $u_*$ habituels, adjoints à gauche et à droite de $u^* : \mathcal{M}(I') \longrightarrow \mathcal{M}(I)$.

\dgrblock{dgrb1184}{Finite completeness and the modern model category theorem}
\begin{dcomment}{status}{p43:Donc la question de l'existence des foncteurs}{Finite completeness and the modern model category theorem}
Example 3.10 isolates the completeness issue. Every map factors as itself followed by an identity, or as an identity followed by itself. A map with the lifting property against every map lifts against itself, and that diagonal is its inverse. Consequently the two displayed pairs really are closed factorization pairs in any category. With finite limits and colimits they give the historical, finitely complete version of Quillen's axioms. Their localization does nothing, so existence of homotopy Kan extensions is precisely existence of ordinary Kan extensions. For a functor between finite categories, the relevant comma categories are finite and the assumed finite limits and colimits suffice. For arbitrary indexing categories those assumptions say nothing about the required infinite constructions. The small-generation calculation for the corresponding example in sets is already given in the commentary to ms.~p.~32.

The terminology now usually includes \emph{all} small limits and colimits in the definition of a model category. Under that convention the requested example in 3.9 does not exist. \citet[Theorem~6.11]{Cisinski2003} proves the derivator assertion for an arbitrary Quillen model category with the required small limits; duality gives the assertion about colimits. Neither cofibrant generation nor local presentability is required. The distinction matters: this theorem supplies homotopy Kan extensions without asserting that the specific projective Quillen structure written in Question 3.6 exists on every diagram category. The bounded-complex example below therefore addresses the older finite-completeness convention, while Question 3.6 retains its more specific lifting-class content.

Grothendieck's suggestion that terminal projections should suffice anticipates the pointwise Kan-extension formulas. In the established derivator framework the value of a Kan extension is a homotopy limit or colimit over a comma category. For an arbitrary prederivator, however, adjoints to terminal projections alone do not supply the coherent pointwise construction or its base-change isomorphisms. Those are additional assertions proved by the model-category theorem, rather than a formal consequence of representability at the terminal category.
\end{dcomment}


\msp{page 44}
On sait bien que l'existence de ces foncteurs pour $u$ donné, fait appel à l'existence de certaines $\limind$ resp. $\limproj$ dans $\mathcal{M}$; p.ex. si $I'=e$, il est nécessaire et suffisant que les $\limind$ (resp. $\limproj$ ) de type $I^\circ$ soient représentables dans $\mathcal{M}$. Mais avec les hypothèses faites sur $\mathcal{M}$, il n'y a pas de problème si $u : I \longrightarrow I'$ est un foncteur entre catégories \textit{finies.} On va donner un exemple où il n'en est plus ainsi.

\textbf{Exemple 3.11.} Soit $k$ un anneau commutatif, et $\mathcal{M}=K^{\mathrm{b}}(k\text{-Ab})$ la catégorie des complexes de cochaînes bornés de $k$-modules, $W$ l'ensemble des quasi-isomorphismes, donc
\[
\Ho_W \ = \ D^{\mathrm{b}}(k\text{-Ab}) \ ,
\]
catégorie dérivée bornée de la catégorie abélienne des $k$-modules. On sait, en utilisant l'existence d'assez de projectifs dans $k$-Ab, que cette structure est associée à une structure de Quillen close $(W,C,F)$, décrite par \textsc{Quillen}
\msp{page 45}
dans Homotopical Algebra -- à cela près qu'il prend $D^-(k)$ au lieu de $D^{\mathrm{b}}(k)$. Peut-être, pour vérifier les axiomes, doit-on supposer que $k$ est de dimension cohomologique finie.

Prenons $I=B_G$, où $G$ est un groupe (p.ex. $G=\{\pm1\}$), et $u : I \longrightarrow e$. Alors $\mathcal{M}(I)$ est la catégorie des complexes bornés de $k[G]$-modules. On voit que $u_!^W$ resp. $u_*^W$ existe si et seulement si pour tout tel complexe borné, ou simplement pour un $k[G]$-module sans plus $M$, le complexe $L\Gamma^G(M)$ resp. $R\Gamma_G(M)$ appartient à $D^{\mathrm{b}}(k)$, ou ce qui revient au même,
\printedpage{25}
% head: A. Grothendieck : Dérivateurs, Ch. XIX, Édition des Universités de Montpellier II et Paris VII
% foot: Version 11 février 2019
si et seulement si on a
\[
\begin{aligned}
H_i(G,M) \ &= \ 0 \quad \text{pour } i \gg 0 \ ,\\
\text{resp.}\quad H^i(G,M) \ &= \ 0 \quad \text{pour } i \gg 0 \ ,
\end{aligned}
\]
pour tout $k[G]$-module $M$. Ce n'est pratiquement jamais le cas; quand on prend p.ex. $k=\mathbf{Z}$ et $G$ groupe fini, cela n'est le cas que si $G=e$. L'exemple le plus économique me paraît $G=\pm1$, $k=\mathbf{F}_2$.

\dgrblock{dgrb1185}{A finite group can force an unbounded resolution}
\begin{dcomment}{examples}{p45:L'exemple le plus économique me paraît}{A finite group can force an unbounded resolution}
The notation $K^{\mathrm b}(k\text{-Ab})$ is explicitly defined here as the category of bounded complexes and chain maps. It is what is now usually denoted $C^{\mathrm b}(k\text{-Mod})$, before quotienting by chain homotopy. The sentence about enough projectives needs the qualification which Grothendieck immediately adds: an unbounded projective resolution does not provide a bounded cofibrant replacement. It is enough for his example to take the proposed field $k=\mathbf F_2$, where every bounded acyclic complex is contractible.

Here is a complete verification of the finite Quillen structure in that case. Put $W=$ quasi-isomorphisms, $C=$ degreewise injections, and $F=$ degreewise surjections. A bounded contractible complex of vector spaces is a finite direct sum of disk complexes and is both projective and injective in the category of complexes. Thus a surjective quasi-isomorphism, whose kernel is contractible, has the right lifting property against every injection; dually an injective quasi-isomorphism has the left lifting property against every surjection. For a bounded complex $X$, the cone $I(X)=\operatorname{Cone}(1_X)$ is contractible and has a natural degreewise injection $X\to I(X)$. Dually $P(Y)=\operatorname{Cone}(1_Y)[-1]$ has a natural degreewise surjection $P(Y)\to Y$. For $f:X\to Y$ the factorizations
\[
 X\xrightarrow{(f,i_X)}Y\oplus I(X)\longrightarrow Y,
 \qquad
 X\longrightarrow X\oplus P(Y)\xrightarrow{(f,p_Y)}Y
\]
have types $(C,F\cap W)$ and $(C\cap W,F)$ respectively. Retract closure and two-out-of-three are immediate from degreewise linear algebra and cohomology. All objects and constructions remain bounded. Finite limits and colimits exist, whereas, for example, a coproduct of $k[n]$ over all integers cannot be bounded: its universal injections and retractions would force nonzero cohomology in every degree.

For $G=C_2$ write $A=k[G]=k[\epsilon]/(\epsilon^2)$, with $\epsilon=g-1$. Since $BG$ has one object, a diagram of bounded complexes is exactly a bounded complex of $A$-modules. There is no uniform-boundedness issue here; for a general category with infinitely many objects, an objectwise bounded diagram need not be bounded uniformly. Thus the relevant restriction functor is
\[
 p^*:D^{\mathrm b}(k)\longrightarrow D^{\mathrm b}(A),
 \qquad V\longmapsto V\text{ with trivial }G\text{-action}.
\]
The exact projective resolution
\[
 \cdots\xrightarrow{\epsilon}A\xrightarrow{\epsilon}A
 \xrightarrow{\epsilon}A\longrightarrow k\longrightarrow0
\]
is exact because $\ker(\epsilon)=\epsilon A$. Tensoring it with $k$, or applying $\operatorname{Hom}_A(-,k)$, kills every differential. Consequently
\[
 \operatorname{Tor}^{A}_n(k,k)=k
 \quad\hbox{and}\quad
 \operatorname{Ext}^{n}_A(k,k)=k\qquad(n\geq0).
\]
This proves nonexistence of the adjoints themselves, not merely failure of one proposed construction. If $R$ were right adjoint to $p^*$, then
\[
 H^n(Rk)=\operatorname{Hom}_{D(k)}(k[-n],Rk)
 \cong\operatorname{Hom}_{D(A)}(k[-n],k)=k
\]
for every $n\geq0$, contradicting boundedness. If $L$ were left adjoint, then
\[
 \operatorname{Hom}_{D(k)}(Lk,k[n])
 \cong\operatorname{Hom}_{D(A)}(k,k[n])=k
\]
for every $n\geq0$. Over a field, the left side vanishes outside a finite range for every bounded $Lk$, another contradiction.

The source's general boundedness criterion can also be justified. The unbounded derived adjunctions are tensoring with the trivial module and derived invariants. If they take every module to a bounded complex, finite truncation triangles show that they do so for every bounded complex. Their restrictions are then the required adjoints. Conversely, the shifts of $k$ detect all cohomology of the proposed right adjoint. For the left adjoint, test maps into shifts of an injective cogenerator of $k$-modules: these detect the homology of the derived tensor product and force its boundedness. For $k=\mathbf Z$ and a nontrivial finite group, choose a subgroup $C_p$ of prime order and induce its trivial $\mathbf F_p$-module to $G$. The cyclic periodic resolution and the induction/coinduction identities give nonzero homology and cohomology in arbitrarily large degrees; see \citet[Lecture~3]{BrownCohomologyLectures}. Thus the stated finite-group obstruction is valid. Its mechanism is the arbitrarily long composable strings in $BG$, despite the finiteness of its object and arrow sets.
\end{dcomment}


\msp{page 46}
Donc ici $u_!^W$ et $u_*^W$ n'existent pas, pour $I \longrightarrow e$ avec $I$ catégorie finie. Mais aussi $\mathcal{M}$ n'est pas stable par petites $\limind$ .

\textbf{Exemple 3.12.} Soit $\mathcal{M}$ une catégorie muni d'un localiseur $W$ fortement saturé, tel que $\mathcal{M} \longrightarrow \underbrace{\mathcal{M}W^{-1}}_{=\mathcal{H}}$ admette un adjoint -- à gauche disons. On a donc une situation de deux foncteurs adjoints % [sic] muni -> munie, agreement with « une catégorie »
\[
\begin{tikzcd}[column sep=huge]
\mathcal{H} \arrow[r, shift left, "v_!\ \text{pleinement fidèle}"] & \mathcal{M} \arrow[l, shift left, "v^*\ \text{localisation}"]
\end{tikzcd}
\ , \tag*{(3.12.1)}
\]
et $W=(v^*)^{-1}(\mathrm{Iso}(\mathcal{H}))$. On trouve que pour toute $I$ dans Cat, on a
\[
\mathrm{HO}_W(I) \ \simeq \ \mathcal{H}(I) \ \overset{\text{déf}}{=} \ \mathrm{Hom}(I^\circ,\mathcal{H}). \tag*{(3.12.2)} % [sic] HO_W -> Ho_W, notation Ho_W used in (3.10.3) and Exemple 3.11
\]
On peut donc songer à faire un exemple où dans $\mathcal{M}$ les petites $\limind$ et $\limproj$ existent, mais non dans $\mathcal{H}$, de sorte que les $u_!^W$ ou $u_*^W$ n'existent pas toujours. Mais ce semble difficile à construire, vu que dans la situation envisagé, on voit facilement que les petites $\limproj$ existent dans $\mathcal{H}$. Donc si dans $\mathcal{H}$ % [sic] envisagé -> envisagée, agreement with « la situation »
\msp{page 47}
il existe une petite famille cogénératrice (et pour ceci, il suffit qu'une telle famille existe dans $\mathcal{M}$), il y aura également des petites $\limind$ , en vertu de SGA 4 I 8.12.7. De même, elles existent si $\mathcal{H}$ est une sous-catégorie pleine de $\mathcal{M}$ \textit{stable par petites} $\limind$ , et il paraît pratiquement impossible de faire un exemple (3.12.1) sans que $\mathcal{H}$ ait cette propriété de stabilité. Sans compter qu'il y a même la perplexité suivante :

\marginpar{! \textbar}\textbf{Question 3.13.} \textit{Existe-t-il une $\mathfrak{U}$-catégorie (\,$\mathfrak{U}$ un univers donné) où les petites $\limind$ soient représentables, sans qu'il en soit de même des petites $\limproj$ ?} \textbf{N.B.} Une telle catégorie n'admet pas de petite famille génératrice, a fortiori elle n'est pas équivalente à une petite catégorie.

\dgrblock{dgrb1186}{A coreflective localization inherits both kinds of limits}
\begin{dcomment}{exposition}{p47:Une telle catégorie n'admet pas de petite famille génératrice}{A coreflective localization inherits both kinds of limits}
The proposed difficulty in Example 3.12 is removed by a formal adjunction argument. Write $L:\mathcal H\to\M$ for the fully faithful left adjoint and $R:\M\to\mathcal H$ for the localization. The unit is an isomorphism. The counit $\varepsilon_X:LRX\to X$ belongs to $W$, since $R\varepsilon_X$ is an isomorphism by the triangle identity. These facts hold objectwise in every diagram category. Therefore $L^I$ and $R^I$ descend to inverse equivalences
\[
 \M(I)[W(I)^{-1}]\simeq\mathcal H(I),
\]
which proves (3.12.2) directly and explains why the example becomes a question about ordinary completeness.

If $\M$ is complete and $D:J\to\mathcal H$ is small, then
\[
 R\bigl(\lim_J LD\bigr)
\]
is its limit in $\mathcal H$: applying $\operatorname{Hom}_{\mathcal H}(H,-)$ and the adjunction identifies the required universal property with the limit universal property in $\M$. If $\M$ is cocomplete, let $X=\operatorname{colim}_J LD$. Every cocone map $LD(j)\to X$ has a unique factorization through $\varepsilon_X:LRX\to X$, because its source lies in the image of $L$. These factorizations form a cocone, hence induce $g:X\to LRX$ with $\varepsilon_Xg=1_X$. Postcomposition with $\varepsilon_X$ is a bijection on maps from any object in the image of $L$. Applying this with source $LRX$ gives $g\varepsilon_X=1_{LRX}$. Thus $\varepsilon_X$ is an isomorphism and $X$ belongs to the essential image of $L$. In particular that image is automatically closed under all small colimits, and $R\operatorname{colim}LD$ is the colimit in $\mathcal H$.

Consequently a full coreflective subcategory of a bicomplete category is bicomplete. No cogenerating family is needed, and the envisaged counterexample cannot have the adjunction (3.12.1). The source's auxiliary observation about cogenerators is nevertheless correct: if $\{C_s\}$ conservatively detects arrows in $\M$ by $\operatorname{Hom}_{\M}(-,C_s)$, then $\{RC_s\}$ does so in $\mathcal H$, by adjunction and full faithfulness of $L$. The cited representability criterion, \citet[Expos\'e~I, 8.12.7]{SGA4I1972}, supplies another sufficient argument when that hypothesis is available; the colimit argument above shows why it is unnecessary here.
\end{dcomment}


\dgrblock{dgrb1187}{All small suprema without a largest ordinal}
\begin{dcomment}{examples}{p47:Une telle catégorie n'admet pas de petite famille génératrice}{All small suprema without a largest ordinal}
Question 3.13 has a direct affirmative answer, with its universe distinction essential. Let $\operatorname{Ord}_{\mathfrak U}$ have the ordinals belonging to $\mathfrak U$ as objects, with one arrow $\alpha\to\beta$ when $\alpha\leq\beta$. It is locally $\mathfrak U$-small: every Hom-set is empty or a singleton. For a $\mathfrak U$-small diagram its colimit is the supremum of the ordinals occurring in the diagram. This supremum belongs to $\mathfrak U$ by the universe axioms; for the empty diagram it is $0$. There is no terminal object, since $\alpha+1$ is again an object and is larger than $\alpha$. Thus all small colimits exist but even the empty limit does not. The category is necessarily large relative to $\mathfrak U$.

The note about generators is correct with the convention used in SGA~4. In \citet[Expos\'e~I, 7.1]{SGA4I1972}, a generating family means that the associated represented functors jointly \emph{reflect isomorphisms}. Given any small family of ordinals $\{\alpha_s\}$, put $\beta=\sup_s\alpha_s$. Every map
\[
 \operatorname{Hom}(\alpha_s,\beta)
 \longrightarrow\operatorname{Hom}(\alpha_s,\beta+1)
\]
is a bijection between singletons, but $\beta\to\beta+1$ is not an isomorphism. Hence there is no small generating family in this conservative sense. A merely faithful family of Hom-functors would be a different convention: every poset has at most one arrow with given endpoints, so faithfulness alone says almost nothing here. It would be incorrect either to reject Grothendieck's note using that weaker convention or to conclude from this example that a small cocomplete category need not be complete.
\end{dcomment}


\msp{page 48}
\textbf{Question 3.14.} \textsc{Anderson} prétend que pour le prédérivateur défini par une catégorie à cofibrations satisfaisant l'axiome Cof$_4$ transfini, les foncteurs $u_!^W$ existent. Sa démonstration est totalement canulée, néanmoins je n'arrive pas à faire de contre-exemple. Le résultat qu'il annonce serait-il vrai? Je ne connais pas même d'exemple d'un localiseur $W$, dans une catégorie $\mathcal{M}$ \textit{avec petites $\limind$ et petites $\limproj$ ,} tel que $u_!^W$ ou $u_*^W$ n'existent pas toujours -- mais je suis persuadé qu'il y a des exemples (même avec $\mathcal{M}$ une petite catégorie?). J'en suis moins convaincu si on suppose $W$ stable par $\limind$ filtrantes, et au besoin, de plus, que les éléments de $\mathcal{M}$ sont accessibles (plus éventuellement d'autres propriétés catégoriques anodines sur $\mathcal{M}$, telles l'existence d'une filtration
\printedpage{26}
% head: A. Grothendieck : Dérivateurs, Ch. XIX, Édition des Universités de Montpellier II et Paris VII
% foot: Version 11 février 2019
\msp{page 49}
cardinale etc.)

\dgrblock{dgrb1188}{The cofibration assertion is true under the stated continuity hypothesis}
\begin{dcomment}{later}{p48:Le résultat qu'il annonce serait-il vrai?}{The cofibration assertion is true under the stated continuity hypothesis}
The first question in 3.14 now has a positive answer. Modern cofibration-category theory separates this assertion from the existence of projective Quillen structures. In the presence of the initial object and pushouts in the standing framework, Cof$_4$ implies the coproduct and sequential-colimit axioms of an Anderson--Brown--Cisinski cofibration category. To see the coproduct assertion, well order a small family of cofibrant objects, attach its members one at a time by pushout from the initial object, and take colimits at limit stages. The resulting transfinite composite is the coproduct. Applied to a family of cofibrations, the same construction proves closure of cofibrations under coproducts; for acyclic cofibrations use the asserted transfinite closure of $TC$. Countable sequential composites are already a special case of Cof$_4$.

\citet[Lemma~1.6.2 and Theorem~10.3.2]{RadulescuBanu2009} gives the corresponding axioms and proves that every such cofibration category defines a strong derivator with all homotopy left Kan extensions. His Theorem~10.3.3 includes arbitrary Quillen model categories. A complementary formulation uses derivable categories: the dual of \citet[Corollary~6.21]{Cisinski2010} supplies the colimit theorem, while that corollary itself treats the limit direction in Cisinski's conventions. These results require neither a small set of generating cofibrations nor functorial factorization. They settle the existence question raised here, without retroactively validating the proof which Grothendieck criticizes.

The broader request in the same paragraph drops the cofibration structure entirely. It has a different answer, even when the underlying category and localizer have every finiteness and accessibility property explicitly mentioned here. The following finite example makes that distinction precise.
\end{dcomment}


\dgrblock{dgrb1189}{A finite complete lattice whose localization loses binary limits}
\begin{dcomment}{development}{p48:Le résultat qu'il annonce serait-il vrai?}{A finite complete lattice whose localization loses binary limits}
\textbf{An eight-object counterexample to the broader request in 3.14.}\ours
There is a finite complete and cocomplete category $P$, with a strongly saturated localizer $W$ stable under filtered colimits and accessible as a full subcategory of the arrow category, for which neither homotopy Kan extension along the discrete two-object category $D\to e$ exists. This elementary construction was prepared for this commentary. A targeted search for localizations of finite and Boolean lattices and failures of products or coproducts did not establish an attribution for this exact presentation. No priority claim is made. This is not a counterexample to the cofibration-category theorem just cited.

Let $P$ be the Boolean lattice of subsets of $\{1,2,3\}$, ordered by inclusion, and abbreviate subsets by their strings of elements. Put $a=1$, $b=12$, and let $W$ consist of the identities and $s:a\to b$. The two-out-of-three property holds: $a<b$ is a cover, so a factorization of its only nonidentity weak equivalence has an identity factor. The underlying category has all small limits and colimits, computed by intersections and unions.

We first compute the localization as a category, not just its object classes. Adjoin $s^{-1}:b\to a$. Any composable word containing two such inverses has the intervening forward path $a\to b$, which is $s$; it therefore reduces to a word with at most one inverse. Thus a morphism $x\to y$ is represented either by an inclusion or by
\[
 x\longrightarrow b\xrightarrow{s^{-1}}a\longrightarrow y.
\]
If both forms occur, then $x\subseteq b$, $a\subseteq y$, and $x\subseteq y$. Either $x\subseteq a$, or $x$ contains $2$ and therefore $b\subseteq y$. In the first case the inverse cancels on the source side; in the second it cancels on the target side. The two forms consequently represent the same arrow. The localization is thin, with relation
\[
 x\preceq y\quad\Longleftrightarrow\quad
 x\subseteq y\ \text{ or }\ (x\subseteq12\text{ and }1\subseteq y).
\]
This also verifies its universal property: for any functor $F$ inverting $s$, the new arrow must be sent to
$F(a\to y)F(s)^{-1}F(x\to b)$; the reductions just proved verify compositions and uniqueness. Only $a$ and $b$ become isomorphic. A skeleton $Q$ therefore has seven objects, and $W$ is exactly the inverse image of its isomorphisms, hence is strongly saturated.

In $Q$ the common upper bounds of $2$ and $3$ are $13$, $23$, and $123$. The first two are incomparable, so no least upper bound exists. Dually the common lower bounds of $13$ and $23$ are $\emptyset$, $2$, and $3$, and there is no greatest one. Thus $Q$ lacks both a binary coproduct and a binary product. Localization of the product relative category $(P\times P,W\times W)$ gives $Q\times Q$: one may first invert $(w,1)$ and then $(1,w)$, and these arrows generate the same inversions as $W\times W$. The restriction along $D\to e$ consequently identifies with the diagonal $Q\to Q\times Q$. Its left and right adjoints would give exactly the missing binary coproducts and products. Both therefore fail to exist.

Finally, all the accessibility qualifications can be checked directly. A filtered diagram in a finite poset has a finite directed image, hence a largest value which occurs at some stage and is its colimit. Every object of $P$ is therefore finitely presentable. The same observation applies to the finite arrow category $P^{[1]}$. A filtered diagram all of whose objects lie in $W$ has its colimit equal to one of those objects, hence in $W$. The full subcategory on $W$ is itself finitely accessible: it has filtered colimits, every object is finitely presentable, and its finite set of objects supplies a generating set. The inclusion preserves filtered colimits. A cardinal filtration is trivial for this finite category. Thus the counterexample persists under the explicit additional hypotheses in the final sentences of 3.14. What fails is the extra homotopical structure supplied by the cofibration axioms, not any smallness or ordinary completeness condition.
\end{dcomment}


\textbf{Question 3.15.} Cela me ramène à la question particulièrement élusive et irritante, de l'existence des petits ensembles Q-générateurs, dans une partie $\Phi$ de $\Fl(\mathcal{M})$, quand on suppose (disons) que $\mathcal{M}$ est stable par petites $\limproj$ , que ses éléments sont accessibles et qu'il admet une filtration cardinale. Ainsi on a la question : une partie L-accessible de $\Fl(\mathcal{M})$, qu'on suppose au besoin Q-saturé à gauche (ou ce qui revient au même, contenant les isomorphismes, et stable par cochangement de base, par composition et par facteurs directs) admet-elle une \textit{petite partie} $\Phi_0$ Q-génératrice à gauche? % [sic] Q-saturé -> Q-saturée, agreement with « une partie »

\dgrblock{dgrb1190}{An accessible class with no small Q-generating subset}
\begin{dcomment}{development}{p49:Q-génératrice à gauche?}{An accessible class with no small Q-generating subset}
\textbf{A negative answer to the Q-saturated form of Question 3.15.}\ours
Let $\mathcal K$ be the category of directed multigraphs, with sets of edges and vertices and two functions $s,t:E\to V$, and let $\Phi$ be its class of pure monomorphisms. Then $\mathcal K$ is bicomplete, every object is accessible, and cardinality supplies a cardinal filtration. The full subcategory of $\mathcal K^{[1]}$ on $\Phi$ is accessible and is closed under all filtered colimits. The class $\Phi$ contains the isomorphisms and is closed under pushouts, transfinite compositions, and retracts, yet it has no small Q-generating subset. The non-generation theorem is due to \citet[Theorem~5.1]{CoxEtAl2026Pure}; the detailed translation to the hypotheses of this manuscript is the development supplied here. We make no priority claim for that translation.

Write $\mathcal K=\mathbf{Set}^{\mathcal E}$, where $\mathcal E$ has objects $e,v$, the distinct arrows $s,t:e\to v$, and identities. Limits and colimits are computed on the two sets, and finite graphs are finitely presentable. More generally a graph of cardinality less than a regular infinite cardinal $\lambda$ is $\lambda$-presentable: maps out of it and equalities between such maps involve fewer than $\lambda$ elements and equations, all of which occur at a common stage of a $\lambda$-filtered diagram. Thus every graph is presentable at some rank, and the subcategories with bounded cardinality are essentially small.

For an inclusion $A\subseteq B$, purity means that every finite system of equations in vertices and edges, with parameters from $A$, which has a solution in $B$, has a solution in $A$. Equivalently, in a commutative square from an arrow $P\to Q$ of finitely presentable graphs to $A\to B$, the top map $P\to A$ extends to $Q\to A$; equality with the given bottom map into $B$ is not required. These are the equivalent descriptions in \citet[Definition~2.7 and Remark~2.8]{CoxEtAl2026Pure}.

We verify accessibility of the \emph{full arrow subcategory}, whose maps are arbitrary commuting squares. This is distinct from the category with graphs as objects and only pure maps as morphisms. Given a pure inclusion $A\subseteq B$ and countable subsets in its domain and codomain, first close them under $s,t$ and ensure that the chosen part of $A$ is included in the chosen part of $B$. Suppose countable subgraphs $A_n\subseteq A$ and $B_n\subseteq B$, with $A_n\subseteq B_n$, have been chosen. There are only countably many finite equation systems with parameters in $A_n$. For each such system that has a solution in $B_n$, choose a solution in $A$, using purity of $A\subseteq B$. Adjoin all these witnesses to $A_n$, close under $s,t$, and adjoin their images to $B_n$. Repeat for $n<\omega$. The unions $A'\subseteq B'$ are countable. Any finite system over $A'$ with a solution in $B'$ has all its parameters and witnesses in a common stage, and was therefore given a solution in $A'$ at the next stage. Hence $A'\to B'$ is pure.

The poset of countable pure subarrows of $A\to B$ is $\aleph_1$-directed: for a countable collection take the unions of their domains and codomains as seeds and apply the preceding closure construction. Its union is the whole arrow, since any chosen vertex or edge can be included in the seeds. There is a set of isomorphism types of countable graph arrows. Each is $\aleph_1$-presentable in $\mathcal K^{[1]}$, by the same element-and-equation argument already used for graphs, with the two components of the arrow included as additional functions. In particular every pure arrow is a filtered colimit of a fixed small family of bounded-size pure arrows, exactly the filtered-colimit-envelope condition used by the manuscript's earlier accessibility language.

We must also check preservation of the relevant colimits. Let $f_i:A_i\to B_i$ be a filtered diagram of pure arrows with arbitrary commuting-square transition maps. A square from finite $P\to Q$ into its colimit factors through a common stage, and its required equality holds at some later stage, because $P,Q$ are finitely presentable. Purity at that stage gives $Q\to A_i$ extending the top map; composing into the colimit proves purity of the colimit arrow. Filtered colimits of sets preserve injections, so it is again a monomorphism. Thus the full subcategory on $\Phi$ has filtered colimits created in $\mathcal K^{[1]}$. Its countable arrows remain $\aleph_1$-presentable in that full subcategory, and the preceding directed presentation proves that it is $\aleph_1$-accessible. This verifies both parts of the manuscript's L-accessibility requirement, rather than inferring them from accessibility of a different category of pure maps.

For completeness, the closure conditions on $\Phi$ are also those required in the question. Composition follows by solving the same finite system first in the intermediate graph and then in the source. A retract of a pure arrow is pure: carry a parameter system and its solution along the section of the retract diagram, solve it in the larger source, and carry the solution back along the retraction. Arbitrary pushouts of pure monomorphisms are pure, as recorded in \citet[Fact~2.9(iii)]{CoxEtAl2026Pure}. Finally transfinite closure follows from composition at successor stages and the filtered-colimit argument at limit stages, applied to the arrows from the fixed starting object. Thus $\Phi$ is Q-saturated in the sense of the question, and in fact has the additional transfinite stability used elsewhere in the chapter.

The cited classification says that pure monomorphisms in $\mathbf{Set}^{\mathcal E}$ are cofibrantly generated precisely when, for every span $y\leftarrow x\to z$ in $\mathcal E$, one leg factors through the other. The span $v\xleftarrow{s}e\xrightarrow{t}v$ fails this test: the only endomorphism of $v$ is its identity, and $s\ne t$. Hence $\Phi$ is not cofibrantly generated.

To connect that conclusion exactly to Grothendieck's terminology, suppose a small subset $S\subseteq\Phi$ were Q-generating. Its defining property is
\[
 \Phi\subseteq{}^\perp(S^\perp).
\]
The small object argument applies in the locally presentable category $\mathcal K$. It identifies ${}^\perp(S^\perp)$ with $\operatorname{cof}(S)$, the class of retracts of transfinite composites of pushouts of maps in $S$. The closure properties just proved give $\operatorname{cof}(S)\subseteq\Phi$. Both inclusions force equality, contradicting the classification theorem. This proves the asserted negative answer with no additional set-theoretic principle. The example concerns Q-\emph{saturation}; it does not silently strengthen the question to the different hypothesis that $\Phi$ is already Q-\emph{closed}. Nor does the argument identify accessibility of an arrow class with combinatoriality of a model structure, since no model structure has been supplied.
\end{dcomment}


\marginpar{! \textbar}D'autre part, je n'ai pas à présent construit un seul exemple d'un couple $\Phi \longleftrightarrow \Psi$, dans une catégorie $\mathcal{M}$, qui ne soit factorisant! Le cas où $\mathcal{M}$ est spéciale, était bien sûr le plus intéressant. Visiblement, dans tout
\msp{page 50}
ça ce sont les exemples qui manquent. Le clef d'une compréhension semble bien être dans les questions de factorisation entre ensembles $\Phi \longleftrightarrow \Psi$. Il faudrait en faire une étude systématique, mais je recule devant la perspective d'une digression qui me prendrait peut-être encore des semaines! % [sic] Le clef -> La clef, « clef » is feminine


\dgrblock{dgrb1191}{Mutual lifting classes need not factor every arrow}
\begin{dcomment}{examples}{p50:Il faudrait en faire une étude systématique}{Mutual lifting classes need not factor every arrow}
The failure of a mutually orthogonal pair to admit factorizations is already illustrated by \citet[Example~5.7.10]{Taylor1999PracticalFoundations}. The following rational-cut presentation gives a small elementary instance, with no priority claim; it does not meet the stronger condition that the ambient category be special. Regard $\mathbf Q$ as a category under its usual order, and fix an irrational real number $r$, for example $\sqrt2$. Define $\mathcal L$ to contain the identities and the nonidentity arrows $a\to b$ with $b<r$. Define $\mathcal R$ to contain the identities and the nonidentity arrows $x\to y$ with $r<x$.

In a poset a lifting problem from $a\to b$ to $x\to y$ has a solution exactly when $b\leq x$, provided the square exists. Nonidentity arrows in the displayed classes have $b<r<x$, so $\mathcal L\perp\mathcal R$. Conversely, if a nonidentity $a\to b$ is not in $\mathcal L$, then $b>r$. Choose a rational $x$ with $\max(a,r)<x<b$. The square from $a\to b$ to $x\to b\in\mathcal R$, with bottom identity, has no lift. Hence ${}^\perp\mathcal R=\mathcal L$. If a nonidentity $x\to y$ is not in $\mathcal R$, then $x<r$. Choose rational $b$ with $x<b<\min(y,r)$. The square from $x\to b\in\mathcal L$ to $x\to y$, with top identity, has no lift. Thus $\mathcal L^\perp=\mathcal R$.

Nevertheless, an arrow $a\to b$ with $a<r<b$ has no factorization of the indicated type. An intermediate object $z$ with both factors nonidentity would have to satisfy simultaneously $z<r$ and $r<z$; choosing either factor to be an identity also fails. This is a pair of mutually closed lifting classes, but not a weak factorization system. The rational poset is not cocomplete, so the example answers the unrestricted request only. It leaves the final, more demanding special-ambient question logically separate from the counterexample.
\end{dcomment}
\begin{otherlanguage}{english}
\dgroverview{XIX}{Chapter XIX \textperiodcentered\ \textup{XIX} \textperiodcentered\ Overview}
\begingroup
\selectlanguage{english}
\noindent\textbf{Chapter XIX: factorization, diagram categories, and the existence of derived Kan extensions.}



\medskip
Grothendieck has reached the point at which his formalism must acquire its
examples. Given a category $\mathcal M$ and a class $W$ of weak equivalences,
the intended coefficient system is
\[
 \mathbb D_{\mathcal M,W}(I)
   =\operatorname{Fun}(I^{\mathrm{op}},\mathcal M)[W(I)^{-1}],
 \qquad u^*:\mathbb D(J)\longrightarrow\mathbb D(I).
\]
The problem is to construct both adjoints to restriction, prove their pointwise
formulas, and explain which hypotheses on the models are doing the work. The
chapter's fifty manuscript pages therefore move from lifting and factorization
to model structures on diagrams, and then return to the places where that
method appears to stop. In the edition the chapter occupies printed pages 1--26,
after two preliminary leaves.

The historical motivation precedes the terminology of this chapter. In
\emph{Pursuing Stacks}, \S69, Grothendieck relates the passage from a single
derived category to a system of derived diagram categories to the inadequacy
of the former for the variance and trace constructions arising in SGA~5.
His letter to Thomason of 2 April 1991 connects the long investigation of
accessible categories in Chapter XVIII with the foundations needed for this
coefficient formalism \citep{Grothendieck1983,Grothendieck1991Letter}. Thus
accessibility enters here for a concrete reason: the homotopical operations
must be constructed using sets of generators and controlled colimits before
they can be recognized intrinsically after localization.

\noindent\textbf{The two levels of closure.}
The opening pages isolate the relation $i\mathbin{\perp}p$ of admitting a
diagonal in every lifting square. It gives an antitone correspondence between
classes of arrows, hence the closures
${}^{\perp}(S^{\perp})$ and $({}^{\perp}T)^{\perp}$. Pushouts, transfinite
compositions, and retracts describe a second closure process. The small
object argument identifies the two when a set of arrows generates the
lifting condition and their domains have the required smallness. The
retract argument is the decisive step: factor an arrow already having the
left lifting property, then use its lift against the right factor to express
it as a retract of the constructed cellular arrow. This is why the distinction
between an orthogonal pair and a factorizing pair must remain visible.
The source's initial identification of separate closedness of two orthogonal
classes with mutual determination of the classes is too strong; the
factorization hypothesis supplies the missing implication later.

Grothendieck's word ``special'' must also retain its stated meaning. His
ambient category has small colimits and each object is presentable at some
cardinal; his special left classes have a small set of lifting generators.
This does not itself supply the small generating subcategory required for
local presentability. Modern statements about combinatorial model categories
may consequently be used only after the additional ambient hypothesis has
been made explicit.

\noindent\textbf{The asymmetry of projective and injective diagrams.}
For contravariant diagrams, evaluation at $i$ has a left adjoint
$i_!A(j)=\coprod_{\operatorname{Hom}_I(j,i)}A$. Applying these adjoints to
generating cofibrations and generating trivial cofibrations constructs the
projective structure. Its fibrations and weak equivalences are pointwise;
its cofibrations encode compatibility with the arrows of the indexing
category. Theorem 2.3 then derives $u_!$, and Proposition 1.12 supplies the
restriction property needed to compute its values on comma categories.
In the variance used here the formulas are
\[
 (\mathbf L u_!X)_j\simeq
   \operatorname*{hocolim}_{(j\downarrow u)^{\mathrm{op}}}X,
 \qquad
 (\mathbf R u_*X)_j\simeq
   \operatorname*{holim}_{(u\downarrow j)^{\mathrm{op}}}X.
\]
Writing the opposite categories is necessary: the manuscript's diagrams
are presheaves, whereas much modern derivator notation uses covariant
diagrams.

For the injective structure the cofibrations are pointwise, and it is the
right class whose compatibility is no longer pointwise. One cannot obtain
small generators merely by placing the old generators at every component.
Question 1.11 has a positive answer for locally presentable ambient
categories, by the generation theorem of \citet[Lemma A.2.8.3]{Lurie2009HTT},
but its final suggested choice of generators already fails for sets with a
nontrivial group action. The commentary gives the sharper elementary
criterion for when the sole generator $\varnothing\to 1$ can be used unchanged
on a diagram category, with a complete proof and an explicit statement of
the literature checked.

Theorems 2.6 and 2.9 secure the injective construction in Grothendieck's
``very special'' situation, in which the cofibrations are all monomorphisms.
The geometric conditions on unions and pushouts, together with accessibility
of $W$, are the mechanism behind the recognition theorem. Here two-out-of-three
does essential work: after building a bounded trivial subextension, it keeps
the remaining extension in $W$. This is a substantially stronger situation
than an arbitrary accessible class of monomorphisms. The proof recalled
from Chapter XV must be read through its corrected argument in \S7, not
through the explicitly abandoned general claim in \S6. The cardinal-stage
condition printed with a projective-limit arrow in XIX is identified in the
apparatus as a source slip: the cited formulations require closure under
the corresponding small inductive limits.

The latching interpretation in section 3 returns to the same care about
testing classes. Proposition 3.5 needs the canonical right class, which the
source explicitly chooses in (3.5.2); an arbitrary presentation by right
tests can give a false equality. The commentary supplies the counterexample
and the valid comparison. It also proves the classwise cellular description
conjectured in (3.5.3) for noetherian posets, with the required colimits stated,
and supplies the short lifting lemma implicit in Proposition 3.7.

\noindent\textbf{What later existence theorems change.}
The existence of a convenient diagram model structure and the existence of
derived Kan extensions are different assertions. Cisinski proves the latter
for an arbitrary Quillen model category with small limits; duality supplies
the colimit statement when small colimits exist
\citep[Theorem 6.11]{Cisinski2003}. Consequently the bicomplete Quillen
case of Question 3.9 admits no counterexample of the requested kind.
For combinatorial models, both diagram model structures give a direct and
symmetric proof; \citet[Theorem 1.28 and Proposition 1.30]{Groth2013} presents
this route. The later existence theorem does not retroactively prove every
factorization assertion contemplated in the manuscript, and it does not
make the adjective ``very special'' necessary.

The final section shows why the size of the permitted diagrams matters.
Its bounded-complex example is effective precisely because bounded derived
categories need not contain the result of a homotopy limit indexed by a
finite category with arbitrarily long composable strings. For
$k=\mathbf F_2$ and $G=C_2$, the resolution of the trivial $k[G]$-module is
periodic and gives nonzero group homology and cohomology in every
nonnegative degree. Testing a hypothetical adjoint to the trivial-action
functor against shifts proves its nonexistence within the bounded category.
This supplies the required obstruction without identifying a bounded
category with a bicomplete modern model category. The distinction between
cofibration-category excision on cofibrant objects and the manuscript's
absolute $W$-cofibrations is equally necessary in reading the Brown lemma
invoked on manuscript page 35.

Some of the closing searches have elementary answers. A fully faithful
coreflective image in a bicomplete category is itself bicomplete: the
coreflector constructs limits, and the image is closed under colimits.
Thus the proposed adjoint-localization construction in Example 3.12 cannot
produce its hoped-for failure. The category of ordinals small relative to
the fixed universe has all small suprema and no terminal object, answering
Question 3.13 with the universe convention made explicit. Its lack of a
small conservative generating family is consistent with the convention of
SGA~4. A rationally ordered category cut at an irrational number supplies
a mutually lifting-closed pair which does not factor arrows across the cut;
it does not have the cocompleteness required of a special ambient category.

The broader localizer question in 3.14 admits a small counterexample.
Localize the Boolean lattice of subsets of $\{1,2,3\}$ at the single cover
$\{1\}\subset\{1,2\}$. The resulting seven-object skeleton has neither
a least upper bound for $\{2\},\{3\}$ nor a greatest lower bound for
$\{1,3\},\{2,3\}$. The localizer is strongly saturated, accessible, and
stable under filtered colimits, and the original lattice is complete and
cocomplete. The commentary computes the localization and proves each
assertion. This answers the broad question even with the additional
smallness hypotheses. The cofibration-category question in the same
paragraph has the opposite answer: the transfinite continuity axiom
supplies the coproduct and sequential-colimit hypotheses of modern
cofibration-category theory, which yields all homotopy left Kan
extensions \citep[Theorem 10.3.2]{RadulescuBanu2009}.

Question 3.15 also benefits from separating the closures introduced at the
start. The class of pure monomorphisms in directed multigraphs is accessible,
closed under filtered colimits, pushouts, transfinite composition, and
retracts. It nevertheless has no small set of Quillen generators. The
non-generation theorem is due to Cox, Feigert, Kamsma, Mazari-Armida, and
Rosick\'y, in their 2026 classification for presheaf categories
\citep[Theorem 5.1]{CoxEtAl2026Pure}; the
commentary verifies the translation into the hypotheses stated here and
supplies the small-subdiagram argument for accessibility. This answers the
question for a Q-saturated class. It does not identify that class with its
double lifting orthogonal, and the stronger question obtained by imposing
that additional condition is not silently substituted for the question
Grothendieck actually wrote.

\medskip
The commentary retains the edition's slips as source evidence, distinguishes
them from errors of transcription, and states repairs at the
point where they are needed. The claim register records the source's proof
status separately from later theorems or counterexamples. Mathematical
dependencies on Chapters XIII, XV, XVII, and XVIII are checked against
the relevant original passages or explicitly identified as recalled results.

\endgroup
\end{otherlanguage}
\dgrchunkend
